\begin{filecontents*}[overwrite]{check_profile_universe_v5.py}
"""Exact finite checks for convolution/gcd catalogue; not a proof checker."""
from fractions import Fraction as F
from math import comb, gcd
from itertools import product

def divisors(n):
    return [d for d in range(1,n+1) if n%d==0]
def phi(n):
    return sum(gcd(a,n)==1 for a in range(1,n+1))
def jordan(n,k):
    ans=n**k
    p=2; m=n
    while p*p<=m:
        if m%p==0:
            ans=ans//p**k*(p**k-1)
            while m%p==0:m//=p
        p+=1
    if m>1:ans=ans//m**k*(m**k-1)
    return ans

def gcd_profile(p,a,r,k):
    z,w=F(1,p),F(1,p**a)
    h=max(r,k); d=abs(r-k)
    return 1+a*(1-z**r) if d==0 else 1+(1-z**min(r,k))*z**d*(1-w**d)/(1-z**d)
def unit_profile(p,a,r,k):
    z,w=F(1,p),F(1,p**a)
    if k>r:return (1-z)**r*w**(k-r)+(1-z**k)*(1-w**(k-r))/(1-z**(k-r))
    if k==r:return (1-z)**r+a*(1-z**r)
    return (1-z)**r+(1-z**k)*z**(r-k)*(1-w**(r-k))/(1-z**(r-k))

def main():
    counts={'factorization_profile':0,'gcd_profile':0,'unit_gcd_profile':0,'gcd_direct':0,'unit_gcd_direct':0,'all_subgroups':0,'cyclic_subgroups':0}
    for p,a,h in product([2,3,5,7],range(1,6),range(1,5)):
        z,w=F(1,p),F(1,p**a)
        value=sum(p**(h*j)*(a-j+1) for j in range(a+1))
        profile=(1-z**h*w**h*((a+2)-(a+1)*z**h))/(1-z**h)**2
        assert F(value,p**(h*a))==profile
        counts['factorization_profile']+=1
    for p,a,r,k in product([2,3,5],range(1,5),range(1,5),range(1,5)):
        n=p**a; H=max(r,k)
        gval=sum(p**(k*j)*jordan(p**(a-j),r) for j in range(a+1))
        assert F(gval,p**(H*a))==gcd_profile(p,a,r,k)
        assert gval==sum(p**(r*j)*jordan(p**(a-j),k) for j in range(a+1))
        counts['gcd_profile']+=1
        mval=phi(n)**r*sum((F(jordan(p**j,k),phi(p**j)**r) for j in range(a+1)),F(0))
        assert mval.denominator==1
        assert mval/F(p**(H*a))==unit_profile(p,a,r,k)
        counts['unit_gcd_profile']+=1
    for n,r,k in product(range(1,13),range(1,4),range(1,5)):
        direct=sum(gcd(n,*aa)**k for aa in product(range(1,n+1),repeat=r))
        conv=sum(d**k*jordan(n//d,r) for d in divisors(n))
        assert direct==conv
        counts['gcd_direct']+=1
        units=[a for a in range(1,n+1) if gcd(a,n)==1]
        direct=sum(gcd(n,*(a-1 for a in aa))**k for aa in product(units,repeat=r))
        formula=phi(n)**r*sum((F(jordan(d,k),phi(d)**r) for d in divisors(n)),F(0))
        assert direct==formula
        counts['unit_gcd_direct']+=1
    for p,a,r,s in product([2,3,5],range(1,5),range(1,5),range(1,6)):
        if r>s:continue
        z,w=F(1,p),F(1,p**a); b=s-r
        direct=sum(p**min(i,j) for i in range(r*a+1) for j in range(s*a+1))
        profile=(b*a+1)*(1-z*w**r)/(1-z)+2*z*(1-(r*a+1)*w**r+r*a*z*w**r)/(1-z)**2
        assert F(direct,p**(r*a))==profile
        counts['all_subgroups']+=1
        direct=1+sum(F(p**(min(j,r*a)+min(j,s*a))-p**(min(j-1,r*a)+min(j-1,s*a)),p**j-p**(j-1)) for j in range(1,s*a+1))
        profile=b*a+1+2*z*(1-w**r)/(1-z)
        assert direct/F(p**(r*a))==profile
        counts['cyclic_subgroups']+=1
    print(counts)
if __name__=='__main__':main()

\end{filecontents*}
\begin{filecontents*}[overwrite]{check_projective_functionals.py}
"""Exact finite checks of the polynomial-kernel examples (not a proof checker)."""
from fractions import Fraction as F
from math import comb, factorial

def add(*polys):
    result = {}
    for p in polys:
        for k,v in p.items(): result[k] = result.get(k,F(0))+v
    return {k:v for k,v in result.items() if v}

def mul(p,q):
    result = {}
    for i,a in p.items():
        for j,b in q.items(): result[i+j] = result.get(i+j,F(0))+a*b
    return {k:v for k,v in result.items() if v}

def d(p): return {k-1:k*v for k,v in p.items() if k}
def J(p): return {k+1:v/(k+1) for k,v in p.items()}

for n in range(1,21):
    A = {n:F(1,factorial(n))}
    lhs = J(mul({2:F(1)},J(mul({0:F(1),1:F(1)},d(d(d(A)))))))
    rhs = add(mul({3:F(1),2:F(1)},d(A)),
              mul({2:F(-4),1:F(-2)},A),
              mul({1:F(8),0:F(2)},J(A)),
              mul({0:F(-8)},J(J(A))),
              {3:-F(int(n==2)-int(n==1),3)})
    assert lhs == rhs, n

for m in range(12):
    coeff = [(-1)**(m-j)*comb(m,j)*comb(m+j,m) for j in range(m+1)]
    for n in range(1,21):
        lhs = sum((F(c,n+j+1) for j,c in enumerate(coeff)),F(0))
        rhs = F(factorial(n)**2,factorial(n-m)*factorial(n+m+1)) if n>=m else F(0)
        assert lhs == rhs, (m,n)

print('20 nested-kernel coefficient identities and 240 Legendre moments passed.')

\end{filecontents*}
\begin{filecontents*}[overwrite]{check_rational_kernels_v5.py}
"""Finite exact checks for the integral identities; not a formal proof."""
from fractions import Fraction as F
from functools import reduce
from itertools import permutations
from math import factorial


def perm_sign(p):
    return (-1) ** sum(p[i] > p[j] for i in range(len(p))
                       for j in range(i + 1, len(p)))


def multiply_linear(poly, constant):
    result = [0] * (len(poly) + 1)
    for i, c in enumerate(poly):
        result[i] += constant * c
        result[i + 1] += c
    return result


def falling_coefficients(a):
    poly = [1]
    for j in range(a):
        poly = multiply_linear(poly, -j)
    return poly


def power(m, exponent):
    return F(m ** exponent) if exponent >= 0 else F(1, m ** -exponent)


def run():
    cube_checks = 0
    for m in range(1, 5):
        ps = list(permutations(range(m)))
        for a in range(1, 4):
            for n in range(7):
                actual = sum((F(perm_sign(p) * perm_sign(q)) *
                              reduce(lambda x, j: x * F(1, n + a + p[j] + m * q[j]),
                                     range(m), F(1))
                              for p in ps for q in ps), F(0))
                constant = factorial(m) * m ** (m * (m - 1) // 2)
                for j in range(1, m):
                    constant *= factorial(j) ** 2
                expected = F(constant)
                for j in range(m * m):
                    expected /= n + a + j
                assert actual == expected
                cube_checks += 1

    residue_checks = 0
    for shifts in [(1,), (1, 3), (2, 7, 11), (8, 1, 4, 17)]:
        for n in range(1, 15):
            lhs = reduce(lambda x, a: x / (n + a), shifts, F(1))
            rhs = F(0)
            for a in shifts:
                den = n + a
                for b in shifts:
                    if b != a:
                        den *= b - a
                rhs += F(1, den)
            assert lhs == rhs
            residue_checks += 1

    shift_checks = 0
    for a in range(1, 8):
        coefficients = falling_coefficients(a)
        assert coefficients[1] == (-1) ** (a - 1) * factorial(a - 1)
        for r in range(1, 7):
            for m in range(1, 30):
                actual = sum((F(c, factorial(m)) * power(m, j - r)
                              for j, c in enumerate(coefficients)), F(0))
                if m < a:
                    expected = F(0)
                elif m == a:
                    expected = F(1, a ** r)
                else:
                    expected = F(1, factorial(m - a) * m ** r)
                assert actual == expected
                shift_checks += 1
    print({"bi_vandermonde": cube_checks,
           "product_partial_fractions": residue_checks,
           "inverse_moment_shift_coefficients": shift_checks})


if __name__ == "__main__":
    run()

\end{filecontents*}
\begin{filecontents*}[overwrite]{check_scope_v5.py}
"""Exact finite checks for angular transforms, counting, and signed blocks.
These checks supplement, and do not replace, the general proofs.
Python standard library only.
"""
from fractions import Fraction as F
from itertools import product
from math import comb, gcd

def trim(p):
    while len(p)>1 and p[-1]==0: p.pop()
    return p

def divmod_poly(p,q):
    p=list(map(F,p)); q=list(map(F,q)); a=[F(0)]*max(1,len(p)-len(q)+1)
    while len(p)>=len(q) and any(p):
        k=len(p)-len(q); c=p[-1]/q[-1]; a[k]=c
        for j,t in enumerate(q): p[k+j]-=c*t
        trim(p)
    return trim(a),trim(p)

def rank(rows):
    a=[list(map(F,row)) for row in rows]; r=0
    for c in range(len(a[0])):
        p=next((i for i in range(r,len(a)) if a[i][c]),None)
        if p is None: continue
        a[r],a[p]=a[p],a[r]; z=a[r][c]; a[r]=[v/z for v in a[r]]
        for i in range(r+1,len(a)):
            if a[i][c]:
                z=a[i][c]; a[i]=[u-z*v for u,v in zip(a[i],a[r])]
        r+=1
        if r==len(a): break
    return r

def angular():
    phi={}
    for n in range(1,31):
        p=[-1]+[0]*(n-1)+[1]
        for d in range(1,n):
            if n%d==0:
                p,remainder=divmod_poly(p,phi[d]); assert not any(remainder)
        phi[n]=p
        rows=[]
        for d in range(1,n+1):
            if n%d: continue
            deg=len(phi[d])-1; block=[[F(0)]*n for _ in range(deg)]
            for k in range(n):
                _,rem=divmod_poly([0]*k+[1],phi[d])
                for j,v in enumerate(rem): block[j][k]=v
            rows+=block
        assert len(rows)==n and rank(rows)==n
    return 30

def factors(n):
    ans={}; p=2
    while p*p<=n:
        while n%p==0: ans[p]=ans.get(p,0)+1; n//=p
        p+=1
    if n>1: ans[n]=ans.get(n,0)+1
    return ans

def surjection_formula(n,m,r):
    v=F(n**(m*r))
    for p in factors(n):
        for j in range(m-r+1,m+1): v*=1-F(1,p**j)
    assert v.denominator==1
    return int(v)

def finite_geometry():
    checks=0
    for n in [2,3,4,5,6,8]:
        for m in [1,2,3]:
            count=sum(gcd(n,*v)==1 for v in product(range(n),repeat=m))
            assert count==surjection_formula(n,m,1); checks+=1
        count=sum(gcd(n,a*d-b*c)==1 for a,b,c,d in product(range(n),repeat=4))
        assert count==surjection_formula(n,2,2); checks+=1
        lines={frozenset(((t*a)%n,(t*b)%n) for t in range(n))
               for a,b in product(range(n),repeat=2) if gcd(n,a,b)==1}
        expected=F(surjection_formula(n,2,1),surjection_formula(n,1,1))
        assert len(lines)==expected; checks+=1
    return checks

def signed_row(fac):
    omega=len(fac); big=sum(fac.values()); squarefree=all(a==1 for a in fac.values())
    mu=(-1)**omega if squarefree else 0; lam=(-1)**big
    row=[c**omega for c in range(1,4)]
    row += [lam*c**omega for c in range(1,4)]
    row += [mu*mu*c**omega for c in range(1,4)]
    row += [mu*c**omega for c in range(1,4)]
    for r in range(2,5):
        v=1
        for a in fac.values(): v*=comb(a+r-1,r-1)
        row.append(v)
    row += [c**sum(a for p,a in fac.items() if p>c) for c in range(2,5)]
    row += [int(omega==1 and big==1)]
    return row

def signed():
    primes=[11,13,17,19,23,29,31,37,41]
    samples=[{11:a} for a in range(1,16)]
    for k in range(1,10):
        samples += [{p:2 for p in primes[:k]},
                    {p:(3 if i==0 else 2) for i,p in enumerate(primes[:k])},
                    {p:1 for p in primes[:k]}]
    rows=[signed_row(f) for f in samples]
    assert len(rows[0])==19 and rank(rows)==19
    return len(rows),19

if __name__=='__main__':
    print('Cyclotomic CRT matrices:',angular(),'full rank')
    print('Direct finite-ring counts:',finite_geometry(),'passed')
    rows,cols=signed(); print('Signed arithmetic sample matrix:',rows,'x',cols,'rank',cols)

\end{filecontents*}
\begin{filecontents*}[overwrite]{check_extensions_v4.py}
#!/usr/bin/env python3
"""Exact auxiliary checks for the 2026-10-29 extensions.

These finite checks supplement, but do not prove, the infinite theorems.
Only the Python standard library is used.
"""
from fractions import Fraction as F
from itertools import combinations, product
from math import factorial, gcd, lcm
from pathlib import Path
import json


def mul(a, b):
    c = [F(0)] * (len(a) + len(b) - 1)
    for i, x in enumerate(a):
        for j, y in enumerate(b):
            c[i+j] += x*y
    return c


def ev(a, w):
    ans = F(0)
    for x in reversed(a):
        ans = ans*w+x
    return ans


def sig_profile(z, k):
    c = [F(0)] * (k+1)
    c[0], c[k] = 1/(1-z**k), -z**k/(1-z**k)
    return c


def lattice_profile(z, rank):
    c = [F(1)]
    for j in range(1, rank):
        c = mul(c, [1/(1-z**j), -z**j/(1-z**j)])
    return c


def profiles(z):
    return {
        "sigma1": sig_profile(z, 1),
        "sigma2": sig_profile(z, 2),
        "sigma1_squared": mul(sig_profile(z, 1), sig_profile(z, 1)),
        "sigma1_sigma2": mul(sig_profile(z, 1), sig_profile(z, 2)),
        "phi_cubed": [(1-z)**3],
        "psi_squared": [(1+z)**2],
        "radical": [F(0), 1/z],
        "unitary_sigma2": [F(1), F(0), F(1)],
        "unitary_jordan2": [F(1), F(0), F(-1)],
        "lattices_rank3": lattice_profile(z, 3),
        "lattices_rank4": lattice_profile(z, 4),
        "divisor_d_phi": [1/(1+z), F(0), z/(1+z)],
    }


DEGREES = {
    "sigma1": 1, "sigma2": 2, "sigma1_squared": 2,
    "sigma1_sigma2": 3, "phi_cubed": 3, "psi_squared": 2,
    "radical": 1, "unitary_sigma2": 2, "unitary_jordan2": 2,
    "lattices_rank3": 2, "lattices_rank4": 3, "divisor_d_phi": 2,
}


def U(c, z):
    return (1-z)*(1+sum((a*z**(j+1)/(1-z**(j+1))
                        for j, a in enumerate(c)), F(0)))


def delta(c, z, t):
    return (ev(c, z**t)-1)/U(c, z)


def divisors(n):
    return [d for d in range(1, n+1) if n % d == 0]


def factor(n):
    ans = []
    p = 2
    while p*p <= n:
        if n % p == 0:
            a = 0
            while n % p == 0:
                n //= p
                a += 1
            ans.append((p, a))
        p += 1
    if n > 1:
        ans.append((n, 1))
    return ans


def phi(n):
    return sum(gcd(j, n) == 1 for j in range(1, n+1))


def lattice_count(n, rank):
    """Independent Hermite-normal-form recurrence, not prime-power formula."""
    if rank == 1:
        return 1
    return sum(d**(rank-1)*lattice_count(n//d, rank-1) for d in divisors(n))


def direct(n):
    ds = divisors(n)
    ps = factor(n)
    s1, s2 = sum(ds), sum(d*d for d in ds)
    rad = 1
    psi = n
    for p, _ in ps:
        rad *= p
        psi = psi*(p+1)//p
    us2 = sum(d*d for d in ds if gcd(d, n//d) == 1)
    uj2 = 1
    for p, a in ps:
        uj2 *= p**(2*a)-1
    return {
        "sigma1": s1, "sigma2": s2, "sigma1_squared": s1*s1,
        "sigma1_sigma2": s1*s2, "phi_cubed": phi(n)**3,
        "psi_squared": psi*psi, "radical": rad,
        "unitary_sigma2": us2, "unitary_jordan2": uj2,
        "lattices_rank3": lattice_count(n, 3),
        "lattices_rank4": lattice_count(n, 4),
        "divisor_d_phi": sum(d*phi(d) for d in ds),
    }


def determinant(a):
    a = [[F(x) for x in row] for row in a]
    det = F(1)
    for j in range(len(a)):
        pivot = next(i for i in range(j, len(a)) if a[i][j])
        if pivot != j:
            a[pivot], a[j] = a[j], a[pivot]
            det = -det
        x = a[j][j]
        det *= x
        a[j] = [t/x for t in a[j]]
        for i in range(j+1, len(a)):
            x = a[i][j]
            a[i] = [y-x*z for y, z in zip(a[i], a[j])]
    return det


def encode(x):
    return {"numerator": str(x.numerator), "denominator": str(x.denominator)}


def series_product(a, b, length):
    return (mul(a, b)+[F(0)]*length)[:length]


def reciprocal_series(factors, length):
    """Coefficients of product (1+a*t)^-1 to a fixed exact order."""
    result = [F(1)]+[F(0)]*(length-1)
    for a in factors:
        result = series_product(result, [F((-a)**k) for k in range(length)], length)
    return result


def main():
    counts = {}
    # Compare completely independent arithmetic definitions against profiles.
    count = 0
    for n in range(1, 81):
        expected = direct(n)
        for name, h in DEGREES.items():
            value = F(n**h)
            for p, a in factor(n):
                value *= ev(profiles(F(1, p))[name], F(1, p**a))
            assert value.denominator == 1
            assert value == expected[name], (n, name, value, expected[name])
            count += 1
        # Known lcm-sum identity supplies the bridge to a nonmultiplicative sum.
        assert 2*sum(lcm(j, n) for j in range(1, n+1)) == n*(1+expected["divisor_d_phi"])
    counts["arithmetic_profile_comparisons"] = count
    counts["lcm_sum_identities"] = 80

    # Check the infinite local mean via its exact finite geometric remainder.
    # For degree-j monomial, omitted sum_{a>A} z^{a(j+1)} is explicit.
    local_checks = 0
    for p in (2, 3, 5, 11, 101):
        z = F(1, p)
        for c in profiles(z).values():
            assert U(c, z) > 0
            for A in (1, 2, 5, 10):
                partial = (1-z)*(1+sum(ev(c, z**a)*z**a for a in range(1, A+1)))
                rest = (1-z)*sum(a*z**((A+1)*(j+1))/(1-z**(j+1))
                                 for j, a in enumerate(c))
                assert partial+rest == U(c, z)
                for t in (1, 2, 3):
                    on, off = ev(c, z**t)/U(c, z), 1/U(c, z)
                    assert on-off == delta(c, z, t)
                    local_checks += 1
    counts["local_mean_and_response_checks"] = local_checks

    # The radical channel vanishes at depth one, but not at depth two.
    for p in (2, 3, 5, 7, 11):
        z = F(1, p)
        c = profiles(z)["radical"]
        assert delta(c, z, 1) == 0
        assert delta(c, z, 2) != 0
    counts["radical_depth_tests"] = 10

    # Choose one nonzero response for a target and a separate fresh-prime
    # response distinguishing it from each other channel. Verify the exact
    # multiplicative interpolation diagonal, including all off-diagonal zeros.
    primes = [101, 103, 107, 109, 113, 127, 131, 137, 139, 149, 151, 157, 163]
    names = list(DEGREES)
    certificates = {}
    isolation_checks = 0
    for target in names:
        z0 = F(1, primes[0])
        cs0 = profiles(z0)
        t0 = next(t for t in range(1, 8) if delta(cs0[target], z0, t))
        selectors = []
        for p, other in zip(primes[1:], [n for n in names if n != target]):
            z, cs = F(1, p), profiles(F(1, p))
            t = next(t for t in range(1, 8)
                     if delta(cs[target], z, t) != delta(cs[other], z, t))
            selectors.append((p, t, other))
        diag = None
        for name in names:
            value = delta(cs0[name], z0, t0)
            for p, t, other in selectors:
                z, cs = F(1, p), profiles(F(1, p))
                value *= delta(cs[name], z, t)-delta(cs[other], z, t)
            assert bool(value) == (name == target)
            if name == target:
                diag = value
            isolation_checks += 1
        certificates[target] = {"base_prime": primes[0], "base_depth": t0,
                                "selectors": selectors, "diagonal": encode(diag)}
    counts["mixed_prime_profile_isolation_checks"] = isolation_checks

    # Exact mod-four rational transform and explicit inverse equations.
    matrix = [[1, 1, 1, 1], [1, -1, 1, -1], [1, 0, -1, 0], [0, 1, 0, -1]]
    det = determinant(matrix)
    assert det != 0
    residue_checks = 0
    for vec in product(range(-2, 3), repeat=4):
        a, b, c, d = [sum(x*y for x, y in zip(row, vec)) for row in matrix]
        inverse = [(F(a+b, 2)+c)/2, (F(a-b, 2)+d)/2,
                   (F(a+b, 2)-c)/2, (F(a-b, 2)-d)/2]
        assert inverse == list(vec)
        residue_checks += 1
    counts["mod_four_inverse_checks"] = residue_checks

    # Periodic grid: check every fiber's shift, all terminal uniqueness,
    # multiplier coprimalities, and congruence classes at several moduli.
    grid_checks = 0
    for H, L in ((1, 3), (2, 3), (3, 2)):
        radix = H*(L-1)+1
        Fmax = max(radix**H-1, H+1)
        for q in (1, 2, 4, 6):
            D = q*factorial(Fmax)
            B = 1+H*D*Fmax
            rows = []
            for e in product(range(L), repeat=H):
                I = sum(e[a]*radix**a for a in range(H))
                J = sum((a+1)*e[a]*radix**a for a in range(H))
                p, t = B+D*I, D*J
                shifts = [(j+1)*p-t for j in range(H+1)]
                assert p % q == 1 % q and t % q == 0
                assert all(gcd(p, j+1) == 1 for j in range(H+1))
                assert shifts[-1] >= (H+1)*B
                assert max(shifts[:-1]) < (H+1)*B
                rows.append((e, p, t, shifts))
            assert len({row[3][-1] for row in rows}) == len(rows)
            for row, other in combinations(rows, 2):
                assert gcd(row[1], other[1]) == 1
                for j in range(H):
                    if all(row[0][a] == other[0][a] for a in range(H) if a != j):
                        assert row[3][j] == other[3][j]
                grid_checks += 1
    counts["periodic_grid_pair_checks"] = grid_checks

    # Offset grid for antiderivatives: lower shift indices are negative, so
    # the usual CRT n_e=0 mod p_e would fail at l=-1. Test the repaired
    # n_e=R+1 mod p_e condition as well as the signed translation t_e.
    offset_checks = 0
    for R, J, L in ((1, 0, 3), (2, 0, 2), (1, 1, 3), (2, 1, 2), (1, 3, 2)):
        M = R+J
        radix = M*(L-1)+1
        Fmax = max(radix**M-1, M+2)
        C = max(R, J)+1
        for q in (1, 4):
            D = q*factorial(Fmax)
            B = 1+(J+2*C+M+2)*D*Fmax
            rows = []
            for e in product(range(L), repeat=M):
                I = sum(e[a]*radix**a for a in range(M))
                T = sum((a-R+1)*e[a]*radix**a for a in range(M))
                p, t = B+D*I, D*T
                shifts = [(ell+1)*p-t for ell in range(-R, J+1)]
                assert p % q == 1 % q and t % q == 0
                assert shifts[-1] >= (J+1)*B
                assert max(shifts[:-1]) < (J+1)*B
                for ell in range(-R, J+1):
                    assert gcd(p, R+ell+2) == 1
                    for k in (0, 1, 3):
                        n_e = R+1+k*p
                        assert gcd(p, n_e+ell+1) == 1
                rows.append((e, p, t, shifts))
            assert len({row[3][-1] for row in rows}) == len(rows)
            for row, other in combinations(rows, 2):
                assert gcd(row[1], other[1]) == 1
                for a in range(M):
                    if all(row[0][b] == other[0][b] for b in range(M) if b != a):
                        assert row[3][a] == other[3][a]
                offset_checks += 1
    counts["offset_grid_pair_checks"] = offset_checks

    # Verify centered inverse-product coefficients independently by exact
    # formal series multiplication, including negative and zero factors.
    centered_checks = 0
    for r in range(6):
        for j in range(9):
            length = 12
            factors = [r-m for m in range(j+1)]
            inverse = reciprocal_series(factors, length)
            denominator = [F(1)]
            for a in factors:
                denominator = mul(denominator, [F(1), F(a)])
            assert series_product(inverse, denominator, length) == [F(1)]+[F(0)]*(length-1)
            assert inverse[0] == 1
            centered_checks += 1
    counts["shifted_centered_series_checks"] = centered_checks

    # Test the periodic Touchard/falling-factorial coefficient rearrangement
    # term by term. Its n=0 correction must be precisely -R_0(0).
    periodic_poly_checks = 0
    for q in (1, 2, 3, 4, 7):
        c = [[F((3*r+2*j)%7-3, j+1) for j in range(6)] for r in range(q)]
        for n in range(31):
            lhs = sum(c[n % q][j]*F(factorial(n), factorial(n-j))
                      for j in range(min(5, n)+1))
            rhs = sum(c[(s+j) % q][j]*F(factorial(n), factorial(n-j))
                      for s in range(q) for j in range(min(5, n)+1)
                      if (n-j) % q == s)
            assert lhs == rhs
            periodic_poly_checks += 1
    counts["periodic_polynomial_rearrangements"] = periodic_poly_checks

    output = {"status": "passed", "scope": "finite exact auxiliary checks only",
              "counts": counts, "mod_four_determinant": encode(det),
              "profile_isolation_certificates": certificates}
    target = Path(__file__).with_name("extensions_v4_checks.json")
    target.write_text(json.dumps(output, indent=2)+"\n")
    print(json.dumps({"status": "passed", "counts": counts,
                      "mod_four_determinant": str(det), "output": str(target)}, indent=2))


if __name__ == "__main__":
    main()

\end{filecontents*}
\begin{filecontents*}[overwrite]{check_dyadic_height.py}
#!/usr/bin/env python3
"""Finite exact checks for the dyadic-height proof, not its infinite proof."""
from fractions import Fraction as F
from math import factorial
from pathlib import Path
import json


def v2(n):
    assert n > 0
    return (n & -n).bit_length() - 1


def sigma5(n):
    return sum(d**5 for d in range(1, n + 1) if n % d == 0)


def encode(q):
    return {"numerator": str(q.numerator), "denominator": str(q.denominator)}


def main():
    fac = [factorial(n) for n in range(129)]
    differences = 0
    for n in range(129):
        assert v2(fac[n]) == n - n.bit_count()
        for m in range(n + 1):
            assert v2(fac[n] // fac[m]) - (n - m) == m.bit_count() - n.bit_count()
            differences += 1
    for n in range(100001):
        assert 2**n.bit_count() <= n + 1
    for j in range(1, 61):
        n = 2**j - 1
        assert 2**n.bit_count() == n + 1
    lattice = 0
    total_terms = 0
    for b in (1, 2, 3, 5):
        for sign in (-1, 1):
            x = F(sign * 2, b)
            for n in range(41):
                scale = 2**n.bit_count()
                for m in range(n + 1):
                    term = scale * F(fac[n], fac[m]) * x**(m - n)
                    assert term.denominator == 1
                    lattice += 1
                for q in (1, 3, 7, 11):
                    for u in (-3, 2):
                        term = q * scale * fac[n] * x**(-n) * F(u, q)
                        assert term.denominator == 1
                        total_terms += 1
    cutoff = 80
    lower = sum((F(2**n * sigma5(n), factorial(n))
                 for n in range(1, cutoff + 1)), F(0))
    first = cutoff + 1
    ratio = F(2 * (first + 1)**4, first**5)
    assert ratio < F(1, 2)
    # For all n>0, the logarithmic derivative of (n+1)^4/n^5
    # is 4/(n+1)-5/n<0; thus subsequent majorant ratios decrease.
    error = F(4 * 2**first * first**5, factorial(first))
    upper = lower + error
    digits = 50
    scale = 10**digits
    lo_digits = lower.numerator * scale // lower.denominator
    hi_digits = upper.numerator * scale // upper.denominator
    assert lo_digits == hi_digits
    whole, fraction = divmod(lo_digits, scale)
    prefix = f"{whole}.{fraction:0{digits}d}"
    result = {
        "scope": "Finite exact arithmetic checks; not a proof of irrationality",
        "factorial_valuation_checks": 129,
        "factorial_quotient_valuation_checks": differences,
        "binary_height_checks": 100001,
        "height_equality_cases": 60,
        "partial_sum_lattice_checks": lattice,
        "rational_total_lattice_checks": total_terms,
        "constant": "sum_{n>=1} 2^n sigma_5(n)/n!",
        "cutoff": cutoff,
        "lower": encode(lower),
        "upper": encode(upper),
        "tail_bound": encode(error),
        "first_majorant_ratio": encode(ratio),
        "certified_decimal_prefix": prefix,
        "all_checks": "passed",
    }
    target = Path(__file__).with_name("dyadic_height_checks.json")
    target.write_text(json.dumps(result, indent=2) + "\n")
    print(json.dumps({k: v for k, v in result.items()
                      if k not in ("lower", "upper", "tail_bound", "first_majorant_ratio")}, indent=2))
    print("certificate:", target)


if __name__ == "__main__":
    main()

\end{filecontents*}
\begin{filecontents*}[overwrite]{check_prime_responses.py}
#!/usr/bin/env python3
"""Exact finite audits for factorial moments and mixed-prime separation.

Python standard library only.  Every matrix calculation uses Fraction.
These finite checks do NOT prove the all-degree independence theorem;
the manuscript's general mixed-prime argument supplies that proof.
"""
from fractions import Fraction as F
from math import prod
from pathlib import Path
import json


def factor(n):
    assert n >= 1
    result = {}
    p = 2
    while p*p <= n:
        while n % p == 0:
            result[p] = result.get(p, 0) + 1
            n //= p
        p = 3 if p == 2 else p+2
    if n > 1:
        result[n] = 1
    return result


def is_prime(n):
    return n >= 2 and factor(n) == {n: 1}


def primes_after(bound, count):
    answer = []
    p = bound+1
    while len(answer) < count:
        if is_prime(p):
            answer.append(p)
        p += 1
    return answer


def rank(rows):
    a = [[F(x) for x in row] for row in rows]
    if not a:
        return 0
    height, width = len(a), len(a[0])
    r = 0
    for c in range(width):
        pivot = next((i for i in range(r, height) if a[i][c]), None)
        if pivot is None:
            continue
        a[r], a[pivot] = a[pivot], a[r]
        q = a[r][c]
        a[r] = [v/q for v in a[r]]
        for i in range(r+1, height):
            q = a[i][c]
            if q:
                a[i] = [x-q*y for x, y in zip(a[i], a[r])]
        r += 1
        if r == height:
            break
    return r


def determinant(rows):
    a = [[F(x) for x in row] for row in rows]
    n = len(a)
    assert all(len(row) == n for row in a)
    result = F(1)
    for c in range(n):
        pivot = next((i for i in range(c, n) if a[i][c]), None)
        if pivot is None:
            return F(0)
        if pivot != c:
            a[c], a[pivot] = a[pivot], a[c]
            result = -result
        q = a[c][c]
        result *= q
        for i in range(c+1, n):
            ratio = a[i][c]/q
            for j in range(c+1, n):
                a[i][j] -= ratio*a[c][j]
            a[i][c] = F(0)
    return result


def encode(q):
    q = F(q)
    return {"numerator": str(q.numerator), "denominator": str(q.denominator)}


def response_row(H, prime):
    z = F(1, prime)
    # At total degree H, d=H-k for k=1,...,H.
    # The first H columns are sigma responses, the last H Jordan responses.
    return [z**k for k in range(1, H+1)] + [
        -z**k/(1-z**(k+1)) for k in range(1, H+1)
    ]


def check_responses():
    records = []
    for H in range(1, 7):
        primes = primes_after(H+1, 2*H)
        rows = [response_row(H, p) for p in primes]
        initial_rank = rank(rows)
        # A fallback is allowed, but the actual run records whether it is used.
        while rank(rows) < 2*H:
            assert len(primes) < 8*H, "Increase the search bound, not the theorem."
            p = primes_after(primes[-1], 1)[0]
            primes.append(p)
            rows.append(response_row(H, p))
        # Select an independent square submatrix, in the order sampled.
        selected_rows, selected_primes = [], []
        for p, row in zip(primes, rows):
            if rank(selected_rows+[row]) > len(selected_rows):
                selected_rows.append(row)
                selected_primes.append(p)
            if len(selected_rows) == 2*H:
                break
        det = determinant(selected_rows)
        assert det != 0
        records.append({
            "H": H,
            "columns": [f"sigma_k={k},d={H-k}" for k in range(1, H+1)]
                       + [f"Jordan_k={k},d={H-k}" for k in range(1, H+1)],
            "sample_primes": primes,
            "first_2H_rank": initial_rank,
            "verified_rank": rank(rows),
            "selected_primes": selected_primes,
            "determinant": encode(det),
            "matrix": [[encode(x) for x in row] for row in selected_rows],
        })
    return records


def falling(n, j):
    return prod(n-r for r in range(j))


def check_falling_factorials(max_degree=12):
    # S[d][j]: x^d = sum_j S(d,j) x^(falling j).
    # s[d][j]: x^(falling d) = sum_j s(d,j) x^j.
    S = [[0]*(max_degree+1) for _ in range(max_degree+1)]
    s = [[0]*(max_degree+1) for _ in range(max_degree+1)]
    S[0][0] = s[0][0] = 1
    for d in range(1, max_degree+1):
        for j in range(1, d+1):
            S[d][j] = S[d-1][j-1]+j*S[d-1][j]
            s[d][j] = s[d-1][j-1]-(d-1)*s[d-1][j]
    n = max_degree+1
    assert all(S[d][d] == s[d][d] == 1 for d in range(n))
    assert all(S[d][j] == s[d][j] == 0 for d in range(n) for j in range(d+1, n))
    for A, B in ((S, s), (s, S)):
        assert all(sum(A[i][k]*B[k][j] for k in range(n)) == int(i == j)
                   for i in range(n) for j in range(n))
    for d in range(n):
        for x in range(21):
            assert x**d == sum(S[d][j]*falling(x, j) for j in range(n))
            assert falling(x, d) == sum(s[d][j]*x**j for j in range(n))
    return {"maximum_degree": max_degree,
            "triangular_diagonal": [1]*n,
            "both_inverse_matrix_products": "identity",
            "integer_evaluation_checks": 2*n*21,
            "Stirling_second_kind": S, "signed_Stirling_first_kind": s}


def valuation(a, p):
    assert a > 0
    v = 0
    while a % p == 0:
        v += 1
        a //= p
    return v


def normalized_mean(kind, k, Q, a):
    """Return mu_sigma/zeta(k+1) or mu_Jordan*zeta(k+1).

    Only finitely many local Euler factors differ from the unrestricted
    mean.  Thus this value is an exact rational number.
    """
    assert k >= 1 and Q >= 1 and a >= 1
    result = F(1)
    for p, vQ in factor(Q).items():
        z = F(1, p)
        if kind == "sigma":
            va = valuation(a, p)
            if va < vQ:
                local = sum(z**(k*u) for u in range(va+1))
            else:
                local = sum(z**(k*u) for u in range(vQ))
                local += z**(k*vQ)/(1-z**(k+1))
            result *= local*(1-z**(k+1))
        elif kind == "Jordan":
            local = 1-z**k if a % p == 0 else F(1)
            result *= local/(1-z**(k+1))
        else:
            raise ValueError(kind)
    return result


def check_local_means():
    records = []
    for k in range(1, 7):
        for Q in (1, 6, 12, 30, 72):
            ell = primes_after(max(Q, 13), 1)[0]
            z = F(1, ell)
            for a in sorted(set((1, 2, 3, 5, 6, 12, 18, 36, 72, Q, 2*Q))):
                off = a+Q*(((1-a)*pow(Q, -1, ell)) % ell)
                on = a+Q*(((-a)*pow(Q, -1, ell)) % ell)
                assert off > 0 and on > 0
                assert off % ell == 1 and on % ell == 0
                assert off % Q == a % Q and on % Q == a % Q
                for kind in ("sigma", "Jordan"):
                    base = normalized_mean(kind, k, Q, a)
                    miss = normalized_mean(kind, k, Q*ell, off)
                    hit = normalized_mean(kind, k, Q*ell, on)
                    assert base > 0 and miss > 0 and hit > 0
                    if kind == "sigma":
                        expected_off = 1-z**(k+1)
                        expected_on = expected_off+z**k
                        response = z**k
                    else:
                        expected_off = 1/(1-z**(k+1))
                        expected_on = (1-z**k)/(1-z**(k+1))
                        response = -z**k/(1-z**(k+1))
                    assert miss/base == expected_off
                    assert hit/base == expected_on
                    assert (hit-miss)/base == response
                    records.append({"kind": kind, "k": k, "Q": Q,
                                    "a": a, "fresh_prime": ell,
                                    "a_off": off, "a_on": on,
                                    "normalized_base_mean": encode(base),
                                    "relative_mean_difference": encode(response)})
    return records


def poly_trim(a):
    a = list(a)
    while len(a) > 1 and a[-1] == 0:
        a.pop()
    return tuple(a)


def poly_add(a, b):
    return poly_trim([(a[i] if i < len(a) else 0)
                      +(b[i] if i < len(b) else 0)
                      for i in range(max(len(a), len(b)))])


def poly_scale(c, a):
    return poly_trim([c*x for x in a])


def poly_mul(a, b):
    result = [0]*(len(a)+len(b)-1)
    for i, x in enumerate(a):
        for j, y in enumerate(b):
            result[i+j] += x*y
    return poly_trim(result)


def poly_eval(a, z):
    value = F(0)
    for c in reversed(a):
        value = value*z+c
    return value


EULER_POLYNOMIALS = {
    "phi": (1, -1),
    "psi": (1, 1),
    "J3": (1, 0, 0, -1),
    "phi_cubed": (1, -3, 3, -1),
    "psi_squared": (1, 2, 1),
}
MIXED_LABELS = ("sigma1", "phi", "psi", "J3", "phi_cubed", "psi_squared")


def mixed_response(kind, prime):
    z = F(1, prime)
    if kind == "sigma1":
        return z
    p = poly_eval(EULER_POLYNOMIALS[kind], z)
    return (p-1)/(1+z*(p-1))


def mixed_local_factors(kind, prime):
    """Fresh-prime off/on means divided by the unrefined mean."""
    z = F(1, prime)
    if kind == "sigma1":
        off = 1-z*z
        return off, off+z
    p = poly_eval(EULER_POLYNOMIALS[kind], z)
    denominator = 1+z*(p-1)
    return 1/denominator, p/denominator


def mixed_contrast(kind, primes):
    """Exact on/off subset expansion, with only these primes refined.

    Each subset of primes uses its own refinement modulus.  In particular,
    this does not incorrectly introduce off-factors at unused primes.
    """
    assert primes
    factors = [mixed_local_factors(kind, p) for p in primes]
    result = F(0)
    constant_contrast = 0
    for mask in range(1 << len(primes)):
        hits = mask.bit_count()
        sign = (-1)**(len(primes)-hits)
        constant_contrast += sign
        result += sign*prod(pair[(mask >> j) & 1]
                            for j, pair in enumerate(factors))
    assert constant_contrast == 0
    return result


def check_mixed_prime_interpolation():
    # Verify psi = -phi + 2 J3 by clearing denominators as polynomials.
    # The identity being checked is psi+phi-2J3=0.
    n_phi, d_phi = (0, -1), (1, 0, -1)
    n_psi, d_psi = (0, 1), (1, 0, 1)
    n_J3, d_J3 = (0, 0, 0, -1), (1, 0, 0, 0, -1)
    numerator = poly_add(
        poly_add(poly_mul(n_psi, poly_mul(d_phi, d_J3)),
                 poly_mul(n_phi, poly_mul(d_psi, d_J3))),
        poly_scale(-2, poly_mul(n_J3, poly_mul(d_phi, d_psi))))
    assert numerator == (0,)

    # Canonical polynomial Euler data must identify phi*psi with J2.
    J2 = (1, 0, -1)
    assert poly_mul(EULER_POLYNOMIALS["phi"], EULER_POLYNOMIALS["psi"]) == J2
    assert poly_mul(poly_mul((1, -1), (1, -1)), (1, -1)) == EULER_POLYNOMIALS["phi_cubed"]
    assert poly_mul((1, 1), (1, 1)) == EULER_POLYNOMIALS["psi_squared"]
    canonical_input = {"phi_times_psi": poly_mul((1, -1), (1, 1)), "J2": J2}
    canonical_classes = {}
    for label, polynomial in canonical_input.items():
        canonical_classes.setdefault(poly_trim(polynomial), []).append(label)
    assert len(canonical_classes) == 1
    assert len(set(EULER_POLYNOMIALS.values())) == len(EULER_POLYNOMIALS)

    sample_primes = primes_after(2, 12)
    one_prime_rows = [[mixed_response(label, p) for label in MIXED_LABELS]
                      for p in sample_primes]
    dependent_rows = [[mixed_response(label, p) for label in ("phi", "psi", "J3")]
                      for p in sample_primes]
    assert rank(dependent_rows) == 2
    assert rank(one_prime_rows) == 5
    for p in sample_primes:
        assert mixed_response("psi", p) == -mixed_response("phi", p)+2*mixed_response("J3", p)
        for label in MIXED_LABELS:
            off, on = mixed_local_factors(label, p)
            assert off > 0 and on > 0
            assert on-off == mixed_response(label, p)

    targets = []
    isolation_matrix = []
    contrast_checks = 0
    for target in MIXED_LABELS:
        base_prime = 3
        assert mixed_response(target, base_prime) != 0
        others = [s for s in MIXED_LABELS if s != target]
        separators = []
        previous = base_prime
        for other in others:
            p = primes_after(previous, 1)[0]
            while mixed_response(target, p) == mixed_response(other, p):
                p = primes_after(p, 1)[0]
            separators.append((other, p))
            previous = p
        assert len({base_prime, *(p for _, p in separators)}) == len(MIXED_LABELS)

        # Expand x_base * product_other (x_p(other)-delta_other(p(other))).
        terms = []
        for mask in range(1 << len(separators)):
            coefficient = F(1)
            primes = [base_prime]
            for j, (other, p) in enumerate(separators):
                if (mask >> j) & 1:
                    primes.append(p)
                else:
                    coefficient *= -mixed_response(other, p)
            terms.append((coefficient, primes))

        row = []
        for tested in MIXED_LABELS:
            expanded = sum(coefficient*prod(mixed_response(tested, p) for p in primes)
                           for coefficient, primes in terms)
            factored = mixed_response(tested, base_prime)*prod(
                mixed_response(tested, p)-mixed_response(other, p)
                for other, p in separators)
            assert expanded == factored
            reconstructed = F(0)
            for coefficient, primes in terms:
                contrast = mixed_contrast(tested, primes)
                assert contrast == prod(mixed_response(tested, p) for p in primes)
                reconstructed += coefficient*contrast
                contrast_checks += 1
            assert reconstructed == expanded
            assert (expanded != 0) == (tested == target)
            row.append(expanded)
        isolation_matrix.append(row)
        targets.append({
            "target": target, "base_prime": base_prime,
            "separating_primes": [{"excluded_channel": s, "prime": p}
                                  for s, p in separators],
            "subset_expansion": [{"coefficient": encode(c), "primes": primes}
                                 for c, primes in terms],
            "values_on_all_channels": [encode(x) for x in row],
        })
    assert rank(isolation_matrix) == len(MIXED_LABELS)
    assert determinant(isolation_matrix) != 0
    return {
        "scope": "Six exact finite examples; not an all-function proof",
        "channels": list(MIXED_LABELS),
        "Euler_polynomials_ascending_coefficients": EULER_POLYNOMIALS,
        "one_prime_identity": "delta_psi = -delta_phi + 2 delta_J3",
        "cleared_polynomial_numerator": list(numerator),
        "one_prime_sample_primes": sample_primes,
        "phi_psi_J3_one_prime_rank": rank(dependent_rows),
        "all_six_one_prime_rank": rank(one_prime_rows),
        "mixed_prime_isolation_rank": rank(isolation_matrix),
        "nonempty_subset_contrast_checks": contrast_checks,
        "canonical_duplicate_check": {
            "identity": "P_phi * P_psi = P_J2",
            "canonical_coefficients": list(J2),
            "duplicate_labels": ["phi_times_psi", "J2"],
            "distinct_canonical_channels": len(canonical_classes),
            "policy": "These two labels must never be counted as independent channels",
        },
        "targets": targets,
        "isolation_matrix": [[encode(x) for x in row] for row in isolation_matrix],
        "isolation_determinant": encode(determinant(isolation_matrix)),
        "modulus_convention": "Each subset refines only its own listed primes",
    }


def main():
    responses = check_responses()
    falling_transform = check_falling_factorials()
    local_means = check_local_means()
    mixed_interpolation = check_mixed_prime_interpolation()
    output = {
        "scope": "Exact finite audits; not a proof of the all-degree theorem",
        "arithmetic": "Python standard-library Fraction and exact integers",
        "response_matrices": responses,
        "falling_factorial_transform": falling_transform,
        "fresh_prime_local_mean_checks": local_means,
        "mixed_prime_interpolation": mixed_interpolation,
        "all_checks": "passed",
    }
    path = Path(__file__).with_name("prime_response_checks.json")
    path.write_text(json.dumps(output, indent=2)+"\n")
    print(json.dumps({
        "response_ranks": [{"H": r["H"], "rank": r["verified_rank"],
                            "first_2H_rank": r["first_2H_rank"],
                            "primes": r["sample_primes"]} for r in responses],
        "falling_factorial_maximum_degree": falling_transform["maximum_degree"],
        "falling_factorial_evaluation_checks": falling_transform["integer_evaluation_checks"],
        "local_mean_checks": len(local_means),
        "mixed_one_prime_rank": mixed_interpolation["all_six_one_prime_rank"],
        "mixed_prime_isolation_rank": mixed_interpolation["mixed_prime_isolation_rank"],
        "mixed_subset_contrast_checks": mixed_interpolation["nonempty_subset_contrast_checks"],
        "all_checks": "passed", "json": str(path),
    }, indent=2))


if __name__ == "__main__":
    main()

\end{filecontents*}
\begin{filecontents*}[overwrite]{check_joint_grid.py}
#!/usr/bin/env python3
"""Exact finite audit of the K=2 joint-divisor factorial-tail grid.

Uses only the Python standard library.  This verifies one auxiliary grid;
it is not a numerical proof of infinite-series linear independence.
"""
from fractions import Fraction as F
from itertools import product
from math import gcd, lcm, prod
from pathlib import Path
import json


def factor(n):
    out = {}
    d = 2
    while d * d <= n:
        while n % d == 0:
            out[d] = out.get(d, 0) + 1
            n //= d
        d = 3 if d == 2 else d + 2
    if n > 1:
        out[n] = out.get(n, 0) + 1
    return out


def sigma(fac, k):
    return prod(sum(p ** (k*j) for j in range(v+1)) for p,v in fac.items())


def nullspace(rows):
    a = [[F(x) for x in row] for row in rows]
    nr, nc = len(a), len(a[0])
    pivots = []
    r = 0
    for c in range(nc):
        pivot = next((j for j in range(r,nr) if a[j][c]), None)
        if pivot is None:
            continue
        a[r],a[pivot] = a[pivot],a[r]
        val = a[r][c]
        a[r] = [x/val for x in a[r]]
        for j in range(nr):
            if j != r and a[j][c]:
                val = a[j][c]
                a[j] = [x-val*y for x,y in zip(a[j],a[r])]
        pivots.append(c)
        r += 1
        if r == nr:
            break
    free = [c for c in range(nc) if c not in pivots]
    vectors = []
    for c in free:
        v = [F(0)]*nc
        v[c] = F(1)
        for row,pivot in enumerate(pivots):
            v[pivot] = -a[row][c]
        vectors.append(v)
    return pivots, vectors


def is_prime(n):
    # Exact trial division; no probable-prime assertion is used.
    return n >= 2 and factor(n) == {n: 1}


def as_json_fraction(x):
    x = F(x)
    return {"numerator":str(x.numerator), "denominator":str(x.denominator)}


def main():
    # General-construction K=1 instance: p=(5,7), t=(0,2), w=(21,-20).
    assert F(21*5,6)-F(20*7,8)==0
    assert [h*pp-tt for h in (1,2) for pp,tt in ((5,0),(7,2))]==[5,5,10,12]
    assert 100%25==0 and 100%49==2
    K,L,B = 2,4,7
    E = list(product(range(L), repeat=K))
    I = (L-1)*(1+B)
    D = prod(p for p in range(2,max(I,K+1)+1) if is_prime(p))
    P = 1+K*D*I
    p = {e:P+D*(e[0]+B*e[1]) for e in E}
    t = {e:D*(e[0]+2*B*e[1]) for e in E}
    fac = {e:factor(p[e]) for e in E}
    sig = {(k,e):sigma(fac[e],k) for k in (1,2) for e in E}
    assert all(prod(q**v for q,v in fac[e].items()) == p[e] for e in E)
    assert all(is_prime(q) for e in E for q in fac[e])
    assert all(gcd(p[e],p[f])==1 for ix,e in enumerate(E) for f in E[ix+1:])
    assert all(gcd(p[e],h)==1 for e in E for h in range(1,K+2))

    # k=2,i=1,h=1 has S(1,0)=0 and is omitted legitimately.
    types = [(1,0,1),(2,0,1),(2,1,2)]
    rows=[]
    labels=[]
    for k,i,h in types:
        for fixed in range(L):
            labels.append([k,i,h,fixed])
            rows.append([F(p[e]**(i+1),sig[k,e]) if e[2-h]==fixed else F(0) for e in E])
    pivots,ns = nullspace(rows)
    assert ns
    v=ns[0]
    denominator=lcm(*(x.denominator for x in v))
    w=[int(x*denominator) for x in v]
    common=gcd(*w)
    w=[x//common for x in w]
    assert any(w)
    assert all(sum(x*y for x,y in zip(row,w))==0 for row in rows)
    shifts={(e,h):h*p[e]-t[e] for e in E for h in range(1,K+2)}
    assert all(r>0 for r in shifts.values())
    assert max(r for (e,h),r in shifts.items() if h<=K)<(K+1)*P
    assert min(shifts[e,K+1] for e in E)==(K+1)*P
    assert len({shifts[e,K+1] for e in E})==len(E)
    active=[e for e,x in zip(E,w) if x]
    star_e=active[0]
    star=shifts[star_e,K+1]
    assert [(e,h) for (e,h),r in shifts.items() if r==star]==[(star_e,K+1)]
    for h in range(1,K+1):
        for fixed in range(L):
            assert len({shifts[e,h] for e in E if e[2-h]==fixed})==1

    # Exact simultaneous residue class and the required multiplicativity.
    Q=prod(p[e]**2 for e in E)
    A=sum(t[e]*(Q//p[e]**2)*pow(Q//p[e]**2,-1,p[e]**2) for e in E)%Q
    assert all((A-t[e])%p[e]**2==0 for e in E)
    N=A+Q
    for e in E:
        n=(N-t[e])//p[e]
        assert n>0 and n%p[e]==0
        for h in range(1,K+2):
            assert p[e]*(n+h)==N+shifts[e,h]
            assert gcd(p[e],n+h)==1

    # Compute mu_{Q,2}(A+star)/zeta(3) exactly through Euler factors.
    qfac={q:2*v for e in E for q,v in fac[e].items()}
    assert prod(q**v for q,v in qfac.items())==Q
    ratio=F(1)
    for q,vQ in qfac.items():
        a=A+star
        va=0
        while a%q==0:
            va+=1
            a//=q
        if va<vQ:
            local=sum(F(1,q**(K*j)) for j in range(va+1))
        else:
            local=sum(F(1,q**(K*j)) for j in range(vQ))
            local+=F(1,q**(K*vQ))/(1-F(1,q**(K+1)))
        ratio*=local*(1-F(1,q**(K+1)))
    assert ratio>0
    ell=max(shifts.values())+1
    while not is_prime(ell):
        ell+=1
    assert ell<2**64 and gcd(ell,Q)==1
    v0=(-A*pow(Q,-1,ell))%ell
    v1=(-(A+star)*pow(Q,-1,ell))%ell
    assert all((Q*v0+A+r)%ell!=0 for r in shifts.values())
    assert all(((Q*v1+A+r)%ell==0)==(r==star) for r in shifts.values())
    top_coeff=F(w[E.index(star_e)]*p[star_e]**(K+1),sig[K,star_e])
    diff_ratio=top_coeff*ratio/F(ell**K)
    assert diff_ratio!=0
    out={
        "scope":"One exact K=2 auxiliary grid; not a proof by finite computation",
        "parameters":{"K":K,"L":L,"B":B,"I":I,"D":D,"P":P},
        "rows":len(rows),"columns":len(E),"rank":len(pivots),"nullity":len(ns),
        "constraint_labels_k_i_h_fixed":labels,
        "vertices":[{"e":e,"p":p[e],"t":t[e],"factorization":fac[e],
            "sigma1":sig[1,e],"sigma2":sig[2,e],"weight":str(w[j])}
            for j,e in enumerate(E)],
        "isolated_endpoint":{"vertex":star_e,"shift":star,"top_coefficient":as_json_fraction(top_coeff)},
        "crt":{"modulus":str(Q),"residue":str(A)},
        "fresh_prime":{"ell":ell,"v_miss":v0,"v_hit":v1,
            "mean_difference_divided_by_zeta3":as_json_fraction(diff_ratio)},
        "all_checks":"passed"
    }
    target=Path(__file__).with_name("joint_grid_K2.json")
    target.write_text(json.dumps(out,indent=2)+"\n")
    print(json.dumps({k:out[k] for k in ("parameters","rows","columns","rank","nullity","all_checks")},indent=2))
    print("weight_support_size",len(active))
    print("max_weight_digits",max(len(str(abs(x))) for x in w))
    print("isolated_vertex",star_e,"isolated_shift",star)
    print("fresh_prime",ell)
    print("json",target)


if __name__=="__main__":
    main()
\end{filecontents*}
\documentclass[11pt,a4paper]{article}
\usepackage[T1]{fontenc}
\usepackage[utf8]{inputenc}
\usepackage{lmodern,microtype}
\usepackage[margin=27mm]{geometry}
\usepackage{amsmath,amssymb,amsthm,mathtools,mathrsfs}
\usepackage{booktabs,array,longtable,enumitem,xcolor}
\usepackage{fancyhdr,listings,xurl}
\usepackage[colorlinks=true,linkcolor=blue!40!black,citecolor=blue!40!black,urlcolor=blue!40!black]{hyperref}
\hypersetup{pdftitle={Simultaneous Fiber Cancellation: Arithmetic Convolutions, Signed Sequences, and the Full Factorial Hierarchy},pdfauthor={KEIO}}
\pagestyle{fancy}\fancyhf{}
\fancyhead[L]{\small Factorial series and local prime responses}
\fancyhead[R]{\small KEIO}
\fancyfoot[C]{\thepage}
\setlength{\headheight}{14pt}
\setlength{\emergencystretch}{2em}
\setlist{itemsep=3pt,topsep=5pt}
\allowdisplaybreaks
\newtheorem{theorem}{Theorem}[section]
\newtheorem{lemma}[theorem]{Lemma}
\newtheorem{proposition}[theorem]{Proposition}
\newtheorem{corollary}[theorem]{Corollary}
\theoremstyle{definition}
\newtheorem{problem}[theorem]{Auxiliary problem}
\newtheorem{definition}[theorem]{Definition}
\theoremstyle{remark}
\newtheorem{remark}[theorem]{Remark}
\newcommand{\Q}{\mathbb Q}
\newcommand{\Z}{\mathbb Z}
\newcommand{\N}{\mathbb N}
\newcommand{\Stir}[2]{\left\{\begin{matrix}#1\\#2\end{matrix}\right\}}
\newcommand{\V}{\mathcal V}
\DeclareMathOperator{\rank}{rank}
\lstset{basicstyle=\ttfamily\scriptsize,columns=fullflexible,keepspaces=true,breaklines=true,showstringspaces=false,frame=single,rulecolor=\color{black!15},xleftmargin=4pt,xrightmargin=4pt,aboveskip=8pt,belowskip=8pt}

\title{Simultaneous Fiber Cancellation\\[3pt]
\large Arithmetic convolutions, signed sequences, and the full factorial hierarchy}
\author{KEIO\thanks{Prepared with ChatGPT assistance. The arguments received separate AI reviews and exact finite arithmetic checks. These checks are not independent human peer review or proof-assistant verification.}}
\date{29 October 2026 --- expanded integrated edition}
\begin{document}
\maketitle

\begin{abstract}
We develop irrationality and linear-independence criteria for factorial
series at every nonzero rational point whose reduced numerator has
absolute value one or two. The original all-rank divisor-sum target of
Deajim and Siksek is retained, with the factorial-grid architecture
attributed to Tokengrinder. The expanded arithmetic theorem admits
polynomial dependence on prime valuations, every logarithmic mean order,
and closure under pointwise, Dirichlet, and unitary products. It includes
all weighted ordered factorizations, higher divisor functions, gcd moments,
finite-group and finite-ring counting functions, and rough-prime
total-multiplicity weights. All polynomial moments, residue restrictions,
and positive and negative derivative orders remain jointly independent
after identical normalized profiles are grouped. Rational cyclotomic
coordinates cover every root-of-unity order; actual four- and six-direction
values are independent over the Gaussian and Eisenstein fields.
A complementary elementary criterion characterizes rationality for every
integer coefficient sequence of growth $o(n)$, and gives complete
two-sided hierarchies for all bounded non-eventually-periodic sequences.
Exact normal forms classify nested polynomial integrals and quotients
of affine coordinate forms. A separately attributed classical
$E$-function argument treats all inverse exponential moments and arbitrary
repeated integer-pole weights at nonzero algebraic arguments.
The original proofs are preserved, the new arithmetic proofs are supplied
in full, and eight exact auxiliary verifiers are embedded in this single
source. The manuscript is unrefereed and not Lean-formalized. Historical
priority is not certified; the classical companion is not claimed as a
new transcendence result, and no transcendence assertion is made for the
nontrivial arithmetic-profile values.
\end{abstract}
\clearpage
\setcounter{tocdepth}{1}
\tableofcontents
\clearpage

\section{The unchanged historical target}
For positive integers $k,n$, define
\begin{equation}\label{eq:alpha}
 \sigma_k(n)=\sum_{d\mid n}d^k,
 \qquad \alpha_k=\sum_{n=1}^{\infty}\frac{\sigma_k(n)}{n!}.
\end{equation}
The individual irrationality question is the Erd\H{o}s--Kac problem,
recorded as Erd\H{o}s Problem~252. Deajim and Siksek~\cite{DS} studied
the stronger, specific target
\begin{equation}\label{eq:target}
 1,\alpha_1,\ldots,\alpha_K
 \quad\hbox{linearly independent over $\Q$.}
\end{equation}
Their criterion assumes Schinzel's hypothesis and was checked for
$K=50$. The exact target and this qualification are recorded in the
primary publisher/Warwick abstract and independently in Pratt's
introduction~\cite{Pratt}. We use the descriptive phrase
``the linear-independence question studied by Deajim and Siksek'':
we do not invent a numbered conjecture or claim that this phrase was
already an established eponym.

Tokengrinder's public development~\cite{TG} supplies a proof of
irrationality of each individual $\alpha_k$. The later extension
directory~\cite{TGgen} expressly distinguishes a polynomial perturbation
$\sigma_k(n)+P(n)$ from a polynomial in the divisors,
$\sum_{d\mid n}P(d)$, and identifies the latter, equivalently the joint
target~\eqref{eq:target}, as further work. Individual irrationality
alone does not imply joint independence.

\begin{theorem}\label{thm:main}
For every integer $K\geq1$, the numbers
$1,\alpha_1,\ldots,\alpha_K$ are linearly independent over $\Q$.
Equivalently, for every nonzero polynomial $P\in\Q[X]$ with $P(0)=0$,
\begin{equation}\label{eq:polydivisor}
 \sum_{n=1}^{\infty}\frac{\sum_{d\mid n}P(d)}{n!}\notin\Q.
\end{equation}
\end{theorem}

The proof neither changes the denominator $n!$ nor restricts the
divisor exponents or their rational coefficients. No prime-pattern
conjecture is assumed. The equivalence follows by writing
$P(X)=\sum_{k=1}^K a_kX^k$, interchanging a finite sum with an absolutely
convergent series, and clearing rational coefficients.

Theorem~\ref{thm:two-sided} is the common transfer mechanism.
The enlarged arithmetic input is supplied by
Theorems~\ref{thm:v5-log-profile-main} and~\ref{thm:v5-rough-product}.
The latter permits a rough total-multiplicity factor multiplied by any
admissible polynomial valuation profile. The separate sequence criterion
in Theorem~\ref{thm:v5-sublinear} admits arbitrary signed sublinear
integer coefficients, and Theorem~\ref{thm:v5-aperiodic-hierarchy}
handles every bounded non-eventually-periodic sequence.
These statements have different hypotheses: their families are combined
only where an explicit joint theorem proves the combination, such as
Theorem~\ref{thm:v5-signed-block}.
All earlier arguments are retained in full so that the historical target
and every new dependency can be checked separately.

The progression from the previous revision is substantive: specifying
an exact prime valuation distinguishes profiles invisible to a
divisibility-only test; preserving residues gives simultaneous signed
and oscillatory values; and negative grid coordinates plus extra
global constraints admit integration orders of arbitrary size.
The permitted evaluation magnitudes remain $2/b$, with $b\geq1$.


\subsection{How to use the expanded results}
For the arithmetic theorem, compute a normalized prime-power profile,
group equal profiles, and check the stated positivity, integrality, and
growth conditions. Theorem~\ref{thm:v5-log-profile-main} and
Section~\ref{v5:profile-universe} supply entire closed arithmetic classes.
One may then use any fixed residue class, all two-sided jets, and the
rational weights with simple negative-integer poles. The stronger
rough-prime product theorem has its own growth restriction and does not
authorize arbitrary products of rough factors.

For signed coefficients of growth $o(n)$, eventual equality of sequences
decides all rational relations between their factorial values. Bounded
aperiodic sequences automatically have independent full jet hierarchies.
For nested polynomial kernels, reduce to the unique normal form of
Theorem~\ref{thm:operator-normal-form}; for quotients of affine forms,
use the exact rank test of Theorem~\ref{thm:projective-rank}.
The classical inverse-moment result in
Section~\ref{sec:v5-classical-inverse} is a separate, stronger
transcendence theorem for the pure exponential channel. It is not a
mixed-independence assertion with the other arithmetic channels.

These distinctions prevent double counting. Two integrals with the same
rational coefficient multiplier are one coordinate, and two counting
functions with the same normalized profile belong to the same arithmetic
block. The concrete examples are instances of infinite families, not a
claim that each was a previously named open constant.

\subsection{What is added to the existing method}
The factorial expansion, squared-congruence alignment, integer limit,
and progression-mean separation are the architecture of~\cite{TG}.
The new step of this argument is to solve all normalized cancellation
equations simultaneously. Multiplicativity introduces the different
factors $1/\sigma_k(p)$ for different $k$; these factors cannot simply
be ignored or replaced by a common binomial weight. A sufficiently
wide finite grid instead supplies a rational kernel. Terminal-shift
injectivity makes every nonzero kernel vector usable.

The supplied corpus was inspected at commit
\nolinkurl{510a32c9a3fd2aa0b5c78d3f6812ded723b9ac8f}~\cite{Corpus}.
Its fixed-box relation methods in [013] and mixed congruence kernels in
[078] motivated looking for simultaneous relations. No result in an
unreviewed corpus manuscript is a hypothesis of the proof below.
The source of the factorial-grid architecture is~\cite{TG}, rather
than those Mahler-series papers.

\section{A finite auxiliary problem and its solution}
The finite question is useful independently of the asymptotic argument.
For $L,K\geq1$, put $\V=\{0,\ldots,L-1\}^K$. A coordinate-$j$
fiber is the set of vertices obtained by fixing every coordinate except
$e_j$, for $0\leq j<K$.

\begin{problem}\label{prob:fiber}
Given finitely many rational coefficient arrays on $\V$, each assigned
to one coordinate direction, can one choose nonzero integer vertex
weights whose weighted sum for each array vanishes on every fiber
in its assigned direction? Can such cancellation be arranged without
annihilating the final expression needed for an application?
\end{problem}

\begin{lemma}[Simultaneous fiber cancellation]\label{lem:kernel}
Suppose there are at most $C$ assigned rational arrays, and $L>C$.
There are integers $c_e$, not all zero, satisfying every prescribed
fiber equation. If a remaining expression has one distinct formal
label $r_e$ for each vertex and a nonzero multiplier $u_e$ at that
label, then $\sum_e c_eu_e[r_e]$ is nonzero.
\end{lemma}
\begin{proof}
There are $L^K$ unknowns and at most $CL^{K-1}$ homogeneous equations
over $\Q$. Rank-nullity gives a nonzero rational solution; multiplication
by a common denominator gives integer weights. If $c_{e_*}\ne0$,
the coefficient of the unique label $r_{e_*}$ is $c_{e_*}u_{e_*}\ne0$.
\end{proof}

This is elementary finite-dimensional linear algebra. Its role here is
specific: it replaces an explicit one-function finite difference by
simultaneous equations involving several different multiplicative
normalizations. The terminal labels in the application will be actual
integer shifts, and their distinctness will be proved rather than
assumed. No bound on the size of the integer weights is required:
the entire finite construction is fixed before the asymptotic variable
tends to infinity.

\section{Factorial tails with a sufficiently small remainder}
Let $\tau(m)$ count the positive divisors of $m$. Complementary-divisor
pairing gives $\tau(m)\leq2\sqrt m$. Consequently, for $k\geq1$,
\begin{equation}\label{eq:crude}
 \sigma_k(m)\leq 2m^{k+1/2},\qquad
 \rho_k(m):=\frac{\sigma_k(m)}{m^k}
      =\sum_{d\mid m}d^{-k}\leq2\sqrt m.
\end{equation}
These deliberately weak bounds are sufficient. In particular all the
series~\eqref{eq:alpha} converge absolutely.

Write $\Stir{i}{j}$ for the Stirling numbers of the second kind,
including $\Stir00=1$ and $\Stir{i}{0}=0$ for $i>0$.

\begin{lemma}[Common-order tail expansion]\label{lem:tail}
Fix $K\geq k\geq1$. As $n\to\infty$ through positive integers,
\begin{align}\label{eq:tail}
 T_k(n)&:=n!\sum_{m>n}\frac{\sigma_k(m)}{m!}\\
 &=\sum_{j=0}^{K}\sum_{i=j}^{K}
      \Stir{i}{j}\frac{\sigma_k(n+j+1)}{(n+j+1)^{i+1}}
       +O_K(n^{-3/2}).\nonumber
\end{align}
In particular the remainder remains $o(1)$ after multiplication by $n$.
\end{lemma}
\begin{proof}
For fixed $0\leq j\leq K$ and $x\to\infty$, the generating-function
identity
\[
 \prod_{r=1}^{j}(1-rz)^{-1}
       =\sum_{u=0}^{\infty}\Stir{u+j}{j}z^u
\]
gives
\begin{equation}\label{eq:stirling}
 \frac{1}{x(x-1)\cdots(x-j)}
   =\sum_{i=j}^{K}\Stir{i}{j}x^{-i-1}
       +O_K(x^{-K-2}).
\end{equation}
For $j=0$ the formula is exact. For $j>0$, the generating identity
follows either by multiplying the geometric series or from the
recurrence $\Stir{i+1}{j}=j\Stir{i}{j}+\Stir{i}{j-1}$.
Since the finitely many rational functions are analytic at $z=0$,
their Taylor remainders have the stated bounds.

Apply~\eqref{eq:stirling} to the first $K+1$ summands of $T_k(n)$,
with $x=n+j+1$. After multiplication by $\sigma_k(x)$, the remainder
is $O_K(n^{k-K-3/2})$, hence $O_K(n^{-3/2})$.
For the remaining summands $h\geq K+2$, use the majorants
\[
 U_h=\frac{2(n+h)^{k+1/2}}{(n+1)\cdots(n+h)}.
\]
Their consecutive ratio is
\[
 \frac{U_{h+1}}{U_h}
 =\frac{(1+1/(n+h))^{k+1/2}}{n+h+1}\leq\frac12
\]
for all sufficiently large $n$, uniformly in $h\geq K+2$.
Also $U_{K+2}=O_K(n^{k-K-3/2})$. The omitted tail is at most
$2U_{K+2}$. This proves~\eqref{eq:tail}.
\end{proof}

If integers $a_1,\ldots,a_K$ satisfy
$\alpha=\sum_k a_k\alpha_k\in\Q$, then
\begin{equation}\label{eq:integral}
 T(n):=\sum_{k=1}^{K}a_kT_k(n)
 =n!\alpha-\sum_{m=1}^{n}\frac{n!}{m!}
                         \sum_{k=1}^{K}a_k\sigma_k(m)
 \in\Z
\end{equation}
for all sufficiently large $n$. Indeed the denominator of $\alpha$
divides $n!$ once $n$ is sufficiently large. The coefficients $a_k$
may have either sign; positivity of the combined tail is unnecessary.

\section{A coprime grid with distinct terminal shifts}
Fix $K\geq1$ and set
\begin{equation}\label{eq:gridconstants}
\begin{gathered}
 C=\frac{K^2(K+1)}2,\qquad L=C+1,\qquad b=K(L-1)+1,\\
 F=\max(b^K-1,K+1),\qquad D=F!,\qquad B=1+KDF.
\end{gathered}
\end{equation}
For $e=(e_0,\ldots,e_{K-1})\in\V=\{0,\ldots,L-1\}^K$, define
\begin{align}\label{eq:grid}
 I_e&=\sum_{a=0}^{K-1}e_ab^a,&
 J_e&=\sum_{a=0}^{K-1}(a+1)e_ab^a,\\
 p_e&=B+DI_e,&t_e&=DJ_e,\nonumber\\
 s_{e,j}&=(j+1)p_e-t_e&& (0\leq j\leq K).\nonumber
\end{align}
The notation $p_e$ denotes positive integer multipliers; they need not
be primes.

\begin{lemma}[Grid geometry and arithmetic]\label{lem:grid}
The construction has the following properties.
\begin{enumerate}
\item The $p_e$ are pairwise coprime, and $\gcd(p_e,h)=1$ for
every $1\leq h\leq K+1$.
\item $0\leq t_e<B\leq p_e$, so all $s_{e,j}$ are positive.
\item For $j<K$, $s_{e,j}$ is independent of $e_j$ along its fiber.
\item All terminal shifts $s_{e,K}$ are distinct. Each is at least
$(K+1)B$, whereas every $s_{e,j}$ with $j<K$ is smaller than $(K+1)B$.
\end{enumerate}
\end{lemma}
\begin{proof}
Base-$b$ encoding makes the $I_e$ distinct and gives $0\leq I_e\leq F$.
Also $J_e\leq KI_e$, whence $t_e\leq KDF=B-1$.
Every $p_e$ is congruent to $1$ modulo $D$. If a prime $\ell$ divided
two distinct multipliers, it would divide $D(I_e-I_f)$ but not $D$.
It would therefore divide the nonzero integer $I_e-I_f$, of absolute
value at most $F$. Hence $\ell\leq F$ and $\ell\mid F!=D$, a
contradiction. Since $F\geq K+1$, every $h\leq K+1$ divides $D$,
giving the other coprimality assertion.

Expanding the shifts gives
\begin{equation}\label{eq:shift}
 s_{e,j}=(j+1)B+D\sum_{a=0}^{K-1}(j-a)e_ab^a.
\end{equation}
For $j<K$, the coefficient of $e_j$ is zero. At $j=K$, the digits
$(K-a)e_a$ lie between $0$ and $K(L-1)=b-1$. Uniqueness of base-$b$
digits makes $e\mapsto s_{e,K}$ injective; those digits are
nonnegative, so $s_{e,K}\geq(K+1)B$.
For $j<K$ we have
\[
 s_{e,j}\leq Kp_e\leq K(B+DF)=(K+1)B-1.
\]
This also proves separation from every lower shift.
\end{proof}

\begin{lemma}[Arithmetic cancellation weights]\label{lem:weights}
There are integers $c_e$, not all zero, such that for
$1\leq k\leq K$, $0\leq i<K$, $0\leq j\leq i$, and each
coordinate-$j$ fiber,
\begin{equation}\label{eq:constraints}
 \sum_{e_j=0}^{L-1}c_e\frac{p_e^{i+1}}{\sigma_k(p_e)}=0.
\end{equation}
Consequently, for every function $g$ on the positive integers,
\begin{equation}\label{eq:fibercancel}
 \sum_{e\in\V}c_e\frac{p_e^{i+1}}{\sigma_k(p_e)}g(s_{e,j})=0
\end{equation}
for the same triples $(k,i,j)$.
\end{lemma}
\begin{proof}
There are $K\sum_{i=0}^{K-1}(i+1)=C$ rational coefficient arrays.
Each denominator $\sigma_k(p_e)$ is a positive integer. Apply
Lemma~\ref{lem:kernel}, since $L>C$. On each fiber the shift
$s_{e,j}$ is fixed, so multiplication of~\eqref{eq:constraints}
by $g(s_{e,j})$ and summation over the fibers proves
\eqref{eq:fibercancel}.
\end{proof}

\section{Progression means distinguish an isolated shift}
We first prove the final obstruction independently of the grid.
For $k,Q,a\geq1$, put
\begin{equation}\label{eq:mean}
 \mu_k(Q,a)=
 \sum_{\substack{d\geq1\\\gcd(d,Q)\mid a}}
          \frac{\gcd(d,Q)}{d^{k+1}}.
\end{equation}
It is finite, depends only on $a\bmod Q$, and is positive because
its $d=1$ summand is $1$.

\begin{lemma}[Exact progression average]\label{lem:average}
For positive integers $k,Q,a$,
\begin{equation}\label{eq:cesaro}
 \lim_{M\to\infty}\frac1M\sum_{v=0}^{M-1}\rho_k(a+Qv)
       =\mu_k(Q,a).
\end{equation}
If $\ell\nmid Q$ is prime and $a'\equiv a\pmod Q$, then
\begin{equation}\label{eq:refined}
 \mu_k(Q\ell,a')=
 \bigl(1-\ell^{-k-1}+\mathbf1_{\ell\mid a'}\ell^{-k}\bigr)
              \mu_k(Q,a).
\end{equation}
\end{lemma}
\begin{proof}
Use $\rho_k(m)=\sum_{d\mid m}d^{-k}$. For each fixed $d$, the
congruence $d\mid a+Qv$ has no solution unless $\gcd(d,Q)\mid a$;
otherwise its density is $\gcd(d,Q)/d$. Truncating the divisor sum
at $H$ therefore gives the desired limit for the truncated expression.
The average of the omitted positive terms is at most
\[
 Q\sum_{d>H}d^{-k-1}
       +\frac1M\sum_{d\leq a+QM}d^{-k}.
\]
The second expression tends to zero, including when $k=1$, when it
is $O(\log M/M)$. The first tends to zero as $H\to\infty$.
This justifies the limit without assuming that $\rho_1$ is bounded.

To prove~\eqref{eq:refined}, write $d=\ell^u d_0$ with
$\ell\nmid d_0$. In~\eqref{eq:mean} for modulus $Q$, the sum over
$u\geq0$ contributes $(1-\ell^{-k-1})^{-1}$. For modulus $Q\ell$,
only $u=0$ is allowed if $\ell\nmid a'$. If $\ell\mid a'$, the
terms $u\geq1$ instead have the additional factor $\ell$ from
$\gcd(d,Q\ell)$. Thus the additional factor, relative to the
original mean, is
\[
 (1-\ell^{-k-1})\sum_{u\geq1}\ell^{1-u(k+1)}=\ell^{-k}.
\]
The condition involving $d_0$ is unchanged because $a'\equiv a\pmod Q$.
\end{proof}

\begin{proposition}[No constant limit with an isolated shift]\label{prop:nolimit}
Fix $k,Q\geq1$, an integer $A$, finitely many integer shifts $r_i$,
and real coefficients $\beta_i$. Suppose one shift $r_*$ occurs
exactly once and its coefficient $\beta_*$ is nonzero. Then
\[
 \sum_i\beta_i\rho_k(A+Qv+r_i)
\]
cannot converge to any finite constant as $v\to\infty$.
Only sufficiently large $v$, for which all arguments are positive,
are relevant.
\end{proposition}
\begin{proof}
Choose a prime $\ell$ larger than $Q$, the number of distinct shifts,
and all differences $|r_i-r_j|$. Distinct shifts are then distinct
modulo $\ell$. Choose one residue for $N\bmod\ell$ missing all
$-r_i$, and another residue equal to $-r_*$. By the Chinese
remainder theorem both can be combined with $N\equiv A\pmod Q$.
Take positive representatives $A_{\rm off},A_{\rm on}$ large enough
for all shifted arguments to be positive.

If the asserted finite limit existed, its value would also be the
Ces\`aro mean along both refined progressions. By
Lemma~\ref{lem:average}, however, the difference between those two
means is exactly
\begin{equation}\label{eq:difference}
 \beta_*\ell^{-k}\mu_k(Q,A+r_*)\ne0.
\end{equation}
Here a positive representative is used in the mean if needed.
Every other shift has the same factor $1-\ell^{-k-1}$ in both
progressions, while only $r_*$ acquires the extra factor
$\ell^{-k}$. Positivity of the original mean gives the contradiction.
\end{proof}

\section{Proof of the joint linear-independence theorem}
Suppose a nontrivial rational relation with $1$ exists. Clear its
denominators and choose the largest index $K$ with a nonzero
coefficient. We then have
\[
 \alpha=\sum_{k=1}^{K}a_k\alpha_k\in\Q,
 \qquad a_k\in\Z,\quad a_K\ne0.
\]
Construct the grid and integer weights from
Lemmas~\ref{lem:grid} and~\ref{lem:weights}. Put
\[
 Q=\prod_{e\in\V}p_e^2.
\]
Pairwise coprimality supplies an integer $A$ with
\begin{equation}\label{eq:crt}
 A\equiv t_e\pmod{p_e^2}\qquad(e\in\V).
\end{equation}
Let $N=A+Qv$ tend to infinity and set
\[
 n_e=\frac{N-t_e}{p_e}.
\]
These are eventually positive integers, divisible by $p_e$. Hence,
for $0\leq j\leq K$,
\[
 \gcd(p_e,n_e+j+1)=\gcd(p_e,j+1)=1,
 \qquad p_e(n_e+j+1)=N+s_{e,j}.
\]
The divisor-sum function is multiplicative on coprime arguments:
each divisor of a product of coprime integers decomposes uniquely
as a product of divisors. It follows exactly that
\begin{equation}\label{eq:rescale}
 \frac{\sigma_k(n_e+j+1)}{(n_e+j+1)^{i+1}}
  =\frac{p_e^{i+1}}{\sigma_k(p_e)}
       \frac{\sigma_k(N+s_{e,j})}{(N+s_{e,j})^{i+1}}.
\end{equation}

By~\eqref{eq:integral}, the quantity
\[
 W(N)=\sum_{e\in\V}c_eT(n_e)
\]
is an integer for all sufficiently large $N$ in this progression.
Substitute the common-order expansion~\eqref{eq:tail} and
the rescaling~\eqref{eq:rescale}. For each $i<K$, the sum over
$e$ is zero by~\eqref{eq:fibercancel}, using
$g(s)=\sigma_k(N+s)/(N+s)^{i+1}$. We obtain
\begin{align}\label{eq:survivor}
 W(N)={}&\sum_{k=1}^{K}a_k\sum_{e\in\V}\sum_{j=0}^{K}
 \frac{c_ep_e^{K+1}\Stir{K}{j}}{\sigma_k(p_e)}
 \frac{\sigma_k(N+s_{e,j})}{(N+s_{e,j})^{K+1}}
   +O(N^{-3/2}).
\end{align}
The implicit constant can depend on the fixed grid, weights, and
coefficients. This causes no difficulty: $K$ and all those objects
are fixed before $N$ varies. In particular $n_e\asymp N$ for each
fixed vertex, so the scaled errors still have order $N^{-3/2}$.

By~\eqref{eq:crude}, each main term in~\eqref{eq:survivor} is
$O(N^{k-K-1/2})$, hence $O(N^{-1/2})$. Thus $W(N)\to0$.
An integer sequence converging to zero is eventually zero. Multiplying
\eqref{eq:survivor} by $N$, the error tends to zero, and the terms
with $k<K$ tend to zero because $k-K+1/2\leq-1/2$.
Define
\[
 \beta_{e,j}=\frac{a_Kc_ep_e^{K+1}\Stir{K}{j}}{\sigma_K(p_e)}.
\]
The remaining assertion is
\[
 \sum_{e,j}\beta_{e,j}
       \frac{N}{N+s_{e,j}}\rho_K(N+s_{e,j})\longrightarrow0.
\]
For every fixed shift $s$, replacing $N/(N+s)$ by $1$ changes a
summand by at most a constant times
$|s|\rho_K(N+s)/(N+s)=O(N^{-1/2})$. Consequently
\begin{equation}\label{eq:limit}
 \sum_{e\in\V}\sum_{j=0}^{K}
       \beta_{e,j}\rho_K(N+s_{e,j})\longrightarrow0.
\end{equation}

Choose $e_*$ with $c_{e_*}\ne0$. Since $\Stir{K}{K}=1$ and
$a_K\ne0$, the coefficient $\beta_{e_*,K}$ is nonzero.
Lemma~\ref{lem:grid} says that its shift $s_{e_*,K}$ occurs
nowhere else among the displayed summands. Thus~\eqref{eq:limit}
contradicts Proposition~\ref{prop:nolimit}. This proves
Theorem~\ref{thm:main}.\qed

\subsection{A further consequence involving the exponential constant}
\begin{corollary}\label{cor:e}
For every $K\geq1$, the numbers
$1,e,\alpha_1,\ldots,\alpha_K$ are linearly independent over $\Q$.
\end{corollary}
\begin{proof}
Suppose $\lambda e+\sum_{k=1}^{K}a_k\alpha_k\in\Q$ with integral
coefficients and with some $a_k\ne0$; let $K$ be the largest such
index. A constant numerator contributes to the combined factorial
tail the additional term
\[
 \lambda U(n),\qquad
 U(n)=n!\sum_{m>n}\frac1{m!}
     =\frac1{n+1}+O(n^{-2}).
\]
The combined tails are still eventually integral; the omitted
$m=0$ contribution is rational. Using the same grid and weights,
the additional contribution to $W(N)$ tends to zero, while after
multiplication by $N$ it tends to $\lambda\sum_ec_ep_e$.
The preceding proof therefore forces the expression in
\eqref{eq:limit} to converge to a finite constant instead of zero.
Proposition~\ref{prop:nolimit} excludes that as well.
If all $a_k=0$, a nontrivial relation would make $e$ rational.
The elementary factorial-tail proof rules this out: rationality of
$e$ would make $U(n)$ an integer eventually, whereas $0<U(n)<1$
for all sufficiently large $n$.
\end{proof}

This corollary is a consequence of the proved mechanism, not a
replacement for the original target. It does not assert algebraic
independence or transcendence of any particular $\alpha_k$.

\section{Several primes separate distinct local responses}\label{sec:transfer}
The preceding argument uses one fresh prime to separate a surviving
divisor sum. We now give a transfer principle in which several primes
separate different arithmetic functions at the same terminal shift.
This permits polynomial weights, Jordan and Dedekind functions, and
products of these functions. The separation is stronger than linear
independence of one-prime responses: distinct responses suffice.

For a function $\rho:\N\to\mathbb R$, write
\[
 M_\rho(Q,a)=\lim_{X\to\infty}\frac1X
        \sum_{v=0}^{X-1}\rho(Qv+a),\qquad Q,a\geq1,
\]
when the limit exists. This mean depends only on $a$ modulo $Q$.

\begin{lemma}[Separation by several fresh primes]\label{lem:multimean}
Let $\rho_1,\ldots,\rho_s$ have positive means $M_u(Q,a)$ on every
arithmetic progression. Suppose that for every prime
$\ell\nmid Q$, and $a'\equiv a\pmod Q$,
\begin{equation}\label{eq:localresponse}
 M_u(Q\ell,a')=M_u(Q,a)
 \left[b_u(1/\ell)+\mathbf1_{\ell\mid a'}\delta_u(1/\ell)\right].
\end{equation}
Here the same functions $b_u,\delta_u$ work for every $Q,a$, and
the $\delta_u$ are pairwise distinct, nonzero rational functions.
Consider a finite sum
\[
 V(N)=\sum_{u=1}^s\sum_r\beta_{u,r}\rho_u(N+r)
\]
on $N\equiv A\pmod Q$. If a shift $r_*$ occurs at a distinguished
location, and some $\beta_{u,r_*}\ne0$, then $V(N)$ cannot have
a finite limit. Repeated shifts in the sum are first grouped together.
\end{lemma}
\begin{proof}
Suppose that the limit is $L$. Call a prime admissible if it exceeds
$Q$, the number of distinct shifts, and every nonzero absolute
difference between two shifts, and avoids all poles required below.
At an admissible prime choose an off-residue avoiding every $-r$,
and an on-residue equal to $-r_*$; the on-residue hits no other shift.

For a nonempty finite set $S$ of admissible primes, CRT gives
$2^{|S|}$ refinements of $N\equiv A\pmod Q$, one for every choice
of on/off states. The mean of $V$ on each refinement equals $L$.
Take their alternating sum, with sign $(-1)^{\text{number of off states}}$.
Every other shift has an off-factor at every prime, independently
of the state, and cancels. Iterating~\eqref{eq:localresponse} at the
distinguished shift gives
\begin{equation}\label{eq:mixedresponse}
 \sum_{u=1}^s B_u\prod_{\ell\in S}\delta_u(1/\ell)=0,
 \qquad B_u=\beta_{u,r_*}M_u(Q,A+r_*).
\end{equation}
If necessary the starting residue is increased by a multiple of $Q$
to make all arguments positive. Crucially, each subset $S$ refines
only its own primes. Thus the coefficients $B_u$ are the same for
every subset; no unused off-factor is incorporated into them.

Fix an index $t$. Choose an admissible base prime $\ell_0$ with
$\delta_t(1/\ell_0)\ne0$. For each $u\ne t$, choose a distinct
admissible prime $\ell_u$, also different from $\ell_0$, such that
$\delta_t(1/\ell_u)\ne\delta_u(1/\ell_u)$.
These choices are possible because a nonzero rational function has
only finitely many zeros, while there are infinitely many primes.
Expand the product in
\begin{equation}\label{eq:interpolation}
 \sum_{v=1}^s B_v\delta_v(1/\ell_0)
   \prod_{u\ne t}
   \left[\delta_v(1/\ell_u)-\delta_u(1/\ell_u)\right]=0.
\end{equation}
Every term in its expansion is a constant multiple of an equation
\eqref{eq:mixedresponse} for a subset containing $\ell_0$, so the
displayed equality follows. For $v\ne t$, the factor indexed by
$u=v$ is zero. The multiplier of $B_t$ is nonzero. Hence $B_t=0$.
As $t$ was arbitrary and the means are positive, every
$\beta_{u,r_*}=0$, a contradiction.
\end{proof}

The hypothesis about the distinguished location is automatic for a
terminal grid shift: its vertex and shift index are unique, although
several arithmetic functions may contribute there. No linear
independence assumption on the rational functions $\delta_u$ is used.
For example,
\[
 \frac{z}{1+z^2}=\frac{z}{1-z^2}
                    -\frac{2z^3}{1-z^4}
\]
is a genuine one-prime linear dependence. It causes no difficulty for
Lemma~\ref{lem:multimean}, since the three functions are distinct.

\begin{theorem}[Transfer from local responses to factorial values]
\label{thm:transfer}
Let $\mathcal F$ be a collection of positive integer-valued
multiplicative functions. Assign each $f\in\mathcal F$ an integer
$h_f\geq1$ such that $f(n)=O_f(n^{h_f+1/2})$, and put
$\rho_f(n)=f(n)/n^{h_f}$. Suppose that these normalized functions
have positive progression means and satisfy~\eqref{eq:localresponse}.
Within each fixed degree $h$, their response functions are nonzero
and pairwise distinct rational functions.
Then, for every fixed $c\in\Z\setminus\{0\}$, the family
\begin{equation}\label{eq:transferfamily}
 \{1,e^{1/c}\}\ \cup\
 \left\{\sum_{n\geq1}\frac{f(n)}{c^n n!}:f\in\mathcal F\right\}
\end{equation}
is linearly independent over $\Q$. Independence of an infinite family
means independence of every finite subfamily.
\end{theorem}
\begin{proof}
It suffices to consider a finite proposed relation, with integral
coefficients $a_f$ and $\lambda$, in which
$\lambda e^{1/c}+\sum_f a_f\sum_{m\geq1}f(m)/(c^m m!)$ is rational.
If some $a_f$ is nonzero, retain just these functions and put
$H=\max h_f$, and let $s$ be their number. Define
\[
 T_{f,c}(n)=c^n n!\sum_{m>n}\frac{f(m)}{c^m m!}.
\]
The same falling-product expansion proved earlier, now with the
factor $c^{-j-1}$, gives
\begin{equation}\label{eq:generaltail}
 T_{f,c}(n)=\sum_{j=0}^H\sum_{i=j}^H
 \frac{\Stir{i}{j}}{c^{j+1}}
 \frac{f(n+j+1)}{(n+j+1)^{i+1}}+O(n^{-3/2}).
\end{equation}
Indeed each finite truncation error is
$O(n^{h_f+1/2-H-2})=O(n^{-3/2})$; the omitted tail is bounded
by a geometric series with first term of the same order.
The inequality $|c|\geq1$ preserves this bound even when $c<0$.

For the exponential term put
\[
 U_c(n)=c^n n!\sum_{m>n}\frac1{c^m m!}
       =\frac1{c(n+1)}+O(n^{-2}).
\]
The combined tail $T(n)=\lambda U_c(n)+\sum_fa_fT_{f,c}(n)$
is an integer for all large $n$. The scaled rational total is an
integer eventually, and every scaled partial-sum term is an integer,
since $c^{n-m}n!/m!\in\Z$ for $m\leq n$.

Use the grid in~\eqref{eq:gridconstants}--\eqref{eq:grid} with
$K=H$, but choose $L>sH(H+1)/2$. Its geometric proof works for
every such $L$, with $b=H(L-1)+1$ and the same formulas for
$F,D,B$. Impose, for every $f$, $0\leq i<H$, $0\leq j\leq i$,
and coordinate-$j$ fiber,
\begin{equation}\label{eq:generalweights}
 \sum_{e_j}w_e\frac{p_e^{i+1}}{f(p_e)}=0.
\end{equation}
There are at most $sH(H+1)L^{H-1}/2<L^H$ equations. Since all
denominators are positive integers, Lemma~\ref{lem:kernel} gives
nonzero integral weights $w_e$.

Choose $N\equiv A\pmod Q$, where $Q=\prod_ep_e^2$ and
$A\equiv t_e\pmod{p_e^2}$. Then $n_e=(N-t_e)/p_e$ is an integer,
$p_e\mid n_e$, and multiplicativity gives exactly
\[
 \frac{f(n_e+j+1)}{(n_e+j+1)^{i+1}}
 =\frac{p_e^{i+1}}{f(p_e)}
   \frac{f(N+s_{e,j})}{(N+s_{e,j})^{i+1}}.
\]
The fiber equations remove every $i<H$ term. The eventual integer
$W(N)=\sum_ew_eT(n_e)$ is therefore $O(N^{-1/2})$ and is eventually
zero. Multiplication by $N$ eliminates all functions with $h_f<H$
and the remainders, and forces
\begin{equation}\label{eq:generallimit}
 \sum_{h_f=H}\sum_e\sum_{j=0}^H
 \frac{a_fw_ep_e^{H+1}\Stir{H}{j}}{c^{j+1}f(p_e)}
       \rho_f(N+s_{e,j})
 \longrightarrow -\frac{\lambda}{c}\sum_ew_ep_e.
\end{equation}
Replacing $N/(N+s)$ by $1$ costs $O(N^{-1/2})$, because
$\rho_f(N+s)=O(N^{1/2})$.

Select a vertex $e_*$ with $w_{e_*}\ne0$. Its terminal shift
$s_{e_*,H}$ occurs at no other grid location. For each retained
top-degree function its coefficient there is
$a_fw_{e_*}p_{e_*}^{H+1}/(c^{H+1}f(p_{e_*}))\ne0$.
Lemma~\ref{lem:multimean} contradicts~\eqref{eq:generallimit}.
If there was no nonzero $a_f$, the relation contradicts the
irrationality of $e^{1/c}$: its $c$th power is $e$, whose elementary
irrationality was proved above. This completes the proof.
\end{proof}

\section{Polynomial Euler factors and all polynomial moments}
\label{sec:eulerpolynomials}
Let $\mathcal P$ consist of polynomials $P\in\Z[X]$ satisfying
\begin{equation}\label{eq:admissibleP}
 P(0)=1,\qquad P\ne1,\qquad P(1/p)>0
 \quad\hbox{for every prime }p.
\end{equation}
Positivity on $0<x\leq1/2$ is a convenient sufficient condition.
For $P\in\mathcal P$ put $h_P=\deg P$ and
\begin{equation}\label{eq:FP}
 F_P(n)=n^{h_P}\prod_{p\mid n}P(1/p).
\end{equation}
Distinct polynomials are used only once. This convention avoids
counting two factorizations of the same polynomial as different
arithmetic functions.

\begin{theorem}[All moments of polynomial Euler factors]
\label{thm:allmoments}
For every fixed nonzero integer $c$, the following entire family is
linearly independent over $\Q$:
\begin{equation}\label{eq:allmoments}
 \{1,e^{1/c}\}\ \cup\
 \left\{\sum_{n\geq1}\frac{n^d\sigma_k(n)}{c^n n!}:
                     k\geq1,\ d\geq0\right\}
 \ \cup\
 \left\{\sum_{n\geq1}\frac{n^dF_P(n)}{c^n n!}:
                     P\in\mathcal P,\ d\geq0\right\}.
\end{equation}
The statement is at one fixed $c$ at a time; it does not assert joint
independence across different evaluation points $1/c$.
\end{theorem}

\begin{lemma}[Growth and means for polynomial Euler factors]
\label{lem:Pmeans}
The function $F_P$ is positive integer-valued and multiplicative.
For every $\epsilon>0$,
$F_P(n)=O_{P,\epsilon}(n^{h_P+\epsilon})$. Its normalized phase
$\rho_P(n)=F_P(n)/n^{h_P}$ has the positive progression mean
\begin{equation}\label{eq:Pmean}
 M_P(Q,a)=
 \prod_{p\nmid Q}\left(1+\frac{P(1/p)-1}{p}\right)
 \prod_{p\mid\gcd(Q,a)}P(1/p).
\end{equation}
Its off-factor and response are
\begin{equation}\label{eq:Presponse}
 b_P(z)=\frac1{1+z(P(z)-1)},\qquad
 \delta_P(z)=\frac{P(z)-1}{1+z(P(z)-1)}.
\end{equation}
\end{lemma}
\begin{proof}
For $a\geq1$,
$F_P(p^a)=p^{h_P(a-1)}[p^{h_P}P(1/p)]$ is a positive integer.
This proves integrality and multiplicativity. Since
$P(1/p)=1+O_P(p^{-1})$, each sufficiently large prime factor
contributes at most $p^\epsilon$; the finitely many smaller primes
give a fixed constant. This proves the growth bound.

Define $a_P(1)=1$, set $a_P(d)=0$ for nonsquarefree $d$, and on
squarefree $d$ put
$a_P(d)=\prod_{p\mid d}(P(1/p)-1)$. Then
\[
 \rho_P(n)=\sum_{d\mid n}a_P(d),\qquad
 |a_P(d)|\ll_{P,\epsilon}d^{-1+\epsilon}.
\]
For the bound, use $|P(1/p)-1|\leq C_P/p$ and
$C_P^{\omega(d)}\ll_{P,\epsilon}d^\epsilon$ on squarefree integers;
the latter follows by separating the finitely many primes with
$p^\epsilon<C_P$.

Counting divisibility in an arithmetic progression now yields
\[
 M_P(Q,a)=\sum_{\gcd(d,Q)\mid a}
              a_P(d)\frac{\gcd(d,Q)}d.
\]
To justify passage to the limit, for a cutoff $Y$ bound the remaining
average in absolute value by
\[
 Q\sum_{d>Y}\frac{|a_P(d)|}{d}
       +\frac1X\sum_{d\leq QX+a}|a_P(d)|.
\]
Choose $0<\epsilon<1$. The second term is $O(X^{-1+\epsilon})$
and the first tends to zero with $Y$. The absolutely convergent
Euler product is exactly~\eqref{eq:Pmean}. Every factor is positive:
an unrestricted factor is $(1-1/p)+P(1/p)/p$. Since it is
$1+O_P(p^{-2})$, the infinite product is strictly positive.

At a new prime, replace the unrestricted local factor
$1+(P(1/p)-1)/p$ by $1$ in the off case and by $P(1/p)$ in the
on case. Their ratios and difference give~\eqref{eq:Presponse}.
\end{proof}

\begin{proof}[Proof of Theorem~\ref{thm:allmoments}]
Assign degree $d+k$ to $n^d\sigma_k(n)$ and degree $d+h_P$ to
$n^dF_P(n)$. These functions are positive integer-valued and
multiplicative, with the growth required by Theorem~\ref{thm:transfer}.
Their normalized phases are independent of $d$.
The divisor phase has the previously proved off-factor and response
\[
 b_{\sigma,k}(z)=1-z^{k+1},\qquad
 \delta_{\sigma,k}(z)=z^k.
\]
The polynomial Euler phases satisfy Lemma~\ref{lem:Pmeans}.

Their response functions are pairwise distinct and nonzero.
Indeed the transformation $H\mapsto H/(1+zH)$ has inverse
$\delta\mapsto\delta/(1-z\delta)$. For $\delta_P$ this inverse
is the nonzero polynomial $P-1$, so distinct $P$ give distinct
responses. For $\delta_{\sigma,k}=z^k$ the inverse is
$z^k/(1-z^{k+1})$, which is not a polynomial. Thus a divisor response
cannot equal a polynomial Euler response, and different divisor
indices give different monomials. At any fixed total degree, a
given $P$ or $k$ determines its weight $d$ uniquely, so no duplicate
response occurs in that degree. Apply Theorem~\ref{thm:transfer}.
\end{proof}

\subsection{Named arithmetic functions included in the theorem}
The statement is not limited to newly defined examples. The standard
Jordan and generalized Dedekind functions are
\[
 J_k(n)=n^k\prod_{p\mid n}(1-p^{-k}),\qquad
 \Psi_k(n)=n^k\prod_{p\mid n}(1+p^{-k}).
\]
Here $J_1=\varphi$ and $\Psi_1=\psi$, the ordinary Dedekind psi
function. They correspond to $P(X)=1-X^k$ and $1+X^k$.
More generally, for nonnegative integers $u_k,v_k$ of finite support,
not all zero,
\begin{equation}\label{eq:arithmeticproducts}
 P(X)=\prod_k(1-X^k)^{u_k}(1+X^k)^{v_k}
 \quad\Longrightarrow\quad
 F_P(n)=\prod_kJ_k(n)^{u_k}\Psi_k(n)^{v_k}.
\end{equation}
The required positivity holds at every $1/p$.
Different descriptions in~\eqref{eq:arithmeticproducts} can represent
the same $P$, so equality of polynomials must be checked first.
For instance $J_k(n)\Psi_k(n)=J_{2k}(n)$; these are one function,
not two independent coordinates.

The theorem therefore covers all polynomial moments of all positive powers
$\varphi(n)^m$ and $\psi(n)^m$, and of every distinct product in
\eqref{eq:arithmeticproducts}, jointly with all divisor-sum moments.
For example it proves the irrationality of
\begin{equation}\label{eq:newexample}
 \sum_{n\geq1}
 \frac{(n^2-3n+1)\sigma_2(n)+2\varphi(n)^3-5\psi(n)^2}{n!}.
\end{equation}
The three nonpolynomial arithmetic coordinates here have distinct
base functions, and their displayed coefficients do not all vanish.

Erd\H{o}s and Straus~\cite{ES74}, Theorem~3.7, already imply the
joint independence of $1,e,\sum\sigma(n)/n!$, and
$\sum\varphi(n)/n!$. Thus the degree-one specialization is not
claimed as new. The all-moment and full polynomial-Euler-factor
statements above are additional results of this manuscript, whose
historical priority remains subject to literature review.

\section{All derivative values and an exact rationality criterion}
\label{sec:derivatives}
For $k\geq1$ and $P\in\mathcal P$ define the entire functions
\[
 A_k(z)=\sum_{n\geq1}\frac{\sigma_k(n)}{n!}z^n,
 \qquad E_P(z)=\sum_{n\geq1}\frac{F_P(n)}{n!}z^n.
\]
Their entireness and termwise differentiability follow from the
polynomial growth bounds. Pratt~\cite{Pratt} explicitly discussed
the divisor-sum moment functions and their nonclosing derivative
hierarchy. The next corollary addresses rational linear relations
among all these derivative values at reciprocal integers. It does
not establish Pratt's broader expectation of transcendence at every
nonzero algebraic argument.

\begin{corollary}[Joint independence of all derivative values]
\label{cor:derivatives}
For each fixed $c\in\Z\setminus\{0\}$, the family
\[
 \{1,e^{1/c}\}\ \cup\
 \{A_k^{(j)}(1/c):k\geq1,\ j\geq0\}
 \ \cup\
 \{E_P^{(j)}(1/c):P\in\mathcal P,\ j\geq0\}
\]
is linearly independent over $\Q$.
\end{corollary}
\begin{proof}
For either numerator function $f$, termwise differentiation gives
\[
 \left(\sum_{n\geq1}\frac{f(n)z^n}{n!}\right)^{(j)}_{z=1/c}
 =c^j\sum_{n\geq1}\frac{n^{\underline j}f(n)}{c^n n!}.
\]
The falling factorial $X^{\underline j}$ is a monic polynomial of
degree $j$. The change from powers to falling factorials is triangular
over $\Q$ with diagonal entries $1$; the additional factors $c^j$
are nonzero rational numbers. Thus every finite proposed derivative
relation becomes a nontrivial relation in Theorem~\ref{thm:allmoments}.
\end{proof}

In particular $A_1'(1)/A_1(1)$ is irrational: the denominator is
positive, and a rational ratio would be a forbidden relation between
two derivative values. The inclusion of $j=0$ is essential; these
coordinates are the original factorial constants.

For the exact criterion below, write
\[
 T_d(x)=\sum_{j=0}^d\Stir{d}{j}x^j,\qquad
 \mathcal B_c(R)=\sum_dr_dT_d(1/c)
 \quad\text{if }R(X)=\sum_dr_dX^d.
\]
These are the Touchard polynomials; at $x=1$ they give Bell numbers.
For rational $R$ and integer $c\ne0$, $\mathcal B_c(R)$ is rational
and is computable by finite exact arithmetic.

\begin{corollary}[Complete rationality test for this class]
\label{cor:classification}
Let $R,U_k,V_P\in\Q[X]$, only finitely many nonzero, with equal
polynomial Euler factors grouped together. Set
\[
 S=\sum_{n\geq1}\frac{
 R(n)+\sum_kU_k(n)\sigma_k(n)+\sum_PV_P(n)F_P(n)}{c^n n!}.
\]
Then $S\in\Q$ if and only if every $U_k,V_P$ is the zero
polynomial and $\mathcal B_c(R)=0$. In that case $S=-R(0)$.
\end{corollary}
\begin{proof}
Expanding $n^d$ in falling factorials and summing gives
\[
 \sum_{n\geq1}\frac{R(n)}{c^n n!}
       =e^{1/c}\mathcal B_c(R)-R(0).
\]
Expand the remaining polynomial coefficients into powers of $n$,
and apply the joint independence in Theorem~\ref{thm:allmoments}.
\end{proof}

The qualification about the polynomial part is necessary:
$\sum_{n\geq1}(n-1)/n!=1$. Thus a nonzero polynomial numerator
alone need not produce an irrational value. By contrast, any
nonzero arithmetic polynomial coefficient in the canonical form
of Corollary~\ref{cor:classification} forces irrationality.

This is a complete classification inside the displayed class, not
a decision procedure for arbitrary convergent series. It makes no
claim about Euler's constant $\gamma$, Robin's inequality, or the
transcendence of an individual factorial constant.


\section{A concrete value at two and a denominator bound}
\label{sec:dyadic}
We now prove the concrete irrationality statement
\begin{equation}\label{eq:concrete-two}
 \boxed{C_*:=\sum_{n=1}^{\infty}\frac{2^n\sigma_5(n)}{n!}\notin\Q.}
\end{equation}
This is $A_5(2)$ in the notation of Section~\ref{sec:derivatives}.
It is not an instance of the earlier reciprocal-integer evaluation
theorem: $2$ cannot be written $1/c$ with $c$ a nonzero integer.
An additional denominator estimate, proved below, supplies the extension.

\paragraph{Historical scope.}
Pratt~\cite{Pratt} studies the functions $A_k(z)$ and discusses their
expected transcendence at nonzero algebraic arguments. The stated
geometric theorem in~\cite{TGgen} treats $z=1/c$ with integral $c\ne0$.
The sources examined did not supply an earlier proof of
\eqref{eq:concrete-two}. However, we did not locate a primary source
explicitly designating this particular value as a previously open,
named irrationality problem. Consequently we claim a concrete
irrationality theorem and an extension of the stated evaluation range,
not a certified resolution of a separately named historical conjecture.
In particular this is not a solution of the full algebraic-point
transcendence expectation, nor of Erd\H{o}s Problems 68, 249, or 1049.

Let $s_2(n)$ be the number of ones in the binary expansion of the
nonnegative integer $n$. We will use
\begin{equation}\label{eq:binary-height}
 v_2(n!)=n-s_2(n),\qquad 2^{s_2(n)}\leq n+1.
\end{equation}
For the first identity, sum $\lfloor n/2^j\rfloor$ and expand $n$ in
binary. For the second, the least nonnegative integer with $s$ binary
ones is $2^s-1$.

\begin{lemma}[A denominator bound for a normalized tail]
\label{lem:dyadic-height}
Fix $x=\varepsilon 2/b$, where $b\geq1$ is an integer and
$\varepsilon\in\{1,-1\}$. Suppose that integers $a_m$ of polynomial
growth satisfy $\sum_{m\geq0}a_mx^m/m!=u/q\in\Q$, with $q\geq1$.
Put
\[
 T_x(n)=n!x^{-n}\sum_{m>n}\frac{a_mx^m}{m!}.
\]
Then $q2^{s_2(n)}T_x(n)\in\Z$ for every $n\geq0$.
\end{lemma}
\begin{proof}
Subtract the finite partial sum from $u/q$ and multiply by
$n!x^{-n}$. The total term becomes integral after multiplication by
$q2^{s_2(n)}$, by~\eqref{eq:binary-height}. For $0\leq m\leq n$,
the scaled partial-sum coefficient is
\[
 \frac{n!}{m!}x^{m-n}
 =\varepsilon^{m-n}\frac{n!}{m!}\frac{b^{n-m}}{2^{n-m}}.
\]
There are no odd denominator factors, and its 2-adic valuation is
at least
\[
 v_2(n!)-v_2(m!)-(n-m)=s_2(m)-s_2(n)\geq-s_2(n).
\]
Multiplication by $2^{s_2(n)}$ therefore clears every partial term.
The proof also permits even $b$; those cases include the old
reciprocal-integer arguments.
\end{proof}

\begin{lemma}\label{lem:dyadic-exp}
For every $x=\varepsilon 2/b$ as above, $e^x$ is irrational.
\end{lemma}
\begin{proof}
If $e^x=u/q$ were rational, apply Lemma~\ref{lem:dyadic-height} to
$a_m=1$ and write the resulting tail as $U_x(n)$. Its first terms give
\[
 U_x(n)=\frac{x}{n}+\frac{x^2-x}{n^2}+O(n^{-3}),
 \qquad U_x(n)-U_x(n+1)=\frac{x}{n^2}+O(n^{-3}).
\]
The difference is nonzero eventually. On the other hand, it lies
in the lattice with denominator
$q2^{\max(s_2(n),s_2(n+1))}\leq q(n+2)$.
Multiplication by that denominator yields a nonzero integer tending
to zero, a contradiction. The tail expansion follows by keeping the
first two factorial terms and bounding the remaining terms by
$O(n^{-3})$; it is valid for either sign of $x$.
\end{proof}

\begin{theorem}[Evaluation at rational points of numerator at most two]
\label{thm:dyadic}
Let $\mathcal F_*$ denote the numerator functions
\[
 \{n\mapsto n^d\sigma_k(n):d\geq0,k\geq1\}
 \ \cup\
 \{n\mapsto n^dF_P(n):d\geq0,P\in\mathcal P\}.
\]
For every fixed $x=\varepsilon 2/b$, $b\geq1$ and
$\varepsilon\in\{1,-1\}$, the family
\[
 \{1,e^x\}\ \cup\
 \left\{\sum_{n\geq1}\frac{f(n)x^n}{n!}:f\in\mathcal F_*\right\}
\]
is linearly independent over $\Q$.
Equivalently, this covers every nonzero rational argument whose
reduced numerator has absolute value one or two, at one fixed
argument at a time.
\end{theorem}
\begin{proof}
Consider a nontrivial integral-coefficient relation
$\lambda e^x+\sum_f a_f\sum_{m\geq1}f(m)x^m/m!\in\Q$.
If all $a_f=0$, use Lemma~\ref{lem:dyadic-exp}. Otherwise retain
the finitely many functions with $a_f\ne0$, and let $K\geq1$ be
their largest growth degree. Write the rational total as $u/q$.
The combined coefficients, including the $m=0$ exponential term,
are integers. Thus their combined normalized tail $T_x(n)$ obeys
Lemma~\ref{lem:dyadic-height}.

Set $H=K+1$. In addition to the retained arithmetic functions,
include $g_0(n)=1$ in the fiber constraints. Choose a grid of dimension
$H$ and side length $L>(s+1)H(H+1)/2$, where $s$ is the number
of retained arithmetic functions. Lemmas~\ref{lem:grid} and
\ref{lem:kernel} supply nonzero integral weights satisfying
\[
 \sum_{e_j}w_e\frac{p_e^{i+1}}{g(p_e)}=0
 \quad (0\leq i<H,\ 0\leq j\leq i)
\]
on each coordinate-$j$ fiber, for every retained $g=f$ and for $g_0$.
The squared CRT alignment and the shifts are exactly as in
Theorem~\ref{thm:transfer}.

For every such $g$, the centered expansion through order $H$ is
\begin{equation}\label{eq:dyadic-expansion}
 T_{g,x}(n)=\sum_{j=0}^H\sum_{i=j}^H
 x^{j+1}\Stir{i}{j}\frac{g(n+j+1)}{(n+j+1)^{i+1}}
       +O(n^{-5/2}).
\end{equation}
The error follows from $g(n)=O(n^{K+1/2})$ and
$K+1/2-H-2=-5/2$. The same geometric tail estimate applies for
fixed $|x|\leq2$.

Along the aligned progression, put $n_e=(N-t_e)/p_e$ and
$W(N)=\sum_ew_eT_x(n_e)$. All $i<H$ terms vanish, including those
of the exponential. Each remaining arithmetic term has size
$O(N^{K+1/2-H-1})=O(N^{-3/2})$, so
\[
 W(N)=O(N^{-3/2}).
\]
Unlike the earlier proof, $W(N)$ need not itself be an integer.
Its common denominator divides
\[
 D_N=q2^{\max_e s_2(n_e)}
       \leq q(1+\max_e n_e)=O(N).
\]
All defects are powers of the same prime, so a maximum suffices;
one does not multiply the individual denominators. Therefore
$D_NW(N)\in\Z$ and $D_NW(N)=O(N^{-1/2})\to0$.
It follows that $W(N)=0$ eventually.

Multiply the remaining identity by $N^2$. The errors are
$O(N^{-1/2})$, all degrees below $K$ disappear, and the top degrees
give
\begin{equation}\label{eq:dyadic-limit}
 \sum_{h_f=K}\sum_e\sum_{j=0}^H
 \frac{a_fw_ep_e^{H+1}x^{j+1}\Stir{H}{j}}{f(p_e)}
      \rho_f(N+s_{e,j})\longrightarrow0.
\end{equation}
Indeed $H=K+1$, so the retained normalized factor before taking
the limit is $N^2/(N+s)^2$; replacing it by $1$ costs
$O(N^{-1/2})$. The $g_0$ contribution is $O(N^{-K})$ and vanishes.
Choose $e_*$ with $w_{e_*}\ne0$. Its terminal shift is unique,
and a top-degree coefficient there is
$a_fw_{e_*}p_{e_*}^{H+1}x^{H+1}/f(p_{e_*})\ne0$.
The distinct response functions proved in
Theorem~\ref{thm:allmoments} and Lemma~\ref{lem:multimean}
contradict~\eqref{eq:dyadic-limit}.
\end{proof}

Taking $x=2$, $d=0$, and $k=5$ proves~\eqref{eq:concrete-two}.

For identification, exact rational interval arithmetic gives
\[
 C_*=3412.81774524922338054544082302233096428831099819146187\ldots.
\]
The first $50$ decimal places are certified by the bounds
\[
 S_{80}<C_*<S_{80}+\frac{4\cdot2^{81}81^5}{81!},
 \qquad S_{80}=\sum_{n=1}^{80}\frac{2^n\sigma_5(n)}{n!}.
\]
Indeed $\sigma_5(n)<2n^5$, and the consecutive-term ratio of the
majorant $2^{n+1}n^5/n!$ is $2(n+1)^4/n^5$, decreasing and below
$1/2$ for $n\geq81$. This is an identification certificate, not the
irrationality proof.

The derivative and rationality-classification corollaries also extend
to each permitted $x$: replace $c^j$ by $x^{-j}$ in the derivative
conversion and evaluate the Touchard polynomials at $x$.

The new argument has a precise limitation. If the reduced numerator
of $x$ contains a prime power $p^a$ with $a>1/(p-1)$, then the
corresponding defect in $n!x^{-n}$ has valuation
\[
 v_p(n!)-an=-\left(a-\frac1{p-1}\right)n+O(\log n).
\]
Its denominator grows exponentially, so the fixed extra cancellation
order used here does not clear it. No extension to numerator three,
to all rational arguments, or to all algebraic arguments is asserted.


\section{Prime-power profiles: products, radicals, and counting functions}
\label{sec:profiles}

Divisibility by a fresh prime observes only one average of the
prime-power values of a multiplicative function.  A refinement that
specifies the exact valuation observes an individual prime-power
value.  This extra freedom gives a substantially larger class than
polynomial Euler factors or single divisor sums.

\subsection{An admissible profile and its progression means}

Write $R\in\Q(z)[w]$ for a polynomial in $w$ whose coefficients are
rational functions of $z$.  We call $(R,h)$ admissible if $h\geq1$
is an integer, the coefficients of $R$ have no poles at $z=1/p$
for primes $p$, and the following conditions hold:
\begin{enumerate}
\item There is a constant $C_R$ such that
\[
 0<R(1/p,p^{-a})\leq C_R\qquad(p\text{ prime},\ a\geq1).
\]
\item $R(z,z)=1+O(z)$ as $z\to0$ through positive real values.
\item $p^{ha}R(1/p,p^{-a})$ is an integer for every prime $p$
and every $a\geq1$.
\end{enumerate}
Define the positive integer-valued multiplicative function $f_{R,h}$ by
\begin{equation}\label{eq:profile-f}
 f_{R,h}(1)=1,\qquad f_{R,h}(p^a)=p^{ha}R(1/p,p^{-a}),
 \qquad \rho_R(n)=\frac{f_{R,h}(n)}{n^h}.
\end{equation}
The normalized function $\rho_R$ is independent of $h$.  For each
$\epsilon>0$, boundedness of the local factors gives
$\rho_R(n)=O_{R,\epsilon}(n^\epsilon)$: separate the finitely many
primes for which $C_R>p^\epsilon$.  Thus the factorial series below
converge absolutely and define entire functions.

\begin{lemma}[Means from a bounded prime-power profile]
\label{lem:profile-means}
For every admissible profile, the normalized function $\rho_R$ has
a finite, strictly positive mean $M_R(Q,a)$ on every arithmetic
progression.  Put $\rho_R(p^0)=1$, and write
$R(z,w)=\sum_{j=0}^{d}r_j(z)w^j$.  The unrestricted local mean is
the rational function
\begin{equation}\label{eq:profile-U}
 U_R(z)=(1-z)\left(1+
           \sum_{j=0}^{d}r_j(z)\frac{z^{j+1}}{1-z^{j+1}}\right).
\end{equation}
At a fresh prime $\ell\nmid Q$, the ratios of refined to original
means are $1/U_R(1/\ell)$ when the new residue is a unit modulo
$\ell$, and $R(1/\ell,\ell^{-t})/U_R(1/\ell)$ when its valuation
is exactly $t\geq1$.
\end{lemma}
\begin{proof}
Let $g=\rho_R*\mu$, where $*$ denotes Dirichlet convolution.
Then $g$ is multiplicative and
\[
 g(p)=O_R(p^{-1}),\qquad
 g(p^a)=\rho_R(p^a)-\rho_R(p^{a-1})=O_R(1)\quad(a\geq2).
\]
Consequently
\[
 \sum_{a\geq1}\frac{|g(p^a)|}{p^a}=O_R(p^{-2}),\qquad
 \sum_{d\geq1}\frac{|g(d)|}{d}<\infty.
\]
The second assertion follows by multiplying the absolutely convergent
nonnegative Euler factors; the finitely many small primes also have
finite factors.  Since $\rho_R(n)=\sum_{d\mid n}g(d)$, counting
divisibility along a progression gives
\begin{equation}\label{eq:profile-mean-series}
 M_R(Q,a)=\sum_{\gcd(d,Q)\mid a}g(d)\frac{\gcd(d,Q)}{d}.
\end{equation}
Here is a justification that requires no pointwise bound on $g(d)$.
The contribution beyond a fixed cutoff $Y$ to a progression average
of length $X$ is bounded by
\[
 Q\sum_{d>Y}\frac{|g(d)|}{d}
       +\frac1X\sum_{d\leq QX+|a|}|g(d)|.
\]
The second term tends to zero.  Indeed split its sum at $Y$; the
fixed initial part is $o(X)$, while the rest divided by $X$ is at
most $(Q+o(1))\sum_{d>Y}|g(d)|/d$.  Let first $X$, then $Y$, tend
to infinity.  This also justifies the absolutely convergent series
in~\eqref{eq:profile-mean-series}.

Its Euler factors are particularly transparent as local conditional
expectations.  At $p\nmid Q$ the factor is
\[
 (1-p^{-1})\sum_{v\geq0}\rho_R(p^v)p^{-v}=U_R(1/p).
\]
At $p^s\Vert Q$, if $v_p(a)=t<s$, the factor is
$\rho_R(p^t)$.  If $p^s\mid a$, it is
$(1-p^{-1})\sum_{v\geq s}\rho_R(p^v)p^{-(v-s)}$.
Every factor is positive.  The unrestricted factors are
$1+O_R(p^{-2})$, so their product is strictly positive.
Expanding $R$ in powers of $w$ and summing geometric series gives
\eqref{eq:profile-U}.  In particular $U_R(z)=1+O_R(z^2)$ near
$z=0$: this holds along $z=1/p$, and the rational function
$U_R(z)-1$ therefore has order at least two at zero.
Fixing the exact valuation at a fresh prime
replaces its unrestricted factor by the corresponding single value
$\rho_R(\ell^t)$; making it a unit replaces that factor by $1$.
\end{proof}

\subsection{A vector of valuation responses separates every profile}

For $t\geq1$ define the rational function
\begin{equation}\label{eq:profile-response}
 \delta_{R,t}(z)=\frac{R(z,z^t)-1}{U_R(z)}.
\end{equation}
It is the relative difference between an exact-valuation-$t$ mean
and a unit mean.  Individual responses may vanish or coincide.
The entire vector of responses has neither defect once identical
profiles have been grouped.

\begin{lemma}[Injectivity of valuation responses]
\label{lem:profile-injective}
If $\delta_{R,t}=\delta_{S,t}$ as rational functions for every
$t\geq1$, then $R=S$ in $\Q(z)[w]$.  The response vector of $R$
vanishes identically if and only if $R=1$.
\end{lemma}
\begin{proof}
For sufficiently small real $0<z<1$, the geometric series converge
and the denominators are nonzero.  From~\eqref{eq:profile-U},
\[
 \sum_{t\geq1}\delta_{R,t}(z)z^t
       =\frac{U_R(z)-1}{(1-z)U_R(z)}.
\]
Equality of response vectors therefore implies $U_R=U_S$ and
then $R(z,z^t)=S(z,z^t)$ for every $t$.  For each such fixed $z$,
the polynomial $R(z,w)-S(z,w)$ has infinitely many distinct roots
$w=z^t$, hence is zero.  Its rational coefficient functions vanish
on an interval, so they are identically zero.  The zero-vector
assertion is the special case $S=1$, for which $U_S=1$.
\end{proof}

\begin{lemma}[Separation by fresh prime powers]
\label{lem:profile-separation}
Let $R_1,\ldots,R_s$ be pairwise distinct admissible profiles,
none equal to $1$.  After identical shifts have been grouped, any
finite sum of their shifted normalized functions with at least one
nonzero coefficient cannot tend to a finite constant on a progression.
\end{lemma}
\begin{proof}
Group repeated shifts and choose a shift $r_*$ with a nonzero
coefficient vector.  Suppose
the sum has limit $L$ on $N\equiv A\pmod Q$.  Write its coefficients
at $r_*$ as $\beta_j$, and put
$B_j=\beta_jM_{R_j}(Q,A+r_*)$.
Choose all fresh primes larger than the absolute differences between
distinct shifts and larger than the number of shifts.  For a chosen
prime $\ell$ and depth $t\geq1$, the on state is
\[
 N+r_*\equiv\ell^t\pmod{\ell^{t+1}}.
\]
It has exactly the desired valuation and makes every other shift a
unit modulo $\ell$.  Choose an off residue modulo $\ell$ missing
all shifts, and lift it to a residue modulo $\ell^{t+1}$.  It makes
all shifts units.  The Chinese remainder theorem allows these two
states independently at distinct primes while preserving the base
progression.

Take progression means on the resulting refinements.  For any
nonempty finite set $S$ of selected primes, the alternating sum over
its $2^{|S|}$ states cancels the constant limit and every shift other
than $r_*$.  Lemma~\ref{lem:profile-means} gives
\begin{equation}\label{eq:profile-mixed}
 \sum_{j=1}^s B_j\prod_{\ell\in S}
          \delta_{R_j,t_\ell}(1/\ell)=0.
\end{equation}
For each subset only its own primes refine the original progression;
thus the same $B_j$ occur in all these identities.

Fix a target index $j_0$.  By Lemma~\ref{lem:profile-injective},
choose a depth $t_0$ for which $\delta_{R_{j_0},t_0}$ is nonzero.
For each $j\ne j_0$, choose a depth $t_j$ for which
$\delta_{R_{j_0},t_j}-\delta_{R_j,t_j}$ is nonzero.  Choose distinct
fresh primes $\ell_0,\ell_j$ avoiding the finitely many zeros and
poles of these rational functions.  A linear combination of
\eqref{eq:profile-mixed}, over subsets containing $\ell_0$, yields
\[
 \sum_v B_v\delta_{R_v,t_0}(1/\ell_0)
  \prod_{j\ne j_0}
   \bigl(\delta_{R_v,t_j}(1/\ell_j)
              -\delta_{R_j,t_j}(1/\ell_j)\bigr)=0.
\]
Every term except $v=j_0$ vanishes, and the remaining multiplier is
nonzero.  Thus $B_{j_0}=0$.  All progression means are positive, so
all coefficients at the distinguished shift vanish, a contradiction.
\end{proof}

\subsection{The larger irrationality theorem}

\begin{theorem}[All moments of distinct prime-power profiles]
\label{thm:profile-transfer}
Let $\mathcal R$ be any collection of pairwise distinct profiles
$R\ne1$, and choose for each $R$ an integer $h_R\geq1$ such that
$(R,h_R)$ is admissible.  For each fixed
$x=\varepsilon2/b$, where $b\geq1$ and $\varepsilon\in\{1,-1\}$,
the family
\begin{equation}\label{eq:profile-independent}
 \{1,e^x\}\ \cup\
 \left\{\sum_{n\geq1}\frac{n^d f_{R,h_R}(n)x^n}{n!}:
                    R\in\mathcal R,\ d\geq0\right\}
\end{equation}
is linearly independent over $\Q$.  The assertion is at one fixed
evaluation point at a time.
\end{theorem}
\begin{proof}
Only finitely many coordinates occur in a proposed relation.
Each numerator $f(n)=n^df_{R,h_R}(n)$ is positive integer-valued and
multiplicative, has degree $h_R+d\geq1$, and satisfies
$f(n)=O(n^{h_R+d+1/2})$.  Its normalized phase is precisely $\rho_R$.
These are all the arithmetic and growth inputs used in the proof
of Theorem~\ref{thm:dyadic} before its final separation step.

For completeness, let $K$ be the largest retained degree and use
the same grid with $H=K+1$, including $1$ in the cancellation
constraints.  The normalized-tail expansion has error $O(n^{-5/2})$.
All lower fiber terms cancel, giving $W(N)=O(N^{-3/2})$.  Its
denominator divides $q2^{\max_e s_2(n_e)}=O(N)$, so $W(N)=0$
eventually.  Multiplication by $N^2$ gives exactly the limiting
identity~\eqref{eq:dyadic-limit}, now with phases $\rho_R$.
At degree $K$, each profile occurs at most once: its weight is
uniquely $d=K-h_R$.  The terminal shift has a nonzero coefficient
because the grid weight and $f(p_e)$ are nonzero.  Apply
Lemma~\ref{lem:profile-separation} to contradict that limiting
identity.  If no arithmetic coordinate was present, use
Lemma~\ref{lem:dyadic-exp}.
\end{proof}

The existing derivative and Touchard-polynomial arguments now apply
to this entire class.  In particular, if
$F_R(z)=\sum_{n\geq1}f_{R,h_R}(n)z^n/n!$, all numbers
$F_R^{(j)}(x)$, together with $1,e^x$, are jointly independent.
For a finite expression
\[
 S=\sum_{n\geq1}\left[V(n)+\sum_R V_R(n)f_{R,h_R}(n)\right]
                      \frac{x^n}{n!},\qquad V,V_R\in\Q[X],
\]
the exact rationality criterion is: every $V_R$ vanishes, and
$\sum_d v_dT_d(x)=0$ for $V(X)=\sum_dv_dX^d$.  Then $S=-V(0)$.
Profiles must first be grouped by equality in $\Q(z)[w]$; this is
an exact finite algebraic test after clearing the $z$-denominators.
If the same profile is written with different base degrees, first
factor out the corresponding power of $n$.

\subsection{Named families covered by the theorem}

The following table gives the growth degree and the normalized
profile.  All parameters are positive integers unless otherwise
specified.  The notation $p^a\Vert n$ specifies the exact prime
power dividing $n$.
\begin{center}
\small
\begin{tabular}{p{0.37\textwidth}c p{0.44\textwidth}}
\toprule
Numerator function & $h$ & $R(z,w)$\\
\midrule
$\sigma_k(n)$ & $k$ & $(1-z^kw^k)/(1-z^k)$\\[5pt]
$J_k(n)$ & $k$ & $1-z^k$\\[5pt]
$\Psi_k(n)$ & $k$ & $1+z^k$\\[5pt]
$\sigma_k^*(n)=\prod_{p^a\Vert n}(1+p^{ak})$ & $k$ & $1+w^k$\\[5pt]
$J_k^*(n)=\prod_{p^a\Vert n}(p^{ak}-1)$ & $k$ & $1-w^k$\\[5pt]
$\operatorname{rad}(n)=\prod_{p\mid n}p$ & $1$ & $w/z$\\[5pt]
$a_r(n)$, index-$n$ subgroups of $\Z^r$, $r\geq2$
 & $r-1$ & $\displaystyle\prod_{j=1}^{r-1}\frac{1-z^jw}{1-z^j}$\\[9pt]
$g(n)=\sum_{d\mid n}d\varphi(d)$ & $2$ & $(1+zw^2)/(1+z)$\\
\bottomrule
\end{tabular}
\end{center}

The original polynomial Euler factors are the profiles $R(z,w)=P(z)$.
Every finite pointwise product of entries in the table and of such
factors is again admissible: multiply the profiles and add the
degrees.  This includes arbitrary powers and mixed products of
divisor sums, ordinary and unitary Jordan or Dedekind factors,
radicals, and subgroup counts.  Positivity and integrality follow
from their definitions.  Each displayed profile is uniformly
bounded at $(z,w)=(1/p,p^{-a})$ and equals $1+O(z)$ when $w=z$.
For example, $w/z=p^{1-a}\leq1$, so the pole of this coefficient
at $z=0$ causes no difficulty.  All assertions survive multiplication
of any fixed finite number of factors.

These are products of arithmetic functions \emph{inside} a factorial
series.  The theorem does not assert algebraic independence of the
resulting constants or irrationality of their arbitrary products.
Nor may algebraically identical profiles be counted as separate
coordinates.  For example,
\[
 J_k\Psi_k=J_{2k},\qquad J_k^*\sigma_k^*=J_{2k}^*,
 \qquad a_2=\sigma_1.
\]

Here is a useful nonredundant subfamily.  Let $P$ range over the
polynomials of Section~8, now permitting $P=1$, and let
$\mathbf m=(m_1,m_2,\ldots)$ be a finitely supported sequence of
nonnegative integers.  Exclude $(P,\mathbf m)=(1,\mathbf0)$ and put
\[
 f_{P,\mathbf m}(n)=F_P(n)\prod_{k\geq1}\sigma_k(n)^{m_k},
 \quad F_1(n)=1,\quad
 h=\deg P+\sum_k km_k.
\]
Its profile is
\begin{equation}\label{eq:sigma-product-profile}
 R_{P,\mathbf m}(z,w)=P(z)
       \prod_k\left(\frac{1-z^kw^k}{1-z^k}\right)^{m_k}.
\end{equation}
All these profiles are distinct.  Indeed equality of two such
profiles, first at $w=0$ and then at arbitrary $w$, would give
$\prod_k(1-(zw)^k)^{m_k-n_k}=1$.  The largest index with a
nonzero exponent has a cyclotomic factor not supplied by smaller
indices, a contradiction.  Thus $\mathbf m=\mathbf n$, and then
$P=Q$.  Consequently this whole family, with every polynomial
moment, is covered simultaneously without further duplicate tests.

\subsection{Lattice counts, least common multiples, and explicit constants}

For clarity the classical subgroup formula~\cite{Zou06} in the table can be obtained
without using its interpretation as a black box.  Hermite normal
form gives
\[
 a_r=1*\operatorname{id}*\operatorname{id}^2*\cdots*
                \operatorname{id}^{r-1},\qquad
 \sum_{a\geq0}a_r(p^a)T^a=
                \prod_{j=0}^{r-1}(1-p^jT)^{-1}.
\]
The coefficient is the Gaussian binomial
\[
 a_r(p^a)=p^{(r-1)a}
           \prod_{j=1}^{r-1}\frac{1-p^{-(a+j)}}{1-p^{-j}},
\]
which proves the displayed profile and its integrality.  Distinct
$r\geq2$ give distinct profiles by their different degrees in $w$.
In particular
\[
 \sum_{n\geq1}\frac{a_3(n)2^n}{n!}\notin\Q,
\]
and all its derivative values join the same independent family.
This is a factorial-weighted count of finite-index sublattices of
the ordinary integer lattice $\Z^3$.

Define the standard least-common-multiple sum~\cite{IM14}
$L(n)=\sum_{j=1}^{n}\operatorname{lcm}(j,n)$.  Grouping terms by
$\gcd(j,n)$, and using the pairing of reduced residues modulo $m$,
gives the exact identity
\begin{equation}\label{eq:lcm-sum-identity}
 L(n)=\frac n2\left(1+\sum_{d\mid n}d\varphi(d)\right)
          =\frac n2(1+g(n)).
\end{equation}
For $m>1$ the reduced residues sum to $m\varphi(m)/2$; the term
$m=1$ contributes $1$, which accounts for the extra $1$ above.
The multiplicative function $g$ has
\[
 g(p^a)=1+\sum_{j=1}^a p^j(p^j-p^{j-1})
        =\frac{p^{2a+1}+1}{p+1},
\]
so its profile is $(1+zw^2)/(1+z)$, as claimed.  Therefore
\begin{equation}\label{eq:lcm-constant}
 \sum_{n\geq1}\frac{x^n}{n!}
                  \sum_{j=1}^{n}\operatorname{lcm}(j,n)
       =\frac{x e^x}{2}+
          \frac12\sum_{n\geq1}\frac{n g(n)x^n}{n!}\notin\Q
\end{equation}
at every permitted $x$.  Its mean connection with zeta values is
also exact:
\[
 U_g(z)=\frac{1-z^2}{1-z^3},\qquad
 \lim_{X\to\infty}\frac1X\sum_{n\leq X}\frac{g(n)}{n^2}
        =\frac{\zeta(3)}{\zeta(2)}.
\]
The appearance of this ratio in the mean does not establish its
irrationality; the theorem concerns the factorial series.

For additional concrete instances, each of
\[
 \sum_{n\geq1}\frac{2^n\operatorname{rad}(n)}{n!},\qquad
 \sum_{n\geq1}\frac{2^n\sigma_1(n)^2}{n!},\qquad
 \sum_{n\geq1}\frac{2^n\varphi(n)^3\sigma_5(n)^2}{n!}
\]
is irrational, and their polynomially weighted versions are
covered by the same theorem.  The corresponding series with
unitary totients $J_1^*$ or unitary divisor sums $\sigma_1^*$ are
also covered.  These examples illustrate the breadth of the
proved implication; their individual historical open status and
priority have not been established by a literature search.


\section{Every residue class, signed twists, and four real transforms}
\label{sec:periodic}
A periodic restriction is not absorbed by changing the evaluation point
in the preceding theorem. The grid can nevertheless preserve an entire
residue class, giving independence of the separate pieces of every series.
For an integer $q\geq1$, a residue $0\leq r<q$, and $x\ne0$, put
\[
 E_{q,r}(x)=\sum_{\substack{n\geq0\\n\equiv r\ (q)}}\frac{x^n}{n!},
 \qquad
 B_{f,q,r}(x)=\sum_{\substack{n\geq1\\n\equiv r\ (q)}}
                       \frac{f(n)x^n}{n!}.
\]
The $n=0$ convention is intentional: exponential terms include it,
whereas arithmetic functions are defined on positive integers.

\begin{lemma}[Periodic exponential independence]\label{lem:periodic-exp}
For every fixed $q\geq1$ and $x=\varepsilon2/b$, $b\geq1$,
$\varepsilon\in\{1,-1\}$, the numbers
$1,E_{q,0}(x),\ldots,E_{q,q-1}(x)$ are linearly independent over $\Q$.
\end{lemma}
\begin{proof}
A nontrivial relation would give a nonzero integer-valued $q$-periodic
sequence $a(n)$ such that $\sum_{n\geq0}a(n)x^n/n!=u/v\in\Q$.
Define $U(n)=n!x^{-n}\sum_{m>n}a(m)x^m/m!$.
Lemma~\ref{lem:dyadic-height} gives
$v2^{s_2(n)}U(n)\in\Z$. Since $a$ is bounded, expansion of its first
two tail terms, followed by a geometric bound for the rest, gives
\[
 U(n)=\frac{x a(n+1)}n+
 \frac{x^2a(n+2)-xa(n+1)}{n^2}+O(n^{-3}).
\]
Choose a residue of $n$ on which $a(n+1)\ne0$.
Periodicity makes both displayed coefficients identical for $n$ and
$n+q$, whence on this infinite progression
\[
 U(n)-U(n+q)=\frac{qx a(n+1)}{n^2}+O(n^{-3}).
\]
This difference is eventually nonzero. Multiplying it by
$v2^{\max(s_2(n),s_2(n+q))}\leq v(n+q+1)$ gives a nonzero integer
tending to zero. The contradiction proves the claim. The term $a(0)$
is already included in the denominator lemma and creates no exception.
\end{proof}

\begin{theorem}[Periodic transfer principle]\label{thm:periodic-transfer}
Let $\mathcal F$ consist of positive integer-valued multiplicative
functions with assigned integers $h_f\geq0$ and
$f(n)=O_f(n^{h_f+1/2})$. Suppose that, within each fixed degree, the
normalized phases $\rho_f(n)=f(n)/n^{h_f}$ have the isolated-shift
separation property: no finite shifted linear combination with a
nonzero grouped coefficient at some shift has a finite limit on an
arithmetic progression. This is the conclusion of
Lemma~\ref{lem:multimean}; a proved stronger local-response separation
criterion may be used in its place.
Then, for every fixed $q\geq1$ and $x=\varepsilon2/b$ as above,
\begin{equation}\label{eq:periodic-family}
 \{1\}\ \cup\ \{E_{q,r}(x):0\leq r<q\}
 \ \cup\ \{B_{f,q,r}(x):f\in\mathcal F,\ 0\leq r<q\}
\end{equation}
is linearly independent over $\Q$.
In particular this applies to the arithmetic family of
Theorem~\ref{thm:dyadic}, including every polynomial moment.
The modulus and the positive magnitude of the evaluation point are
fixed; different moduli are not listed as independent coordinates.
\end{theorem}
\begin{proof}
Clear denominators in a proposed finite relation. It says that
\[
 \sum_{m\geq0}a_0(m)\frac{x^m}{m!}
 +\sum_f\sum_{m\geq1}a_f(m)f(m)\frac{x^m}{m!}\in\Q,
\]
where every $a_f$ and $a_0$ is an integer-valued $q$-periodic sequence.
If all arithmetic sequences vanish, apply
Lemma~\ref{lem:periodic-exp}. Otherwise discard the zero sequences,
put $K=\max h_f$, and set $H=K+1$. Include $g_0(n)=1$ in the
cancellation constraints, just as in Theorem~\ref{thm:dyadic}.

Use the grid construction with this $H$, a sufficiently large $L$,
and the same $F$, but replace $D=F!$ by $D=qF!$ and retain
$B=1+HDF$. Every proof in Lemma~\ref{lem:grid} is unchanged:
$F!\mid D$, and distinct multipliers remain pairwise coprime.
Additionally,
\[
 p_e\equiv1\pmod q,\qquad t_e\equiv0\pmod q.
\]
The nonzero integer weights obey
\[
 \sum_{e_j}w_e\frac{p_e^{i+1}}{g(p_e)}=0
 \quad(0\leq j\leq i<H)
\]
on every coordinate-$j$ fiber, for every retained arithmetic function
and for $g_0$. Pick a top-degree function $f_*$ and a residue $A$
such that $a_{f_*}(A+H+1)\ne0$.
Since $\gcd(q,p_e)=1$, CRT simultaneously imposes
\[
 N\equiv t_e\pmod{p_e^2}\quad\hbox{for every }e,
 \qquad N\equiv A\pmod q.
\]
Writing $n_e=(N-t_e)/p_e$, we have $n_e\equiv A\pmod q$.
Consequently the periodic multiplier in every centered tail term is
\[
 a_f(n_e+j+1)=a_f(A+j+1),
\]
independent of the vertex. Thus every lower fiber cancellation remains
exact despite the periodic coefficients. Multiplicativity still gives
$f(n_e+j+1)=f(N+s_{e,j})/f(p_e)$.

The common expansion through order $H$ has remainder $O(n^{-5/2})$,
since the periodic coefficients are bounded. After cancellation,
$W(N)=\sum_e w_e T(n_e)=O(N^{-3/2})$. Its denominator divides
$v2^{\max_e s_2(n_e)}=O(N)$, where $v$ is the denominator of the
rational total. Hence $W(N)=0$ eventually. Multiplication by $N^2$
removes all degrees below $K$ and all errors, giving a finite limit for
\begin{equation}\label{eq:periodic-limit}
 \sum_{h_f=K}\sum_e\sum_{j=0}^H
 \frac{w_ep_e^{H+1}x^{j+1}\Stir{H}{j}}{f(p_e)}
     a_f(A+j+1)\rho_f(N+s_{e,j}).
\end{equation}
For $K\geq1$ the limit is zero. For $K=0$, the auxiliary $g_0$ terms
converge to a constant instead; this is harmless because the separation
hypothesis excludes every finite limit. As before, replacing
$N^2/(N+s)^2$ by $1$ costs $O(N^{-1/2})$.

Choose $e_*$ with $w_{e_*}\ne0$. The terminal shift $s_{e_*,H}$
is unique. Its $f_*$ coefficient in~\eqref{eq:periodic-limit} is
\[
 \frac{w_{e_*}p_{e_*}^{H+1}x^{H+1}}{f_*(p_{e_*})}
                 a_{f_*}(A+H+1)\ne0.
\]
This contradicts the assumed isolated-shift separation and proves the
theorem, including the possible degree-zero case.
\end{proof}

\subsection{Simultaneous positive, negative, and oscillatory values}
Write $B_f(z)=\sum_{n\geq1}f(n)z^n/n!$.
For a positive permitted $x=2/b$, the modulus-two theorem and the
invertible rational change of coordinates
\[
 B_f(x)=B_{f,2,0}(x)+B_{f,2,1}(x),\qquad
 B_f(-x)=B_{f,2,0}(x)-B_{f,2,1}(x)
\]
prove joint independence of
\[
 \{1,e^x,e^{-x}\}\ \cup\ \{B_f(x),B_f(-x):f\in\mathcal F\}.
\]
There is an equally explicit four-direction statement.

\begin{corollary}[Four real transforms at once]\label{cor:four-directions}
For every positive $x=2/b$, the real numbers in
\[
 \{1,e^x,e^{-x},\cos x,\sin x\}
 \ \cup\
 \{B_f(x),B_f(-x),\operatorname{Re}B_f(ix),
                         \operatorname{Im}B_f(ix):f\in\mathcal F\}
\]
are jointly linearly independent over $\Q$.
\end{corollary}
\begin{proof}
Apply Theorem~\ref{thm:periodic-transfer} with $q=4$. Relative to the
four residue coordinates, the displayed transforms have matrix
\[
 \begin{pmatrix}
 1&1&1&1\\
 1&-1&1&-1\\
 1&0&-1&0\\
 0&1&0&-1
 \end{pmatrix},
\]
whose determinant is $-8$. The same matrix takes the four exponential
residue sums to $e^x,e^{-x},\cos x,\sin x$. Apply this invertible
rational transformation separately to every function block.
\end{proof}
For example, with $f(n)=\sigma_5(n)$, this proves that both oscillatory
series
\[
 \sum_{m\geq1}\frac{(-1)^m\sigma_5(2m)\,2^{2m}}{(2m)!},
 \qquad
 \sum_{m\geq0}\frac{(-1)^m\sigma_5(2m+1)\,2^{2m+1}}{(2m+1)!}
\]
are irrational, jointly with $B_f(2)$ and $B_f(-2)$ and the listed
classical exponential and trigonometric values. These are statements
about real rational linear relations. They do not assert linear
independence over algebraic numbers, algebraic independence, or
transcendence of the arithmetic values.

\subsection{Periodic signs and conditional Poisson expectations}
Any nonzero rational-valued periodic weight $w(n)$ now gives
\[
 \sum_{n\geq1}\frac{w(n)f(n)x^n}{n!}\notin\Q.
\]
This includes real Dirichlet characters and the integer Ramanujan
weights $c_d(n)=\sum_{1\leq a\leq d,\ (a,d)=1}e^{2\pi ian/d}$.
More generally, a finite collection of periodic weights which is
linearly independent over $\Q$ as functions on $\Z/q\Z$ gives
independent weighted values within each arithmetic block.

A probabilistic interpretation supplies further explicit constants.
Let $N$ have the Poisson distribution of positive mean $x=2/b$, and
extend $f$ by $f(0)=0$. Then
\begin{equation}\label{eq:conditional-poisson}
 \mathbb E\!\left[f(N)\mid N\equiv r\pmod q\right]
       =\frac{B_{f,q,r}(x)}{E_{q,r}(x)}\notin\Q.
\end{equation}
The denominator is strictly positive, and a rational quotient would
contradict independence of those two coordinates. Thus the expected
totient, divisor-power sum, or any covered moment, even after
conditioning on an arbitrary residue class, is irrational. This is an
exact consequence of independence, not of a numerical approximation.

\subsection{A complete test for periodic polynomial numerators}
For each residue $r$, let $R_r,U_{k,r},V_{P,r}\in\Q[X]$, with only
finitely many nonzero arithmetic coefficients and canonical grouping
of equal polynomial Euler factors. Consider
\[
 S=\sum_{n\geq1}\left[R_{n\bmod q}(n)
       +\sum_kU_{k,n\bmod q}(n)\sigma_k(n)
       +\sum_PV_{P,n\bmod q}(n)F_P(n)\right]\frac{x^n}{n!}.
\]
Expand the ordinary polynomial part in falling factorials,
$R_r(X)=\sum_j c_{r,j}(X)_{\underline j}$, and set
\[
 C_s(x)=\sum_{j\geq0}c_{s+j\bmod q,j}\,x^j
       \qquad(0\leq s<q).
\]
Then $S$ is rational if and only if every $U_{k,r}$ and $V_{P,r}$
is zero and every $C_s(x)$ is zero. In that case $S=-R_0(0)$.
Indeed the ordinary polynomial contribution is exactly
\[
 \sum_{s=0}^{q-1}C_s(x)E_{q,s}(x)-R_0(0),
\]
because $(n)_{\underline j}/n!=1/(n-j)!$ for $n\geq j$.
Apply~\eqref{eq:periodic-family} to all remaining moment coordinates.
The same test applies to any additional base arithmetic family whose
complete moment family satisfies the transfer hypothesis.

\subsection{All derivatives and rational differential relations}
Let $\mathcal B$ be a base arithmetic family whose full moment family
$\{n\mapsto n^df(n):f\in\mathcal B,d\geq0\}$ satisfies
Theorem~\ref{thm:periodic-transfer}. For fixed $q$ and permitted $x$,
\[
 \{1,E_{q,r}(x):0\leq r<q\}\ \cup\
 \{B_{f,q,r}^{(j)}(x):f\in\mathcal B,\ 0\leq r<q,\ j\geq0\}
\]
is linearly independent over $\Q$. In fact
\[
 B_{f,q,r}^{(j)}(x)
  =x^{-j}\sum_{\substack{n\geq1\\n\equiv r\ (q)}}
                  (n)_{\underline j}f(n)\frac{x^n}{n!},
\]
and the falling factorials give an invertible triangular change from
ordinary moments, separately for each base function and residue.

Consequently the analytic functions $1$, all $E_{q,r}(z)$, and all
$B_{f,q,r}^{(j)}(z)$ are linearly independent over $\Q(z)$.
To prove this, evaluate a putative finite rational-function relation at
$x=2/b$, avoiding the finitely many poles and the finitely many zeros
of one nonzero coefficient. The value relation contradicts the
preceding independence. In particular no covered base residue function
satisfies a nontrivial linear differential equation over $\Q(z)$,
even with an inhomogeneous term in the $\Q(z)$-span of the exponential
residue sums and $1$. This differential consequence does not exclude
nonlinear differential equations, and no separate historical priority
claim is made for the special non-D-finiteness consequences.


\section{A logarithmic channel: the divisor-counting function}
\label{sec:tauchannel}
The divisor-counting function $\tau(n)=\sum_{d\mid n}1$ has a logarithmically
divergent mean. It therefore does not satisfy the finite-mean hypothesis of
Theorem~\ref{thm:transfer}. Nevertheless its entire polynomial-moment family
can be added, by separating logarithmic growth before the finite local means.

\begin{lemma}[Logarithmic progression mean]\label{lem:taumean}
For positive integers $Q,a$, put
\[
 L_\tau(Q,a)=\frac1Q\sum_{r=1}^Q
       \gcd(r,Q)\,\mathbf1_{\gcd(r,Q)\mid a}>0.
\]
Then
\[
 \frac1X\sum_{v=0}^{X-1}\tau(Qv+a)
       =L_\tau(Q,a)\log X+O_{Q,a}(1).
\]
For every prime $\ell\nmid Q$ and $a'\equiv a\pmod Q$,
\begin{equation}\label{eq:tauresponse}
 L_\tau(Q\ell,a')=L_\tau(Q,a)
      \bigl(1-\ell^{-1}+\mathbf1_{\ell\mid a'}\bigr).
\end{equation}
\end{lemma}
\begin{proof}
The bounded $Q$-periodic function
$w(d)=\gcd(d,Q)\mathbf1_{\gcd(d,Q)\mid a}$ has mean $L_\tau(Q,a)$.
Partial summation of its bounded partial-sum discrepancy gives
\[
 \sum_{d\leq Z}\frac{w(d)}d=L_\tau(Q,a)\log Z+O_Q(1).
\]
Divisor pairing, with a correction of at most $\sqrt Y$ for squares, gives
\[
 \sum_{\substack{n\leq Y\\n\equiv a\ (Q)}}\tau(n)
  =2\sum_{d\leq\sqrt Y}
       \#\{m\geq d:dm\leq Y,\ dm\equiv a\pmod Q\}+O(\sqrt Y).
\]
Writing $g=\gcd(d,Q)$, the inner count is zero unless $g\mid a$;
otherwise it equals $(Y/d-d)g/Q+O(1)$. Thus the last display equals
\[
 \frac{2Y}{Q}\sum_{d\leq\sqrt Y}\frac{w(d)}d+O_Q(Y)
   =\frac YQ L_\tau(Q,a)\log Y+O_Q(Y).
\]
Taking $Y=Q(X-1)+a$ proves the claimed average, up to a bounded initial
segment. The terms with $\gcd(r,Q)=1$ prove positivity. Finally CRT
makes residues modulo $Q$ and $\ell$ independent in the defining finite
average: the new prime contributes weight $1$ on its $\ell-1$ nonzero
residues and weight $\ell\mathbf1_{\ell\mid a'}$ on its zero residue.
Their average is precisely the multiplier in~\eqref{eq:tauresponse}.
\end{proof}

\begin{lemma}[Separation of a logarithmic phase]\label{lem:tauseparation}
Let $\rho_1,\ldots,\rho_s$ have finite means on every arithmetic
progression. A finite expression
\[
 V(N)=\sum_r b_r\tau(N+r)+\sum_{u,r}c_{u,r}\rho_u(N+r)
\]
cannot have a finite limit on $N\equiv A\pmod Q$ if, after grouping
repeated shifts, at least one $b_r$ is nonzero.
\end{lemma}
\begin{proof}
Choose $r_*$ with $b_{r_*}\ne0$ and a fresh prime larger than the number
of shifts and all their nonzero absolute differences. Compare an off-residue
avoiding every $-r$ with an on-residue equal to $-r_*$. On either refined
progression the mean of $V$, divided by $\log X$, tends to zero if $V$
converges. The finite-mean terms disappear. By Lemma~\ref{lem:taumean},
subtracting the two resulting equalities gives
$b_{r_*}L_\tau(Q,A+r_*)=0$, contrary to positivity. Positive representatives
of the starting residues may be used throughout.
\end{proof}

\begin{theorem}[Adjoining all divisor-counting moments]\label{thm:taumoments}
At each fixed $x=\pm2/b$, $b\geq1$, the family in
Theorem~\ref{thm:dyadic} remains linearly independent over $\Q$ after
adjoining all the values
\[
 T_d(x)=\sum_{n\geq1}\frac{n^d\tau(n)x^n}{n!},\qquad d=0,1,2,\ldots.
\]
More generally, let $\mathcal F$ be a family satisfying the hypotheses of
Theorem~\ref{thm:periodic-transfer}, and suppose all its normalized phases
have finite progression means. For each common period $q$, its independent
family remains independent after adjoining all residue-class values
\[
 T_{d,r}^{[q]}(x)=
    \sum_{\substack{n\geq1\\n\equiv r\ (q)}}
          \frac{n^d\tau(n)x^n}{n!},
 \qquad d\geq0,\quad 0\leq r<q.
\]
The periodic base values are the $E_{q,r}(x)$ of that theorem; polynomial
weights on those base values are reduced by the stated finite identities.
\end{theorem}
\begin{proof}
In a hypothetical finite relation assign degree $d$ to $n^d\tau(n)$,
and retain the assigned degrees of the other arithmetic coordinates.
Let $K\geq0$ be the largest degree and use $H=K+1$. Include $g_0=1$ in
the fiber equations. The bound $\tau(n)\leq2\sqrt n$, positivity,
integer values, and multiplicativity supply every algebraic and growth
hypothesis used in the cancellation and dyadic-height steps.
Thus $W(N)=O(N^{-3/2})$, its denominator is $O(N)$, and $W(N)=0$
eventually. After multiplication by $N^2$, the top-degree normalized
shifted phases have a finite limit; every lower-degree phase vanishes.
If $K=0$, the constant phase contributes a finite constant rather than
vanishing, which is permitted here.

If there is no top-degree divisor-counting coordinate, use the finite-mean
separation argument. Otherwise it is the unique coordinate $n^K\tau(n)$.
At a vertex $e_*$ with nonzero grid weight, its isolated terminal-shift
coefficient is
\[
 \frac{a_\tau w_{e_*}p_{e_*}^{H+1}x^{H+1}}
      {p_{e_*}^K\tau(p_{e_*})}\ne0.
\]
All other top-degree phases have finite progression means.
Lemma~\ref{lem:tauseparation} supplies the contradiction.

For the general periodic statement, adjoin $n^d\tau(n)$ with assigned
degree $d$. At any fixed degree there is exactly one additional normalized
phase, namely $\tau$. If a grouped shifted combination contains a nonzero
$\tau$ coefficient, Lemma~\ref{lem:tauseparation} excludes a finite limit;
otherwise the assumed separation for $\mathcal F$ applies. Thus the
enlarged family satisfies every hypothesis of
Theorem~\ref{thm:periodic-transfer}, proving the assertion. Equivalently,
in its grid proof one chooses the progression residue $A$ so that the
nonzero top-degree periodic coefficient is nonzero at $A+H+1$; its isolated
terminal coefficient survives exactly as above.
\end{proof}

Let $T(z)=\sum_{n\geq1}\tau(n)z^n/n!$. Every derivative of $T$ at a
permitted point joins the independent derivative family, by the same
falling-factorial change of basis. The exact rationality criterion also
extends: a polynomial coefficient multiplying $\tau(n)$ must vanish,
as must every nonconstant arithmetic coordinate already present.
The corresponding statements hold residue by residue in the periodic
setting. These are joint results; no transcendence assertion is inferred.
Individual divisor-counting factorial irrationality has classical
precedents, so its specialization alone is not claimed as new.


\clearpage
\section{The full derivative--antiderivative hierarchy}
\label{sec:antiderivatives}

The preceding polynomial-moment results extend to genuine integration
constants.  The new ingredient is a grid whose coordinates include negative
indices, together with different truncation orders for different factorial
shifts.  A finite set of additional equations removes the exponential tail
to sufficiently high order.  These modifications are necessary: a shifted
factorial denominator is not a polynomial weight in the earlier theorem.

\begin{theorem}[Two-sided factorial hierarchy]
\label{thm:two-sided}
Let $\mathcal B$ be a family of positive integer-valued multiplicative
functions.  For each $f\in\mathcal B$ assign an integer $h_f\geq0$ such that
$f(n)=O_f(n^{h_f+1/2})$.  Suppose that the normalized phases
$\rho_f(n)=f(n)/n^{h_f}$, taken across the whole canonical base family,
have the isolated-shift separation property of
Theorem~\ref{thm:periodic-transfer}.  In particular, positive progression
means and pairwise distinct nonzero scalar responses as in
Lemma~\ref{lem:multimean} suffice.  A proved profile-vector separation
criterion, such as Lemma~\ref{lem:profile-separation}, also suffices.
For every fixed $q\geq1$ and $x=\varepsilon2/b$, $b\geq1$ an integer and
$\varepsilon\in\{-1,1\}$, the family
\begin{align}
 &\{1\}\cup\{E_{q,a}(x):0\leq a<q\}\notag\\
 &\ \cup\left\{\sum_{\substack{n\geq1\\n\equiv a\ (q)}}
                 \frac{n^d f(n)x^n}{n!}:
                 f\in\mathcal B,\ d\geq0,\ 0\leq a<q\right\}\notag\\
 &\ \cup\left\{\sum_{\substack{n\geq1\\n\equiv a\ (q)}}
                 \frac{f(n)x^n}{(n+r)!}:
                 f\in\mathcal B,\ r\geq1,\ 0\leq a<q\right\}
 \label{eq:two-sided-family}
\end{align}
is linearly independent over $\Q$.
\end{theorem}

The divisor functions $\sigma_k$ and the distinct polynomial Euler factors
$F_P$ satisfy these hypotheses by the preceding sections.  More generally,
the result applies to any additional family for which these hypotheses have
actually been verified. In particular, choose one arithmetic base
for each distinct nontrivial admissible profile and adjoin $\tau$ with
degree zero. If a shifted combination contains a nonzero $\tau$
coefficient, Lemma~\ref{lem:tauseparation} excludes a finite limit;
otherwise Lemma~\ref{lem:profile-separation} does. Thus this entire
enlarged base family satisfies the theorem, across all effective degrees.
The assertion is at one fixed $x$; arbitrary unrelated evaluation
points are not asserted to be independent.

\begin{proof}
We give the details because both negative shifts and arbitrarily high
integration orders matter.  Consider a finite proposed rational relation.
Its arithmetic channels have one of the canonical forms
\[
 (g,r)=(n^d f,0),\quad h=h_f+d,
 \qquad\hbox{or}\qquad
 (g,r)=(f,r),\quad h=h_f,\quad r\geq1.
\]
Define the effective degree of a channel to be $w=h-r$.
Group residue-class coefficients into periodic functions before assigning
channels. At a fixed effective degree each base function occurs at most
once, so the top normalized phases belong to the separated base family.
If every arithmetic coefficient vanishes, apply
Lemma~\ref{lem:periodic-exp}. Henceforth assume an arithmetic channel remains.
Retain the channels with nonzero coefficients, put
\[
 R=\max r,\qquad w_* = \max(h-r),\qquad
 J=\max(0,w_*+1),\qquad M=J+R.
\]
If there is an arithmetic channel then $M\geq1$.
Reindexing gives
\[
 \sum_{n\geq1}\frac{g(n)x^n}{(n+r)!}
 =x^{-r}\sum_{m\geq r+1}\frac{g(m-r)x^m}{m!}.
\]
Clear all rational coefficients, including the factors $x^{-r}$.
Thus the putative rational relation is an integer-coefficient factorial
series with a $q$-periodic exponential coefficient $a_0(m)$ and
nonzero $q$-periodic arithmetic coefficient functions $a_{g,r}(m-r)$.
These functions are evaluated at the original arithmetic argument.
Write the rational value as $u/v$.
For its normalized tail $T(n)$, Lemma~\ref{lem:dyadic-height} gives
$v2^{s_2(n)}T(n)\in\Z$.
Choose a top effective-degree channel and a residue $A$ such that
$a_{g,r}(A+J+1)\ne0$. Fix this residue before solving for the weights.

\emph{The offset grid.}
Index the coordinates by $\ell=-R,\ldots,J-1$ and let
$e_\ell\in\{0,\ldots,L-1\}$, where $L$ will be chosen below. Put
\begin{gather*}
 b_0=M(L-1)+1,\quad F=\max(b_0^M-1,M+2),\quad D=qF!,
 \quad C=\max(R,J)+1,\\
 I_e=\sum_{\ell=-R}^{J-1}e_\ell b_0^{\ell+R},\qquad
 T_e=\sum_{\ell=-R}^{J-1}(\ell+1)e_\ell b_0^{\ell+R},\\
 B=1+(J+2C+M+2)DF,\qquad
 p_e=B+DI_e,\qquad t_e=DT_e,\\
 s_{e,\ell}=(\ell+1)p_e-t_e\qquad(-R\leq\ell\leq J).
\end{gather*}
The integers $t_e$ and some lower shifts may be negative; all arguments
below are positive for sufficiently large $N$.
Exactly as in Lemma~\ref{lem:grid}, the $p_e$ are pairwise coprime,
and each is coprime to every integer at most $M+2$.  A lower shift
$s_{e,\ell}$ is constant on its coordinate-$\ell$ fiber.  Moreover,
\[
 s_{e,J}=(J+1)B+
 D\sum_{\ell=-R}^{J-1}(J-\ell)e_\ell b_0^{\ell+R}.
\]
Its displayed digits lie between $0$ and $M(L-1)=b_0-1$, so the
terminal shifts are distinct.  They are at least $(J+1)B$, whereas
every lower shift is at most $JB+(J+C)DF<(J+1)B$.
Here we used $0\leq I_e\leq F$ and $|T_e|\leq CF$.

We also have $p_e\equiv1\pmod q$ and $t_e\equiv0\pmod q$.
Take the CRT progression specified by
\[
 N\equiv t_e+(R+1)p_e\pmod {p_e^2}\quad\hbox{for every }e,
 \qquad N\equiv A\pmod q.
\]
Put $n_e=(N-t_e)/p_e$.
Then $n_e\equiv R+1\pmod{p_e}$ and
\[
 \gcd(p_e,n_e+\ell+1)=\gcd(p_e,R+\ell+2)=1,
 \qquad -R\leq\ell\leq J,
\]
since $2\leq R+\ell+2\leq M+2$.  In particular the otherwise
problematic shift $\ell=-1$ is covered.  Multiplicativity therefore gives
the exact rescaling identity
\[
 \frac{g(n_e+\ell+1)}{(n_e+\ell+1)^{i+1}}
 =\frac{p_e^{i+1}}{g(p_e)}
   \frac{g(N+s_{e,\ell})}{(N+s_{e,\ell})^{i+1}}.
\]

Moreover $n_e\equiv A\pmod q$, so the periodic coefficient of the term
at index $\ell$ is $a_{g,r}(A+\ell+1)$, independent of the vertex.
Every fiber cancellation below therefore remains exact.

\emph{Different truncation orders.}
For a channel $(g,r)$ expand its reindexed tail through order $J+r$.
If $y=n+j+1-r$, the relevant denominator is
\[
 \frac{n!}{(n+j+1)!}
 =\frac1{(y+r)(y+r-1)\cdots(y+r-j)}.
\]
Its expansion in $y^{-1}$ has integral coefficients $c_{r,j,i}$,
with $c_{r,j,j}=1$.  The same geometric remainder argument as in
Lemma~\ref{lem:tail} gives
\begin{align}
 T_{g,r}(n)&:=n!x^{-n}
       \sum_{m>n}\frac{g(m-r)x^m}{m!}\notag\\
 &=\sum_{j=0}^{J+r}\sum_{i=j}^{J+r}
       x^{j+1}c_{r,j,i}
       \frac{g(n+j+1-r)}{(n+j+1-r)^{i+1}}
       +O(n^{h-J-r-3/2}).
 \label{eq:offset-tail}
\end{align}
For example, $c_{r,j,i}$ is $(-1)^{i-j}$ times the complete
homogeneous polynomial of degree $i-j$ in $r,r-1,\ldots,r-j$.
The omitted tail has the same bound.  The error in
\eqref{eq:offset-tail} is $O(n^{w-J-3/2})$.

For every term with $\ell=j-r<J$, impose on every coordinate-$\ell$
fiber the rational equation
\[
 \sum_{e_\ell}v_e\frac{p_e^{i+1}}{g(p_e)}=0.
\]
There are a fixed number $C_0$ of coefficient arrays, hence at most
$C_0L^{M-1}$ equations.  These equations cancel every displayed term
except $j=i=J+r$, whose shift is the common terminal index $\ell=J$.

\emph{The exponential at negative effective degrees.}
One cannot simply discard the exponential when $w_*<0$.
For any fixed integer $T\geq2$, the periodic exponential tail, restricted
to $n\equiv A\pmod q$, has a rational expansion
\[
 U_x(n)=n!x^{-n}\sum_{m>n}\frac{a_0(m)x^m}{m!}
       =\sum_{a=1}^{T-1}u_a n^{-a}+O(n^{-T}),\qquad u_a\in\Q.
\]
Substituting $n_e=(N-t_e)/p_e$ yields
\[
 U_x(n_e)=\sum_{v=1}^{T-1}K_v(p_e,t_e)N^{-v}+O(N^{-T}),
\quad
 K_v(p,t)=\sum_{a=1}^v u_a\binom{v-1}{v-a}p^a t^{v-a}.
\]
Set $T=\max(2,J+2-w_*)$ and add the $T-1$ global equations
\[
 \sum_e v_e K_v(p_e,t_e)=0\qquad(1\leq v<T).
\]
Choose $L>C_0+T$.  The total number of equations is less than
$L^M$, so rank--nullity supplies a nonzero integer vector $(v_e)$
satisfying all the fiber and global constraints.  All its data are fixed
before $N$ tends to infinity.

\emph{The integer limit.}
The weighted combined tail $W(N)=\sum_e v_eT(n_e)$ now satisfies
\begin{align*}
 W(N)={}&\sum_{g,r}\sum_e
 \frac{a_{g,r}(A+J+1)x^{J+r+1}v_ep_e^{J+r+1}}{g(p_e)}
 \frac{g(N+s_{e,J})}{(N+s_{e,J})^{J+r+1}}\\
 &\quad+O(N^{w_*-J-3/2})+O(N^{-T}).
\end{align*}
The terminal terms are $O(N^{w_*-J-1/2})$, so $W(N)=O(N^{-3/2})$.
Its denominator divides
\[
 v2^{\max_e s_2(n_e)}=O(N).
\]
Multiplication by this integer gives an integer tending to zero.
Consequently $W(N)=0$ eventually.

Multiply by $N^{J+1-w_*}$.  The errors tend to zero, every smaller
effective degree disappears, and replacing the fixed shifted power
ratio by $1$ costs $O(N^{-1/2})$.  We obtain
\[
 \sum_{h-r=w_*}\sum_e
 \frac{a_{g,r}(A+J+1)x^{J+r+1}v_ep_e^{J+r+1}}{g(p_e)}
 \rho_g(N+s_{e,J})\longrightarrow0.
\]
Choose a vertex with $v_e\ne0$.  Its terminal shift is isolated,
and the top channel used to choose $A$ has a nonzero coefficient there.
The separation property of the canonical bases is a contradiction.
If the original relation had no arithmetic channel, use
Lemma~\ref{lem:periodic-exp}.  This proves the theorem.
\end{proof}

\begin{corollary}[All derivatives and all normalized antiderivatives]
\label{cor:two-sided-jet}
First take $q=1$.
For $f\in\mathcal B$, put $A_f(z)=\sum_{n\geq1}f(n)z^n/n!$.
For $r\geq1$, let
\[
 A_f^{[-r]}(x)=\frac1{(r-1)!}\int_0^x(x-t)^{r-1}A_f(t)\,dt
             =\sum_{n\geq1}\frac{f(n)x^{n+r}}{(n+r)!}.
\]
For nonnegative $j$ use $A_f^{[j]}=A_f^{(j)}$.
At every fixed permitted $x$, the family
\[
 \{1,e^x\}\ \cup\
 \{A_f^{[j]}(x):f\in\mathcal B,\ j\in\Z\}
\]
is linearly independent over $\Q$.
For general fixed $q$, the same assertion holds with all residue-restricted
functions
\[
 A_{f,q,a}(z)=\sum_{\substack{n\geq1\\n\equiv a\ (q)}}
                  \frac{f(n)z^n}{n!}\qquad(0\leq a<q)
\]
in place of $A_f$, and with $E_{q,0}(x),\ldots,E_{q,q-1}(x)$ in place
of $e^x$.
\end{corollary}
\begin{proof}
Negative orders multiply the second family in
\eqref{eq:two-sided-family} by the nonzero rational number $x^r$.
For nonnegative orders, the falling-factorial change of basis from
Corollary~\ref{cor:derivatives} is triangular and invertible over $\Q$.
Apply Theorem~\ref{thm:two-sided} to the finite set of coordinates in
any proposed relation.
The change of basis acts separately on each residue-restricted function,
so the same proof establishes the last assertion.
\end{proof}


\begin{corollary}[The two-sided hierarchy in four real directions]
\label{cor:four-two-sided}
For a positive permitted $x=2/b$, the following real numbers are
jointly linearly independent over $\Q$:
\begin{gather*}
 1,\ e^x,\ e^{-x},\ \cos x,\ \sin x,\\
 A_f^{[j]}(x),\quad A_f^{[j]}(-x),\quad
 \operatorname{Re}A_f^{[j]}(ix),\quad
 \operatorname{Im}A_f^{[j]}(ix)
 \qquad(f\in\mathcal B,\ j\in\Z).
\end{gather*}
For negative orders, the antiderivative is normalized to vanish at
zero together with its lower derivatives; its entire series defines
it also at complex arguments.
\end{corollary}
\begin{proof}
Use the residue-restricted statement of Corollary~\ref{cor:two-sided-jet}
with modulus four. For $\zeta\in\{1,-1,i\}$, the power series gives
\[
 A_f^{[j]}(\zeta x)=\zeta^{-j}
       \sum_{a=0}^3\zeta^a A_{f,4,a}^{[j]}(x),\qquad j\in\Z.
\]
This identity is valid for derivatives and normalized antiderivatives
alike. For each fixed $j$, its real transformation is the invertible
rational matrix in Corollary~\ref{cor:four-directions}, followed by
sign changes and possibly interchanging the real and imaginary rows.
The exponential block uses the same matrix. Independence is therefore
preserved block by block in every finite proposed relation.
\end{proof}

\begin{corollary}[Rational weights with simple negative-integer poles]
\label{cor:simple-poles}
Let finitely many $R_f\in\Q(X)$ have arbitrary polynomial parts and
at most simple poles, all among $-1,-2,\ldots$.  If at least one
$R_f$ is nonzero, then, for every $\lambda,q\in\Q$,
\[
 q+\lambda e^x+
 \sum_f\sum_{n\geq1}\frac{R_f(n)f(n)x^n}{n!}\notin\Q.
\]
In particular the quantities
\[
 M_{f,a}(x):=\int_0^1t^{a-1}A_f(xt)\,dt
           =\sum_{n\geq1}\frac{f(n)x^n}{n!(n+a)},\qquad a\geq1,
\]
are jointly independent together with $1,e^x$ and every polynomial
moment in the first family of \eqref{eq:two-sided-family}.
\end{corollary}
\begin{proof}
For each $R\geq1$, the rational functions
\[
 \frac1{(X+1)(X+2)\cdots(X+r)},\qquad 1\leq r\leq R,
\]
form a basis of the rational functions vanishing at infinity with
at most simple poles in $\{-1,\ldots,-R\}$.  Indeed the residue at
$-R$ first determines the last coefficient, and downward induction
determines the others; the dimension is $R$ by partial fractions.
Combining this basis with ordinary polynomials converts the claim
to Theorem~\ref{thm:two-sided}.  Termwise integration is justified
by absolute uniform convergence on the compact integration interval.
For the last assertion, the functions $(X+a)^{-1}$ are themselves
linearly independent, again by their residues.
\end{proof}

The rational-weight statement also holds separately and jointly on all
residue classes of any fixed modulus. Apply the same basis conversion
to each $R_{f,a}$, and then use the residue-restricted family
\eqref{eq:two-sided-family}.  Consequently the class permits, at once,
periodic signs, residue restrictions, polynomial weights, simple
negative-integer poles, and canonical arithmetic bases with a proved
separation property.

For example, the theorem includes simultaneously
\[
 \int_0^1t^6A_5(2t)\,dt,
 \qquad \int_0^2(2-t)^9A_5(t)\,dt,
 \qquad A_5^{(17)}(2),
\]
and their counterparts for every admissible arithmetic base, after
the canonical basis conversion just proved.  The displayed integral
with $(2-t)^9$ is $9!A_5^{[-10]}(2)$.  Independence must always be
checked after this conversion: listing two different integral
representations of the same linear functional does not create two
independent coordinates.

This extension does not include a pole at zero, such as the weight
$1/n$, or repeated poles such as $(n+1)^{-2}$.  It also does not infer
irrationality of an arbitrary integral from pointwise irrationality
of its integrand.  Here the precise factorial coefficients and the
new tail argument prove the integral statements directly.  Their
historical priority, including classical special cases of low degree,
requires a separate literature review.


\section{Consequences across arithmetic, lattices, and finite groups}
\label{sec:applications}

The preceding theorems control several operations together: products
of arithmetic functions, residue restrictions, differentiation, and
the specific integral transforms proved in
Corollary~\ref{cor:simple-poles}.  The following examples make the
scope concrete.  The underlying counting functions are classical;
their factorial evaluations are the objects studied here.

\subsection{Cyclic subgroups of finite Abelian groups}

For $r\geq2$ let
\[
 c_r(n)=\#\{\text{cyclic subgroups of }(\Z/n\Z)^r\}.
\]
The identity
\begin{equation}\label{eq:cyclic-count}
 c_r(n)=\sum_{d\mid n}\frac{J_r(d)}{\varphi(d)}
\end{equation}
is a specialization of the subgroup formulas in~\cite{Toth12}.
It also has a short direct proof.  For $d\mid n$, the number of
vectors whose order divides $d$ is $d^r$.  M\"obius inversion shows
that exactly $J_r(d)$ vectors have order $d$.  Each cyclic subgroup
of that order has $\varphi(d)$ generators, proving
\eqref{eq:cyclic-count}, including the trivial subgroup at $d=1$.

The summand $J_r/\varphi$ is multiplicative, so $c_r$ is
multiplicative and integer-valued.  Put $h=r-1$ and
\[
 D_r(z)=\frac{1-z^r}{(1-z)(1-z^h)}.
\]
At a prime power, summing the geometric progression gives
\[
 c_r(p^a)=1+\frac{p^r-1}{p-1}
                       \sum_{j=0}^{a-1}p^{hj},\qquad
 \frac{c_r(p^a)}{p^{ha}}
       =D_r(z)+(1-D_r(z))w^h,
 \quad z=p^{-1},\quad w=p^{-a}.
\]
For $0<z\leq1/2$, $D_r(z)$ is positive and bounded.  The last
expression is a convex combination of $D_r(z)$ and $1$, so it is
positive and uniformly bounded for all $a\geq1$; at $w=z$ it is
$1+O(z)$.  It is therefore an admissible, nonconstant profile.
Different $r$ give different degrees in $w$.  Consequently
Theorem~\ref{thm:profile-transfer} covers all $c_r$ and all their
polynomial moments jointly with $1,e^x$; any further profiles may
be included after the stated equality test.  Their fixed finite
pointwise products are also admissible.

\subsection{Explicit constants and one simultaneous statement}

To keep the formulas readable, write
\[
 \mathcal E[f]=\sum_{n\geq1}\frac{2^nf(n)}{n!},
 \qquad A_5(z)=\sum_{n\geq1}\frac{\sigma_5(n)z^n}{n!}.
\]
Every entry in the left column below is irrational.  The second
column identifies the coefficients or the proved operation.
\begin{center}
\small
\begin{tabular}{p{0.46\textwidth}p{0.45\textwidth}}
\toprule
Constant & Arithmetic or geometric meaning\\
\midrule
$\mathcal E[\operatorname{rad}]$
 & Product of the distinct prime factors of each index\\[4pt]
$\mathcal E[\sigma_1^2]$
 & Square of the divisor sum at each index\\[4pt]
$\mathcal E[\varphi^3\sigma_5^2]$
 & A mixed product of totients and divisor sums\\[4pt]
$\mathcal E[a_3]$
 & Index-$n$ sublattices of $\Z^3$~\cite{Zou06}\\[4pt]
$\mathcal E[c_3]$
 & Cyclic subgroups of $(\Z/n\Z)^3$~\cite{Toth12}\\[4pt]
$\mathcal E[L]$, $L(n)=\sum_{j=1}^n\operatorname{lcm}(j,n)$
 & The classical least-common-multiple sum~\cite{IM14}\\[4pt]
$\displaystyle\sum_{\substack{n\geq1\\ n\equiv1\!\!\pmod3}}
             \frac{2^n\sigma_5(n)}{n!}$
 & One arithmetic progression of indices\\[9pt]
$\displaystyle\int_0^1t^6A_5(2t)\,dt
       =\sum_{n\geq1}\frac{2^n\sigma_5(n)}{(n+7)n!}$
 & A rational-weight integral covered by the new tail argument\\
\bottomrule
\end{tabular}
\end{center}
The first six follow from Theorem~\ref{thm:profile-transfer} and
\eqref{eq:lcm-constant}; the last two follow from
Theorem~\ref{thm:periodic-transfer} and
Corollary~\ref{cor:simple-poles}, respectively.  None of these
statements requires numerical approximation to certify irrationality.

A stronger assertion connects four of these arithmetic settings
with the integral transform.  The following seven numbers are
linearly independent over $\Q$:
\begin{equation}\label{eq:seven-crossfield}
 \boxed{\begin{gathered}
 1,\ e^2,\ A_5(2),\ \int_0^1t^6A_5(2t)\,dt,\\[3pt]
 \mathcal E[a_3],\ \mathcal E[\operatorname{rad}],\
 \mathcal E[L].
 \end{gathered}}
\end{equation}
Indeed the profiles of $\sigma_5$, $a_3$, $\operatorname{rad}$,
and $g(n)=\sum_{d\mid n}d\varphi(d)$ are distinct.  Their degrees
in $w$ are $5,2,1,2$; the $a_3$ profile has a nonzero coefficient
of $w$, whereas the $g$ profile has none.  Within the $\sigma_5$
coordinate, the rational weights $1$ and $(n+7)^{-1}$ are
independent.  Finally, \eqref{eq:lcm-sum-identity} gives
\[
 \mathcal E[L]=e^2+\tfrac12\mathcal E[n g],
\]
an invertible rational change of coordinates in a family supplied
by Corollary~\ref{cor:simple-poles}.  This proves
\eqref{eq:seven-crossfield}; in particular every rational linear
combination having a nonzero coefficient on one of the six
nonconstant coordinates is irrational.

The square in $\mathcal E[\sigma_1^2]$ acts on each coefficient:
it does not denote $\mathcal E[\sigma_1]^2$.  Products of the
resulting constants, algebraic independence, and transcendence
require different arguments.  Likewise, the exact appearance of
$\zeta(3)/\zeta(2)$ in the normalized mean of $g$, established
above, does not transfer this conclusion to that zeta ratio.
Classical low-growth factorial results~\cite{HT10} remain part of
the prior literature.  The catalog illustrates the range of the
proved theorems; it does not certify that each displayed constant
was previously a named open problem.


\section{Polynomial valuation profiles and every logarithmic order}
\label{sec:v5-log-profiles}

The finite-mean profiles and the separate divisor-counting channel fit into
one larger class. The enlargement admits polynomial dependence on a prime
exponent, arbitrarily high powers of the divisor-counting function, all
ordered-factorization functions, and products of these with the earlier
arithmetic bases. We prove the necessary progression averages here; no
unproved distribution hypothesis or general Selberg--Delange theorem is
needed.

\begin{definition}[An admissible valuation profile]
\label{def:v5-log-profile}
Let $\mathcal R(z,w,u)\in\Q(z)[w,u]$. Its coefficient functions have
no poles at $z=1/p$ for primes $p$. We call $\mathcal R$ admissible if:
\begin{enumerate}
\item there are constants $C,D$ such that, for every prime $p$ and
integer $a\geq1$,
\[
 0<\mathcal R(p^{-1},p^{-a},a)\leq C(1+a)^D;
\]
\item for an integer $\kappa_{\mathcal R}\geq1$,
\[
 \mathcal R(z,z,1)=\kappa_{\mathcal R}+O(z)\qquad(z\longrightarrow0);
\]
\item for some integer $h\geq0$, every
$p^{ha}\mathcal R(p^{-1},p^{-a},a)$ is a positive integer.
\end{enumerate}
Set $\rho_{\mathcal R}(1)=1$ and extend
\[
 \rho_{\mathcal R}(p^a)=\mathcal R(p^{-1},p^{-a},a)
\]
multiplicatively. An integral realization is
$f_{\mathcal R,h}(n)=n^h\rho_{\mathcal R}(n)$.
The value at exponent zero is defined to be $1$; it need not equal
$\mathcal R(z,1,0)$. Distinct profiles always mean distinct elements
of $\Q(z)[w,u]$ after simplification.
\end{definition}

The integer $\kappa_{\mathcal R}-1$ is the order of logarithmic growth
of the progression mean. The former admissible profiles are the
$u$-independent case with $\kappa=1$. The profile $u+1$ gives $\tau$
and has $\kappa=2$. The constant profile $c\geq2$ gives
$c^{\omega(n)}$, where $\omega(n)$ counts distinct prime divisors:
this example explains the exponent-zero convention.

\begin{lemma}[Growth and an absolutely convergent convolution]
\label{lem:v5-log-convolution}
For an admissible profile and every $\epsilon>0$,
$\rho_{\mathcal R}(n)=O_{\mathcal R,\epsilon}(n^\epsilon)$.
Writing $\kappa=\kappa_{\mathcal R}$, let
\[
 g=\rho_{\mathcal R}*\mu^{*\kappa},
 \qquad \rho_{\mathcal R}=d_\kappa*g,
\]
where $*$ denotes Dirichlet convolution, $\mu^{*\kappa}$ its
$\kappa$-fold power, and $d_\kappa=1^{*\kappa}$.
For every real $\sigma>1/2$,
\[
 \sum_{n\geq1}\frac{|g(n)|}{n^\sigma}<\infty.
\]
Moreover, with
\[
 \mathcal L_{\mathcal R}(z)
   =1+\sum_{a\geq1}\mathcal R(z,z^a,a)z^a,
 \qquad U_{\mathcal R}(z)=(1-z)\mathcal L_{\mathcal R}(z),
\]
these are rational functions of $z$, and
\[
 G_{\mathcal R}
 =\sum_{n\geq1}\frac{g(n)}n
 =\prod_p(1-p^{-1})^\kappa\mathcal L_{\mathcal R}(p^{-1})>0.
\]
\end{lemma}
\begin{proof}
The polynomial bound in $a$ implies that
$B=\sup_{a\geq1}[C(1+a)^D]^{1/a}$ is finite. If
$p\geq B^{1/\epsilon}$, then $\rho(p^a)\leq p^{\epsilon a}$.
For each of the finitely many remaining primes, the supremum of
$\rho(p^a)p^{-\epsilon a}$ over $a\geq0$ is finite. Multiplying those
finitely many bounds proves the growth assertion.

At a prime, $g(p)=\rho(p)-\kappa=O(p^{-1})$. At higher powers,
\[
 g(p^a)=\sum_{j=0}^{\min(a,\kappa)}
             (-1)^j\binom\kappa j\rho(p^{a-j})
             =O_{\mathcal R}((1+a)^D).
\]
Thus, for fixed $\sigma>1/2$,
\[
 \sum_p\sum_{a\geq1}|g(p^a)|p^{-a\sigma}<\infty.
\]
The absolutely convergent Euler product gives the asserted Dirichlet
series. Each monomial $r(z)w^ju^m$ contributes
$r(z)\sum_{a\geq1}a^m z^{(j+1)a}$ to $\mathcal L-1$.
The series $\sum_{a\geq1}a^m v^a$ is rational, as follows by repeatedly
applying $v\,d/dv$ to $v/(1-v)$. This proves rationality.
Finally every Euler factor defining $G$ is positive and equals
$1+O(p^{-2})$. The convergent product is therefore strictly positive.
\end{proof}

\begin{lemma}[All progression means and their local factors]
\label{lem:v5-log-profile-mean}
For every fixed positive $Q,a$,
\[
 \frac1X\sum_{v=0}^{X-1}\rho_{\mathcal R}(Qv+a)
  =L_{\mathcal R}(Q,a)(\log X)^{\kappa-1}
    +\begin{cases}o(1),&\kappa=1,\\
                   O_{\mathcal R,Q,a}((\log X)^{\kappa-2}),&\kappa\geq2,
      \end{cases}
\]
where $L_{\mathcal R}(Q,a)>0$. More explicitly, put
$R_p(t)=\rho_{\mathcal R}(p^t)$, including $R_p(0)=1$, and $z=1/p$.
For $p^v\Vert Q$ define
\[
 A_{\mathcal R,p^v}(a)=
 \begin{cases}
 \displaystyle\frac{R_p(t)}{U_{\mathcal R}(z)},
       &t=v_p(a)<v,\\[6pt]
 \displaystyle\frac{\sum_{j\geq0}R_p(v+j)z^j}
                       {\mathcal L_{\mathcal R}(z)},
       &p^v\mid a.
 \end{cases}
\]
Then
\[
 L_{\mathcal R}(Q,a)
  =\frac{G_{\mathcal R}}{(\kappa-1)!}
       \prod_{p^v\Vert Q}A_{\mathcal R,p^v}(a).
\]
In particular, at a fresh prime $\ell\nmid Q$, refinement to a unit
class modulo $\ell^{t+1}$ multiplies $L$ by $1/U_{\mathcal R}(1/\ell)$,
whereas refinement to exact valuation $t\geq1$ multiplies it by
$\mathcal R(1/\ell,\ell^{-t},t)/U_{\mathcal R}(1/\ell)$.
\end{lemma}
\begin{proof}
We supply an elementary progression argument to include nonunit classes.
First,
\[
 \sum_{n\leq Y}d_\kappa(n)
   =\frac{Y(\log Y)^{\kappa-1}}{(\kappa-1)!}
       +O_\kappa(Y(1+\log Y)^{\kappa-2})\quad(\kappa\geq2),
\]
while the sum for $\kappa=1$ is $Y+O(1)$.
For $\kappa=2$ this follows from the harmonic sum. Induction uses
$d_\kappa=1*d_{\kappa-1}$ and
\[
 \sum_{m\leq Y}\frac{(\log(Y/m))^j}{m}
       =\frac{(\log Y)^{j+1}}{j+1}+O_j((1+\log Y)^j).
\]
The displayed estimate follows by comparison with the integral.

For a nonprincipal Dirichlet character $\chi$ modulo a fixed integer,
its partial sums are bounded. The $\kappa$-fold hyperbola argument gives
\[
 \sum_{n\leq Y}\chi(n)d_\kappa(n)
       =O_{\chi,\kappa}(Y^{1-1/\kappa})\quad(\kappa\geq1),
\]
where the exponent is zero when $\kappa=1$.
For completeness, in the sum over $\kappa$ factors with product at most
$Y$, at least one factor is at most $Y^{1/\kappa}$.
Use inclusion--exclusion on these $\kappa$ conditions.
For a term with $j<\kappa$ restricted factors, induction on
$\ell=\kappa-j$ bounds the remaining convolution by
$O((Y/(m_1\cdots m_j))^{1-1/\ell})$. Summing absolute values over
$m_i\leq Y^{1/\kappa}$ gives
$O(Y^{1-1/\ell+j/(\kappa\ell)})=O(Y^{1-1/\kappa})$.
The term with all factors restricted is a product of bounded character
sums. This proves the estimate.

For the principal character modulo $Q$, removing the primes dividing
$Q$ from $d_\kappa$ is convolution with the finite Euler polynomial
$\prod_{p\mid Q}(1-p^{-s})^\kappa$. Its leading coefficient is
therefore $\prod_{p\mid Q}(1-p^{-1})^\kappa$.
Now convolve with $g(n)\chi(n)$ from
Lemma~\ref{lem:v5-log-convolution}. For a nonprincipal character choose
$\beta$ with $\max(1/2,1-1/\kappa)<\beta<1$; the character sum is
$O(Y^\beta)$. For the principal character, absolute convergence of
$\sum |g(n)|n^{-\beta}$ also gives all the logarithmic moments
$\sum |g(n)|(\log n)^j/n<\infty$ needed to expand the leading term.
The convolution consequently preserves the stated main term and error.
Character orthogonality now proves, when $(a,Q)=1$, the asserted mean
with coefficient
\[
 \frac{G_{\mathcal R}}{(\kappa-1)!}
       \prod_{p\mid Q}\frac1{U_{\mathcal R}(1/p)}.
\]
For $\kappa=1$ the same proof gives a power-saving error before division
by $X$, which in particular is $o(1)$ afterwards.

For a general class put $c=(a,Q)$, $Q'=Q/c$, and $a'=a/c$.
Then $(a',Q')=1$ and the function
\[
 \rho_c(m)=\frac{\rho_{\mathcal R}(cm)}{\rho_{\mathcal R}(c)}
\]
is multiplicative. Its local factors differ from those of $\rho$ only
at primes dividing $c$, where its values are $R_p(v_p(c)+j)/R_p(v_p(c))$.
Those finitely many modified factors still have polynomial growth in $j$;
the convolution proof just given applies unchanged to them.
Use the unit-class result on $Q'v+a'$ and multiply by $\rho(c)$.
If $v_p(a)<v_p(Q)$, the modified prime is removed from the unit-class
Euler product, leaving $R_p(v_p(a))/U_{\mathcal R}(1/p)$.
If $p^{v_p(Q)}\mid a$, it remains unrestricted and leaves the second
factor in the statement. These are exactly the asserted formulas,
and every factor is positive. The fresh-prime assertions follow at once.
\end{proof}

\begin{lemma}[Injectivity and separation for logarithmic profiles]
\label{lem:v5-log-separation}
For every nontrivial admissible profile $\mathcal R\ne1$, set
\[
 \delta_{\mathcal R,t}(z)
   =\frac{\mathcal R(z,z^t,t)-1}{U_{\mathcal R}(z)},\qquad t\geq1.
\]
The vector $(\delta_{\mathcal R,t})_{t\geq1}$ is nonzero, and distinct
profiles have distinct vectors. Consequently, a nontrivial finite linear
combination of shifted phases belonging to distinct nontrivial admissible
profiles cannot have a finite limit on an arithmetic progression, once
identical shifts for the same profile have been grouped.
\end{lemma}
\begin{proof}
Direct summation gives
\[
 \sum_{t\geq1}\delta_{\mathcal R,t}(z)z^t
       =\frac{U_{\mathcal R}(z)-1}{(1-z)U_{\mathcal R}(z)}.
\]
Hence equality of two response vectors implies equality of their $U$'s,
and then equality of $\mathcal R(z,z^t,t)$ for every positive integer $t$.
If a polynomial $H(z,w,u)\in\Q(z)[w,u]$ vanishes at
$(w,u)=(z^t,t)$ for every such $t$, then $H=0$.
Indeed, specialize $z$ to a sufficiently small positive number avoiding
all coefficient poles and retaining at least one nonzero coefficient.
Writing $H(z,w,u)=\sum_j P_j(u)w^j$, choose the least $j$ with
$P_j\ne0$. Division by $z^{jt}$ would force the nonzero polynomial
$P_j(t)$ to equal an exponentially decaying polynomial-exponential sum
as $t\to\infty$, which is impossible. This also proves that the zero
response vector corresponds only to $\mathcal R=1$.

Suppose now that the shifted combination has a finite limit. Among its
nonzero grouped coefficients choose the largest value $\kappa_*$, and a
shift $r_*$ carrying a profile of that order. On any fixed refinement of
the starting progression, divide the progression average by
$(\log X)^{\kappa_*-1}$ and use
Lemma~\ref{lem:v5-log-profile-mean}. Terms of smaller order disappear.
When $\kappa_*>1$ the resulting limit is zero; when $\kappa_*=1$ it is the
same finite constant on every refinement.

Choose fresh primes exceeding the number of shifts and all their nonzero
absolute differences. At each such prime an off-class avoids all the shifts modulo
the prime, while an on-class gives exact valuation $t\geq1$ to $N+r_*$.
All the other shifted arguments are units in both classes.
For each nonempty finite set of selected primes take the alternating sum
over their on/off choices. The constant limit and all other shifts cancel.
If $B_{\mathcal R}$ denotes the grouped coefficient at $r_*$ multiplied by
$L_{\mathcal R}(Q,A+r_*)$, the resulting equations are
\[
 \sum_{\kappa_{\mathcal R}=\kappa_*}
 B_{\mathcal R}\prod_{\ell\in S}
            \delta_{\mathcal R,t_\ell}(1/\ell)=0.
\]
The $B_{\mathcal R}$ do not depend on the chosen finite subset $S$.
Use one prime and valuation giving a nonzero response for a selected
profile. For each competing profile use a new prime and valuation that
distinguishes its response from the selected one. Such choices exist by
vector injectivity and the fact that a nonzero rational function has only
finitely many zeros. The product-interpolation argument of
Lemma~\ref{lem:multimean}, now allowing a chosen valuation at each prime,
forces the selected $B_{\mathcal R}$ to vanish. Positivity of its
progression mean contradicts its choice. This proves the claim.
\end{proof}

\begin{theorem}[A unified arithmetic hierarchy]
\label{thm:v5-log-profile-main}
Choose one positive integral realization $f_{\mathcal R,h}$ for each
of any collection of distinct nontrivial admissible valuation profiles.
For every fixed $q\geq1$ and $x=\pm2/b$, $b\geq1$, this entire canonical
base collection satisfies Theorem~\ref{thm:two-sided}.
Thus all its residue-class polynomial moments and shifted-factorial values
in~\eqref{eq:two-sided-family} are jointly linearly independent over $\Q$
together with $1$ and the exponential residue values.
All derivative orders, all normalized antiderivative orders, the four real
directions, and the simple-negative-integer-pole weights in the succeeding
corollaries apply simultaneously to this enlarged family.
\end{theorem}
\begin{proof}
Integral realizations are positive integer-valued and multiplicative.
Lemma~\ref{lem:v5-log-convolution} gives
$f_{\mathcal R,h}(n)=O(n^{h+1/2})$, and
Lemma~\ref{lem:v5-log-separation} proves the required separation across the
whole canonical base collection, including different logarithmic orders.
These are exactly the hypotheses of Theorem~\ref{thm:two-sided}.
\end{proof}

Choosing the least possible integral realization degree is convenient but
not necessary. If two constructions have the same profile and different
degrees, they differ by an integral power of $n$ and must first be combined
in the polynomial-weight coordinates. The profile $1$ is also reduced to
the exponential-polynomial block, rather than counted as a new base.
No list below is intended to override this canonical reduction.

\subsection{Ordered factorizations, divisor powers, and multiplicative mixtures}

The number $d_r(n)$ of ordered $r$-tuples of positive integers with product
$n$ satisfies
\[
 d_r(p^a)=\binom{a+r-1}{r-1}.
\]
This follows by distributing the exponent $a$ among the $r$ ordered
factors. Thus every $d_r$, $r\geq2$, is an admissible base of degree zero,
with profile $\binom{u+r-1}{r-1}$ and logarithmic parameter $\kappa=r$.
The collection may be enlarged at once by
\[
 \tau(n)^m\quad(m\geq1),\qquad
 c^{\omega(n)}\quad(c\geq2),\qquad
 \prod_{j=1}^s d_{r_j}(n)^{m_j},
\]
whose respective profiles and parameters are
\[
 (u+1)^m,\quad \kappa=2^m;
 \qquad c,\quad\kappa=c;
 \qquad \prod_{j=1}^s\binom{u+r_j-1}{r_j-1}^{m_j},
       \quad\kappa=\prod_{j=1}^s r_j^{m_j}.
\]
The products are pointwise products, whereas $d_r$ itself was defined by
Dirichlet convolution. Their arithmetic meanings and their different
operations should not be conflated. For instance, $d_4$ and $\tau^2$
both have $\kappa=4$, but their profiles differ; the response-vector proof
establishes their joint independence rather than relying only on different
logarithmic growth orders.

Any finite pointwise product of admissible integral realizations is again
admissible: multiply the profiles, add the realization degrees, and multiply
the integers $\kappa$. Positivity, integrality, and the polynomial bound in
$u$ are preserved. Consequently the above families can be multiplied by
arbitrary finite products of the earlier divisor, Jordan, Dedekind,
unitary, radical, lattice-counting, and cyclic-subgroup bases.
After identical profiles are grouped, their full two-sided hierarchies
remain jointly independent.

Here is one concrete block at $x=2$. Write
$\mathscr A_f(z)=\sum_{n\geq1}f(n)z^n/n!$. The four bases
\[
 f_1=\tau^2,\qquad f_2=d_4,\qquad
 f_3=4^{\omega},\qquad f_4=\sigma_3\tau^2
\]
have profiles
\[
 (u+1)^2,\qquad \binom{u+3}{3},\qquad
 4,\qquad (u+1)^2\frac{1-z^3w^3}{1-z^3},
\]
which are distinct, even though all have $\kappa=4$.
Therefore the family containing
\begin{gather*}
 1,e^2,e^{-2},\cos2,\sin2,\\
 \mathscr A_{f_j}^{[v]}(2),\quad
 \mathscr A_{f_j}^{[v]}(-2),\quad
 \operatorname{Re}\mathscr A_{f_j}^{[v]}(2i),\quad
 \operatorname{Im}\mathscr A_{f_j}^{[v]}(2i)
 \qquad(1\leq j\leq4,\ v\in\Z)
\end{gather*}
is jointly linearly independent over $\Q$.
Every finite choice is covered by a single instance of the proved theorem.
These are unconditional consequences of the manuscript's mechanism;
no separate assertion of historical priority is made for individual
classical specializations.

\section{Closure under divisor convolutions and a larger arithmetic catalogue}
\label{v5:profile-universe}

The valuation-polynomial theorem supplies a structural extension, not just
another finite list of numerators.  The class is closed under the three
standard ways of combining multiplicative counting functions: pointwise
multiplication, Dirichlet convolution, and unitary convolution.  Repeated
local poles, which produce polynomials in the prime-power exponent, are
precisely why the additional valuation variable is needed.

\subsection{Three closure operations}

In this section an admissible valuation profile has the meaning of
Theorem~\ref{thm:v5-log-profile-main}.  We write its data as $(R,h,\kappa)$,
where
\[
 f(p^a)=p^{ha}R(1/p,p^{-a},a),\qquad
 R(z,z,1)=\kappa+O(z),\qquad \kappa\geq1.
\]
Here $h\geq0$ is an integer, $R\in\Q(z)[w,u]$, and the sampled local
values are positive, polynomially bounded in $a$, and give integers after
multiplication by $p^{ha}$.  The normalization is always $f(1)=1$.
Identical normalized profiles are grouped, regardless of the degree used
to realize them.

For multiplicative functions normalized at $1$, define
\[
 (f*g)(n)=\sum_{d\mid n}f(d)g(n/d),\qquad
 (f\mathbin{\star_u}g)(n)=
       \sum_{\substack{d\mid n\\(d,n/d)=1}}f(d)g(n/d).
\]
The second operation is unitary convolution.  It assigns a whole prime
power to one of its two factors.

\begin{proposition}[Arithmetic closure]\label{v5:closure}
Let $f_1,\ldots,f_s$ have admissible valuation profiles
$(R_i,h_i,\kappa_i)$.
\begin{enumerate}
\item The pointwise product $\prod_i f_i$ has profile $\prod_i R_i$,
degree $\sum_i h_i$, and parameter $\prod_i\kappa_i$.
\item Their Dirichlet convolution $f_1*\cdots*f_s$ has an admissible
valuation profile of degree $H=\max_i h_i$ and parameter
\[
 \kappa=\sum_{i:h_i=H}\kappa_i.
\]
\item Their iterated unitary convolution has degree $H$ and the same
parameter as in part 2.  Its profile is explicitly
\[
 R(z,w,u)=\sum_{i=1}^s w^{H-h_i}R_i(z,w,u).
\]
\end{enumerate}
These operations therefore preserve every factorial-series conclusion of
Theorem~\ref{thm:v5-log-profile-main}, after the usual grouping and removal of
the constant profile $1$.
\end{proposition}
\begin{proof}
Multiplicativity, positive integer values, and the asserted first-prime
parameters are immediate for all three operations.  For part 2 put
\[
 G_i(z,T)=1+\sum_{a\geq1}R_i(z,z^a,a)T^a.
\]
Expanding $R_i$ into monomials $w^j u^m$ expresses $G_i$ as a rational
function of $z,T$ whose possible $T$-poles are $T=z^{-j}$, $j\geq0$.
Indeed $\sum_{a\geq1}a^m(z^jT)^a$ has only the pole
$1-z^jT=0$.  Moreover $G_i$ is bounded at $T=\infty$, so its polynomial
part in $T$ is at most constant.  The normalized local generating function
of the convolution is exactly
\begin{equation}\label{v5:convolution-generating}
 G(z,T)=\prod_{i=1}^sG_i(z,z^{H-h_i}T).
\end{equation}
Its poles have the form $T=z^{-d}$ with $d\geq0$.  Partial fractions in
$T$ show that, for $a\geq1$, its coefficient is a finite sum
\[
 \sum_{d,m}c_{d,m}(z)\binom{a+m-1}{m-1}z^{da}.
\]
This is $R(z,z^a,a)$ for a polynomial $R\in\Q(z)[w,u]$.
The additional denominators introduced by partial fractions are powers
of $z$ and differences $z^d-z^e$ with $d\ne e$, so none introduces a
pole at $z=1/p$.

For a uniform bound, suppose
$R_i(1/p,p^{-a},a)\leq C_i(1+a)^{D_i}$, also allowing $a=0$
by increasing $C_i$.  A coefficient of~\eqref{v5:convolution-generating}
is a sum over $a_1+\cdots+a_s=a$.  The additional factors
$p^{-(H-h_i)a_i}$ are at most one.  Thus it is bounded by
\[
 \left(\prod_i C_i\right)(1+a)^{D_1+\cdots+D_s+s-1}
\]
up to a constant depending only on $s$.  The prime coefficient is
\[
 \sum_i z^{H-h_i}R_i(z,z,1)
       =\sum_{i:h_i=H}\kappa_i+O(z).
\]
This proves part 2.  The displayed profile in part 3 follows by assigning
$p^a$ to exactly one factor; its bounds follow at once.  Part 1 follows
by multiplying the local data.
\end{proof}

\subsection{All weighted ordered factorizations}

Let $\mathbf m=(m_0,m_1,\ldots)$ have finite support, with $m_0\geq1$
and $m_j\geq0$.  Put $H=\max\{j:m_j>0\}$ and define
\begin{equation}\label{v5:weighted-factorizations}
 D_{\mathbf m}=\mathop{*}_{j=0}^{H}
               (\operatorname{id}^{H-j})^{*m_j},
 \qquad \operatorname{id}^{b}(n)=n^b.
\end{equation}
A zeroth convolution power is omitted.  Thus $D_{\mathbf m}(n)$ is a
weighted count of ordered factorizations of $n$; the list of factor weights
contains $m_j$ copies of $H-j$.  Replacing every weight by that weight plus
$c\geq0$ multiplies the answer by $n^c$, which is already covered by the
polynomial-moment coordinates.

The degree is $H$, the parameter is $m_0$, and its profile is uniquely
determined by
\begin{equation}\label{v5:weighted-factorization-local}
 1+\sum_{a\geq1}R_{\mathbf m}(z,z^a,a)T^a
       =\prod_{j\geq0}(1-z^jT)^{-m_j}.
\end{equation}
This formula is also an exact algorithm for computing the profile.
It includes arbitrary repeated weights, rather than just the distinct
weights occurring in the subgroup-counting functions.

Define $U_{\mathbf m}$ by replacing every Dirichlet convolution in
\eqref{v5:weighted-factorizations} with unitary convolution.  It counts
only ordered factorizations whose factors are pairwise coprime.  Its
normalized profile is particularly simple:
\begin{equation}\label{v5:unitary-factorization-profile}
 R^u_{\mathbf m}(z,w,u)=\sum_{j\geq0}m_jw^j.
\end{equation}
In particular, weights all equal to zero give
\[
 D_{(r)}(n)=d_r(n),\qquad U_{(r)}(n)=r^{\omega(n)}.
\]
Here $d_r$ counts ordered factorizations into $r$ positive factors,
and $\omega(n)$ is the number of distinct prime divisors.  Both are
normalized to take the value $1$ at $n=1$.

\begin{corollary}[A nonredundant factorization hierarchy]
\label{v5:factorization-hierarchy}
Exclude the single vector $(1,0,0,\ldots)$.  The profiles of all
$D_{\mathbf m}$ and all $U_{\mathbf m}$ are pairwise distinct, both
within and between the two families.  Consequently the family
\[
 \{1,e^x\}\ \cup\
 \left\{\sum_{n\geq1}\frac{n^dD_{\mathbf m}(n)x^n}{n!},\
        \sum_{n\geq1}\frac{n^dU_{\mathbf m}(n)x^n}{n!}:
                d\geq0,\ \mathbf m\ne(1,0,\ldots)\right\}
\]
is linearly independent over $\Q$ at every permitted $x$.
Every residue-class, two-sided jet, and simple-negative-pole extension
in Theorem~\ref{thm:v5-log-profile-main} applies to the same canonical list.
\end{corollary}
\begin{proof}
At the first prime power, both profiles give
$\sum_jm_jz^j$.  Thus equality of profiles forces equality of the
vectors $\mathbf m$.  For a fixed vector, write the positive numbers
$z^j$, repeated $m_j$ times, as $b_1,\ldots,b_\ell$.  At the second
prime power the ordinary and unitary values differ by
\[
 \sum_{1\leq i<j\leq\ell}b_ib_j>0
\]
when $\ell\geq2$.  The excluded vector is the only case with
$\ell=1$ under the convention $m_0\geq1$ and no redundant overall
power of $n$.  Thus all retained profiles are distinct and none is $1$.
The conclusion follows from the valuation-profile theorem.
\end{proof}

For example, if $h,m\geq1$, let
\[
 B_{h,m}(n)=\sum_{d\mid n}d^h d_m(n/d).
\]
This is the subfamily with $m_0=1$, $m_h=m$, and all other entries zero.
The case $m=1$ is $\sigma_h$, whereas $m\geq2$ produces genuine
valuation dependence.  For instance
\begin{equation}\label{v5:B-h-two-profile}
 R_{B_{h,2}}(z,w,u)=
 \frac{1-z^hw^h\bigl((u+2)-(u+1)z^h\bigr)}{(1-z^h)^2}.
\end{equation}
Its coefficient of $u w^h$ is nonzero.  Polynomial-exponential
independence therefore shows that it cannot be represented by a profile
in $\Q(z)[w]$ alone.  The new theorem really enlarges the earlier class.

The mean constants are explicit without becoming new irrationality
claims about zeta values.  For the ordinary hierarchy, the normalized
unrestricted local mean is
\[
 U_{R_{\mathbf m}}(z)
     =(1-z)^{1-m_0}\prod_{j\geq1}(1-z^{j+1})^{-m_j}.
\]
Hence Lemma~\ref{lem:v5-log-profile-mean} gives
\begin{equation}\label{v5:factorization-mean}
 \frac1X\sum_{n\leq X}\frac{D_{\mathbf m}(n)}{n^H}
 \sim\frac{\prod_{j\geq1}\zeta(j+1)^{m_j}}{(m_0-1)!}
                (\log X)^{m_0-1}.
\end{equation}
In particular $B_{h,m}(n)/n^h$ has mean $\zeta(h+1)^m$.
Equation~\eqref{v5:factorization-mean} is a mean-value formula; the
irrational numbers asserted here remain the factorial series.

\subsection{All positive gcd moments, including the diagonal case}

For $r,k\geq1$ put
\begin{equation}\label{v5:gcd-moment-definition}
 G_{r,k}(n)=\sum_{1\leq a_1,\ldots,a_r\leq n}
              \gcd(a_1,\ldots,a_r,n)^k.
\end{equation}
Grouping by the common gcd gives
\begin{equation}\label{v5:gcd-moment-convolution}
 G_{r,k}=\operatorname{id}^{k}*J_r
         =\mu*\operatorname{id}^{k}*\operatorname{id}^{r}.
\end{equation}
The number of tuples with common gcd $d\mid n$ is $J_r(n/d)$:
divide by $d$ and exclude, for each prime dividing $n/d$, the tuples
all of whose coordinates are divisible by that prime.  This proves
\eqref{v5:gcd-moment-convolution} directly.  In particular
$G_{r,k}=G_{k,r}$, so only one order of the pair is retained.

Set $H=\max(r,k)$ and $\Delta=|r-k|$.  The normalized local generating
function is
\begin{equation}\label{v5:gcd-moment-local}
 1+\sum_{a\geq1}\frac{G_{r,k}(p^a)}{p^{Ha}}T^a
       =\frac{1-z^HT}{(1-T)(1-z^\Delta T)},\qquad z=1/p.
\end{equation}
For $\Delta>0$ this gives the bounded profile
\[
 R_{r,k}(z,w,u)=
        1+(1-z^{\min(r,k)})
              \frac{z^\Delta(1-w^\Delta)}{1-z^\Delta}.
\]
For $r=k$ it gives instead
\begin{equation}\label{v5:diagonal-gcd-profile}
 R_{r,r}(z,w,u)=1+u(1-z^r),\qquad \kappa=2.
\end{equation}
This diagonal was absent from the bounded-profile theorem.  Its first
case is the classical gcd-sum, or Pillai, function
$G_{1,1}(n)=\sum_{a=1}^n\gcd(a,n)$.

The profiles for $1\leq k\leq r$ are all distinct: their first-prime
values are $1+z^{r-k}-z^r$ when $k<r$, and $2-z^r$ when $k=r$.
They are also distinct from every ordinary and unitary factorization
profile above, whose first-prime polynomials have nonnegative
coefficients.  Therefore all the $G_{r,k}$, $k\leq r$, can be adjoined
\emph{simultaneously} to Corollary~\ref{v5:factorization-hierarchy}, with
all the same moments, residue restrictions, and two-sided jets.

For reference, their normalized means are
\[
 \frac1X\sum_{n\leq X}\frac{G_{r,k}(n)}{n^H}
 \sim
 \begin{cases}
  \zeta(\Delta+1)/\zeta(H+1),&r\ne k,\\
  (\log X)/\zeta(r+1),&r=k.
 \end{cases}
\]
They follow either from~\eqref{v5:gcd-moment-local} or from the general
mean theorem, with no transfer of irrationality to these zeta ratios.

\subsection{Gcd moments on the unit group}

A second natural family sums over units rather than over all residue
classes.  For $r,k\geq1$ define
\begin{equation}\label{v5:unit-gcd-definition}
 M_{r,k}(n)=
 \sum_{a_1,\ldots,a_r\in(\Z/n\Z)^\times}
       \gcd(a_1-1,\ldots,a_r-1,n)^k.
\end{equation}
Use the convention that the unit group modulo $1$ has one element,
so $M_{r,k}(1)=1$.  The basic identity
$g^k=\sum_{d\mid g}J_k(d)$ and the surjectivity of reduction of units
modulo $n$ onto units modulo $d\mid n$ give
\begin{equation}\label{v5:unit-gcd-divisor}
 M_{r,k}(n)=\varphi(n)^r
                 \sum_{d\mid n}\frac{J_k(d)}{\varphi(d)^r}.
\end{equation}
Indeed, for each $d\mid n$ there are
$(\varphi(n)/\varphi(d))^r$ unit tuples with every coordinate congruent
to $1$ modulo $d$.  Reduction is surjective prime power by prime power,
and then by the Chinese remainder theorem.  This proves the identity,
its multiplicativity, and its positive integrality.  These are classical
Menon-type counting identities; see~\cite{TothMenon23}.  At $r=k=1$
one recovers $M_{1,1}=\varphi\tau$.

The prime-power computation from~\eqref{v5:unit-gcd-divisor} is explicit.
With degree $H=\max(r,k)$, the normalized profiles are
\begin{equation}\label{v5:unit-gcd-profiles}
 R^M_{r,k}(z,w,u)=
 \begin{cases}
 (1-z)^r w^{k-r}+
       (1-z^k)\dfrac{1-w^{k-r}}{1-z^{k-r}},&k>r,\\[6pt]
 (1-z)^r+u(1-z^r),&k=r,\\[6pt]
 (1-z)^r+
       (1-z^k)\dfrac{z^{r-k}(1-w^{r-k})}{1-z^{r-k}},&k<r.
 \end{cases}
\end{equation}
They are bounded in the off-diagonal cases and polynomially bounded
in the diagonal case.  Positivity follows directly from the count.
The parameter is $1$ off the diagonal and $2$ on the diagonal.  None
is the constant profile $1$: in the off-diagonal cases the coefficient
of the displayed nonconstant power of $w$ is nonzero, and on the
diagonal the coefficient of $u$ is nonzero.
Consequently every factorial series formed from $M_{r,k}$ is irrational,
and every finite product or convolution of these functions with the
previous catalogue is covered by Proposition~\ref{v5:closure} whenever
its resulting profile is not $1$.  Joint statements use the same exact
profile-grouping rule; they do not presume that every written counting
formula is a different coordinate.

\subsection{All and cyclic subgroups of rank-two finite groups}

Let $C_m=\Z/m\Z$.  For $1\leq r\leq s$, let $S_{r,s}(n)$ be the
number of all subgroups of $C_{n^r}\times C_{n^s}$, and let $C_{r,s}(n)$
be the number of its cyclic subgroups.  These familiar counting functions
have explicit formulas; see~\cite{Hampejs14}.  The following direct
argument supplies all arithmetic inputs needed here.

\begin{proposition}\label{v5:rank-two-counts}
Both functions are multiplicative and have admissible valuation profiles
of degree $r$ and parameter $\kappa=s-r+1$.  Writing $b=s-r$, the profiles
are
\begin{align}
 R^S_{r,s}(z,w,u)
 &= (bu+1)\frac{1-zw^r}{1-z}
   +\frac{2z\bigl(1-(ru+1)w^r+ruzw^r\bigr)}{(1-z)^2},
       \label{v5:all-subgroups-profile}\\
 R^C_{r,s}(z,w,u)
 &=bu+1+\frac{2z(1-w^r)}{1-z}.
       \label{v5:cyclic-subgroups-profile}
\end{align}
All profiles in the combined $S$ and $C$ list are pairwise distinct
and nonconstant.  Thus all their factorial constants and polynomial
moments are jointly linearly independent with $1,e^x$, and the full
residue-class and two-sided hierarchy applies.
\end{proposition}
\begin{proof}
For groups of coprime orders every subgroup splits into its primary
components, and cyclicity does also.  This proves multiplicativity.
To count subgroups of $C_{p^A}\times C_{p^B}$, choose the projection
onto the first factor, of order $p^i$, and the intersection with the
second factor.  If the quotient of the second factor by that
intersection has order $p^j$, the remaining choice is a homomorphism
$C_{p^i}\longrightarrow C_{p^j}$.  There are $p^{\min(i,j)}$ such
homomorphisms.  Consequently, for $A=ra$, $B=sa$,
\[
 S_{r,s}(p^a)=\sum_{i=0}^{ra}\sum_{j=0}^{sa}p^{\min(i,j)}
   =p^{ra}\sum_{d=0}^{ra}\bigl((s-r)a+2d+1\bigr)p^{-d}.
\]
At $d=0$ there are $(s-r)a+1$ pairs having minimum $ra$; at the
other depths successive rectangular counts give the coefficient
$(s-r)a+2d+1$.  Summing the finite geometric series and its derivative
gives~\eqref{v5:all-subgroups-profile}.

For cyclic subgroups, count elements of each exact order and divide by
the number of generators of a cyclic group of that order.  For orders
$p^j$ with $1\leq j\leq ra$ this gives $(p+1)p^{j-1}$; for
$ra<j\leq sa$ it gives $p^{ra}$.  Including the trivial subgroup yields
\[
 C_{r,s}(p^a)=\bigl((s-r)a+1\bigr)p^{ra}
                   +2\sum_{j=0}^{ra-1}p^j,
\]
which proves~\eqref{v5:cyclic-subgroups-profile}.
Positivity and integrality follow from the counts; the two displayed
profiles are uniformly $O(1+a)$ at the sampled arguments.

At the first prime power their normalized values are, respectively,
\[
 \kappa+\sum_{d=1}^{r}(\kappa+2d)z^d,
 \qquad \kappa+2\sum_{d=1}^{r}z^d.
\]
The polynomial degree and constant term identify $r,s$ within either
list.  Between the lists the coefficient of $z$ differs, since
$\kappa+2\geq3$ in the first list and it is $2$ in the second.
Every profile is nonconstant.  Apply
Theorem~\ref{thm:v5-log-profile-main}.
\end{proof}

In particular the series counting all subgroups of
$C_n\times C_n$, $C_n\times C_{n^2}$, and $C_{n^2}\times C_{n^5}$,
and the corresponding cyclic-subgroup series, are covered at once.
Their arbitrary fixed pointwise products and divisor convolutions
remain within Proposition~\ref{v5:closure}; cross-family lists are
again grouped by their exact profiles before independence is asserted.

\subsection{A compact list of concrete consequences}

Write $\mathcal E_f(x)=\sum_{n\geq1}f(n)x^n/n!$.
At every permitted $x$, each of the following is irrational:
\[
 \mathcal E_{d_r}(x)\ (r\geq2),\qquad
 \mathcal E_{r^{\omega}}(x)\ (r\geq2),\qquad
 \mathcal E_{\operatorname{id}^{h}*d_m}(x)\ (h,m\geq1),
\]
\[
 \mathcal E_{G_{r,k}}(x)\ (r,k\geq1),\qquad
 \mathcal E_{M_{r,k}}(x)\ (r,k\geq1).
\]
For a particularly concrete simultaneous statement, the numbers
\[
 1,\ e^2,\ \mathcal E_{d_3}(2),\ \mathcal E_{3^{\omega}}(2),\
 \mathcal E_{\operatorname{id}*d_2}(2),\
 \mathcal E_{G_{1,1}}(2),\ \mathcal E_{G_{2,1}}(2),\
 \mathcal E_{\varphi\tau}(2)
\]
are linearly independent over $\Q$.
The first-prime profile values are respectively
$3,3,1+2z,2-z,1+z-z^2,2-2z$; the sole collision, between $d_3$
and $3^{\omega}$, is separated at the second prime power by $6\ne3$.
All degree-zero, residue-class, primitive-integral, and higher-derivative
coordinates permitted by the general theorem can be added without
changing this reasoning, after canonical grouping.

The countably infinite statements in this section mean that every
finite subfamily is independent.  They concern values of factorial
series and their specified linear transforms.  They neither assert
algebraic independence of these values nor identify new historically
named open problems for each illustrative constant.

\section{Finite-ring matrices, Grassmannians, and flag counts}
\label{sec:v5-finite-geometry}

The profile class also includes counting functions from finite linear
algebra.  The counting identities below are elementary and classical;
the application here is their insertion into the common factorial-value
theorem.  Put $R_n=\Z/n\Z$, with the usual convention that each counting
function has value one at $n=1$.

\begin{proposition}[Three finite-geometric counting profiles]
\label{prop:v5-finite-counts}
For $m\ge r\ge1$, let $S_{m,r}(n)$ count surjective $R_n$-linear maps
$R_n^m\longrightarrow R_n^r$.  For $m>r\ge1$, let $G_{m,r}(n)$ count
free direct summands of rank $r$ in $R_n^m$.  If
$\mathbf b=(b_1,\ldots,b_s)$ has positive integer entries, $s\ge2$,
and $m=\sum b_i$, let $F_{\mathbf b}(n)$ count flags of free direct
summands whose successive quotient ranks are $b_1,\ldots,b_s$.
All three functions are positive integral and multiplicative.  Their
degree $h$ and normalized prime-power profile are
\[
\begin{array}{c|c|c}
\text{function}&h&R(z)\\ \hline
S_{m,r}&mr&\displaystyle\prod_{j=m-r+1}^{m}(1-z^j)\\[4pt]
G_{m,r}&r(m-r)&\displaystyle
 \frac{\prod_{j=m-r+1}^{m}(1-z^j)}{\prod_{j=1}^{r}(1-z^j)}\\[6pt]
F_{\mathbf b}&\displaystyle\sum_{i<j}b_i b_j&\displaystyle
 \frac{\prod_{j=1}^{m}(1-z^j)}
      {\prod_{i=1}^{s}\prod_{j=1}^{b_i}(1-z^j)}.
\end{array}
\]
Here $f(p^a)=p^{ha}R(1/p)$ for every prime $p$ and $a\ge1$.
\end{proposition}
\begin{proof}
Over $\Z/p^a\Z$, a matrix defines a surjection if and only if its
reduction modulo $p$ has full row rank.  There are
$\prod_{j=0}^{r-1}(p^m-p^j)$ such matrices over $\mathbb F_p$,
and each has $p^{(a-1)mr}$ lifts.  This proves the first formula.
The criterion follows directly from a unit $r\times r$ minor; conversely
a surjection remains surjective after reduction.

A full-column-rank $m\times r$ matrix over this local ring is an ordered
basis of a free direct summand.  Its column basis can be extended to a
basis of the ambient module, again using a unit minor.  Each summand
has exactly $|\mathrm{GL}_r(\Z/p^a\Z)|=S_{r,r}(p^a)$ ordered bases.
Thus $G_{m,r}(p^a)=S_{m,r}(p^a)/S_{r,r}(p^a)$, giving the second formula.

The general linear group acts transitively on the stated flags, since
bases of the successive free quotients can be lifted and concatenated.
The stabilizer of a standard flag consists of block upper triangular
matrices with invertible diagonal blocks of sizes $b_i$.  If
$D=\sum_{i<j}b_i b_j$, its order is
$p^{aD}\prod_i S_{b_i,b_i}(p^a)$.  Dividing $S_{m,m}(p^a)$ by this
order gives the last formula, since $m^2-\sum b_i^2-D=D$.
The Chinese remainder decomposition of $R_n$ proves multiplicativity
and the asserted product formulas for arbitrary $n$.
\end{proof}

Each displayed rational function has no pole on $0<z\le1/2$, is positive
there, is bounded there, and has value $1+O(z)$ at zero.  It is consequently
an admissible profile from Section~\ref{sec:profiles}.  It is nonconstant
for the specified nontrivial ranks.  For $S_{m,r}$ this follows from the
product formula.  For $G_{m,r}$ the profile is the Gaussian binomial
polynomial $\left[\begin{smallmatrix}m\\r\end{smallmatrix}\right]_z$,
of positive degree $r(m-r)$, with constant and leading coefficients one.
The flag profile is a Gaussian multinomial polynomial of positive degree
$D$.  These polynomial descriptions also follow inductively by decomposing
subspaces over a finite field according to their intersection with a fixed
hyperplane; the rational product formulas above are enough for admissibility.

Two especially familiar cases are
\[
 |\mathrm{GL}_m(R_n)|
   =n^{m^2}\prod_{p\mid n}\prod_{j=1}^{m}(1-p^{-j}),
 \qquad
 |\mathrm{SL}_m(R_n)|
   =n^{m^2-1}\prod_{p\mid n}\prod_{j=2}^{m}(1-p^{-j})
 \quad(m\ge2).
\]
The second identity uses the surjective determinant map onto $R_n^\times$.
For complete flags, $\mathbf b=(1,\ldots,1)$, the profile is
$\prod_{j=1}^{m}(1+z+\cdots+z^{j-1})$.

\begin{corollary}[Factorial constants from finite linear algebra]
\label{cor:v5-geometric-irrationality}
For every nontrivial counting function $f$ above, every permitted
$x=\pm2/b$, and every integer $j$, the number $A_f^{[j]}(x)$ is irrational.
The same holds after any fixed residue restriction and with every weight
$n^d$ ($d\ge0$) or $1/(n+a)$ ($a\ge1$).  All these values obey the
joint-independence and exact rationality rules of the canonical profile
theorem.  In particular a single fixed $f$ supplies the entire independent
two-sided hierarchy.
\end{corollary}
\begin{proof}
The preceding proposition verifies all the arithmetic hypotheses.
Choose $f$ itself as the integral realization of its normalized profile
and apply Theorem~\ref{thm:two-sided} and
Corollary~\ref{cor:simple-poles}.  For a collection of counting functions,
choose one integral realization per distinct profile.  Any other
realization of the same profile differs by an integral power of $n$;
after choosing the least degree in a finite collection, this is a
nonnegative polynomial weight.  Shifted-factorial values reduce to
polynomial weights and simple negative-integer poles by partial fractions.
Their coefficient vectors, rather than their different counting names,
therefore determine all joint relations.
\end{proof}

For example, the following are irrational at $x=2$:
\[
 \sum_{n\ge1}\frac{2^n|\mathrm{GL}_3(R_n)|}{n!},\qquad
 \sum_{n\ge1}\frac{2^nG_{6,2}(n)}{n!(n+7)},\qquad
 \int_0^1 t^4(1-t)^3 A_{F_{(2,1,3)}}(2t)\,dt.
\]
There are genuine equalities among these counting functions, which the
profile criterion preserves.  For example
$G_{2,1}(n)=\Psi_1(n)$ and
$|\mathrm{SL}_2(R_n)|=nJ_2(n)$.
Also $G_{m,r}=G_{m,m-r}$, and permutations of the block sizes have the
same flag count.  These are identical arithmetic coefficients, not extra
independent constants.

One further matrix family comes for free.  Let $M_{u,v,r}(n)$ count
$u\times v$ matrices whose image is a free rank-$r$ direct summand,
where $1\le r\le\min(u,v)$ and $G_{r,r}:=1$.
Choosing the image and then a surjection onto it gives
\[
 M_{u,v,r}(n)=G_{u,r}(n)S_{v,r}(n).
\]
Hence its degree is $r(u+v-r)$ and its profile is the product of the
two profiles above.  It is nontrivial, positive, and admissible, so all
the preceding irrationality statements apply.  Finally, the closure theorems of Section~\ref{v5:profile-universe} permit any finite pointwise, ordinary Dirichlet, or unitary
convolution of these counts with the other admitted arithmetic families,
subject to the same explicit grouping of equal normalized profiles.


\section{Total prime-factor multiplicity on rough residue classes}
\label{sec:v5-rough-omega}

Polynomial dependence on a prime exponent includes $c^{\omega(n)}$,
but not $c^{\Omega(n)}$, where $\Omega$ counts prime factors with
multiplicity.  An additional growth observation admits the latter
on every residue class avoiding the finitely many primes at most $c$.
This extension uses an exponential, rather than polynomial, local
dependence on the valuation.

\begin{lemma}[Strictly sublinear normalized growth suffices]
\label{lem:v5-sublinear-growth}
In Theorem~\ref{thm:two-sided}, its growth hypothesis can be replaced
by the following: for each base $f$ there is a real number
$0\leq\theta_f<1$ such that
\[
 f(n)=O_f(n^{h_f+\theta_f}).
\]
The other hypotheses, in particular multiplicativity, positive integer
values, and isolated-shift separation, are unchanged.  All its residue,
derivative, antiderivative, and rational-weight consequences remain valid.
There is no requirement that one common $\theta<1$ work for an infinite
base family.
\end{lemma}
\begin{proof}
In a finite proposed relation choose
$\theta=\max_f\theta_f<1$.  Keep the same notation $r,w,J,T$ and
the same grid as in the proof of Theorem~\ref{thm:two-sided}.
The tail expansion~\eqref{eq:offset-tail} now has remainder
\[
 O(n^{h+\theta-J-r-2})=O(n^{w+\theta-J-2}).
\]
Indeed the first omitted factorial term has denominator order
$J+r+2$, and the same bound holds for the finite denominator expansions.
The surviving terminal terms are
$O(N^{w_*+\theta-J-1})$.  Since $J\geq w_*+1$, their total is
$O(N^{\theta-2})$, and the separately cancelled exponential remainder
is $O(N^{-T})$.  The common dyadic denominator is still $O(N)$.
Multiplication by it therefore gives an integer of size
$O(N^{\theta-1})+O(N^{1-T})$, tending to zero.

After the weighted tail becomes identically zero, multiply by
$N^{J+1-w_*}$ as before.  The expansion error is
$O(N^{\theta-1})$; each smaller effective degree contributes at most
$O(N^{\theta-1})$; and replacement of a fixed shifted power ratio
by $1$ has the same bound.  The exponential remainder tends to zero
by the original choice $T=\max(2,J+2-w_*)$.
The identical isolated-shift contradiction follows.  All subsequent
changes of basis are exact and hence remain applicable.
\end{proof}

For an integer $c\geq2$ define
\[
 M_c=\prod_{p\leq c}p,\qquad
 n_{>c}=\prod_{\substack{p^a\Vert n\\p>c}}p^a,\qquad
 W_c(n)=c^{\Omega(n_{>c})}.
\]
Thus $W_c$ is a positive integer-valued completely multiplicative
function, with prime value $1$ at $p\leq c$ and $c$ at $p>c$.
On $(n,M_c)=1$ it agrees with the ordinary function $c^{\Omega(n)}$.
Let $P_c$ be the least prime exceeding $c$ and put
\[
 \theta_c=\frac{\log c}{\log P_c}<1.
\]
Prime by prime, $c\leq p^{\theta_c}$ for $p>c$, so
\begin{equation}\label{eq:v5-omega-growth}
 1\leq W_c(n)\leq n^{\theta_c}.
\end{equation}

\begin{lemma}[Logarithmic means for the rough multiplicity function]
\label{lem:v5-rough-mean}
The function $W_c$, normalized with $h=0$, has positive progression
means of logarithmic order $c-1$ as in
Lemma~\ref{lem:v5-log-profile-mean}.  Its unrestricted leading
coefficient is $G_c/(c-1)!$, where
\[
 G_c=
 \prod_{p\leq c}(1-p^{-1})^{c-1}
 \prod_{p>c}\frac{(1-p^{-1})^c}{1-c/p}>0.
\]
At a fresh prime $p>c$, the unit and exact-valuation-$t$ refinement
factors are, respectively,
\[
 \frac1{U_c(1/p)},\qquad\frac{c^t}{U_c(1/p)},
 \qquad U_c(z)=\frac{1-z}{1-cz}.
\]
In particular the response vector is
\begin{equation}\label{eq:v5-rough-response}
 \delta_{c,t}(z)=(c^t-1)\frac{1-cz}{1-z},\qquad t\geq1.
\end{equation}
\end{lemma}

\begin{proof}
Set $g_c=W_c*\mu^{*c}$.  Its prime-power generating series is
\[
 \sum_{a\geq0}g_c(p^a)T^a=
 \begin{cases}
 (1-T)^{c-1},&p\leq c,\\
 (1-T)^c/(1-cT),&p>c.
 \end{cases}
\]
At a good prime $p>c$ the linear coefficient vanishes, and the
higher coefficients are $O_c(c^a)$.  Choose any
\[
 \max(1/2,\theta_c)<\sigma<1.
\]
Then $cP_c^{-\sigma}<1$ and
\[
 \sum_{p>c}\sum_{a\geq2}|g_c(p^a)|p^{-a\sigma}
 \ll_{c,\sigma}\sum_{p>c}p^{-2\sigma}<\infty.
\]
The finitely many remaining Euler factors are polynomials.  Therefore
$\sum_n|g_c(n)|n^{-\sigma}<\infty$.

This is the exact convolution input used in the proof of
Lemma~\ref{lem:v5-log-profile-mean}, with $d_c$ in place of
$d_\kappa$.  For the nonprincipal character estimate one chooses
the convolution exponent also larger than $1-1/c$; this is possible
because all three lower bounds are strictly less than $1$.
The principal-character calculation gives $G_c$ above, and the
nonprincipal sums have a power-saving error.  The good Euler factors
of $G_c$ are positive and $1+O_c(p^{-2})$, so the product is positive.

For nonunit progressions there is no new local-growth issue:
$W_c(am)=W_c(a)W_c(m)$ holds for all $a,m$, without a coprimality
restriction.  Dividing a progression by $(a,Q)$ therefore reduces it
to the unit-class calculation with the same function, multiplied by
a positive constant.  The local conditional factors are those in
Lemma~\ref{lem:v5-log-profile-mean}.  At $p>c$ its local series is
$1/(1-cz)$, which gives $U_c$ and~\eqref{eq:v5-rough-response}.
All errors are $O_{c,Q,a}((\log X)^{c-2})$ after division by $X$.
\end{proof}

\begin{theorem}[Adjoining every rough total-multiplicity channel]
\label{thm:v5-rough-omega}
Adjoin $W_c$, with growth degree $h=0$, for every integer $c\geq2$
to the canonical family of nontrivial admissible valuation profiles
in Theorem~\ref{thm:v5-log-profile-main}.  The enlarged family satisfies
the two-sided hierarchy and all its residue-class conclusions.
It also satisfies the cyclotomic-coordinate and quadratic-field
conclusions of Section~\ref{sec:cyclotomic-v5}.
\end{theorem}

\begin{proof}
The strict growth condition is~\eqref{eq:v5-omega-growth}, and the
progression means were just proved.  It remains to verify joint
isolated-shift separation; individual irrationality alone would not
suffice.

Use the largest logarithmic order in a finite proposed limit, exactly
as in Lemma~\ref{lem:v5-log-separation}.  Two different $W_c$ have
different orders.  At a fixed order $c-1$, compare $W_c$ with a
polynomial valuation profile $\mathcal R$ of $\kappa_{\mathcal R}=c$.
Their response vectors cannot agree.  Otherwise, for every $t\geq1$,
\[
 \mathcal R(z,z^t,t)
  =1+\frac{U_{\mathcal R}(z)}{U_c(z)}(c^t-1).
\]
Fix a sufficiently small real $z\in(0,1/c)$ avoiding all poles and
zeros of the nonzero rational prefactor.  The left side has at most
polynomial growth in $t$, because $\mathcal R$ is polynomial in
its last two variables and $z^t\to0$.  The right side has nonzero
exponential growth.  This is a contradiction.

Thus the response vectors within every logarithmic order are distinct
and nonzero.  All their fixed-$t$ components are rational functions
of $z$.  The same finite fresh-prime interpolation as in
Lemma~\ref{lem:v5-log-separation} separates the surviving coefficients.
Primes at most the largest retained $c$ are simply excluded from
the finite set of allowed fresh primes.  This proves isolated-shift
separation for the enlarged family.
Apply Lemma~\ref{lem:v5-sublinear-growth} and then the exact changes
of basis already established.
\end{proof}

\subsection{Products with every polynomial valuation profile}

The rough multiplicity construction can be combined with all the
earlier arithmetic profiles, not only placed beside them.  Set
$W_1(n)=1$ and $M_1=1$.  For $c\geq1$ and an admissible valuation
profile $\mathcal R$, choose one integral realization
$f_{\mathcal R,h}=n^h\rho_{\mathcal R}$ and put
\[
 F_{c,\mathcal R}(n)=W_c(n)f_{\mathcal R,h}(n),\qquad
 \rho_{c,\mathcal R}(n)=W_c(n)\rho_{\mathcal R}(n).
\]
Permit $\mathcal R=1$ here with $h=0$, but exclude the single pair
$(c,\mathcal R)=(1,1)$.  Profiles equal as elements of
$\Q(z)[w,u]$ are grouped before selecting the realization.

\begin{theorem}[Rough exponential times polynomial profiles]
\label{thm:v5-rough-product}
The entire canonical family $F_{c,\mathcal R}$, over all the pairs
just specified, satisfies the two-sided hierarchy, its residue-class
and rational-weight consequences, and the cyclotomic and quadratic-field
consequences.  Distinct pairs give separated normalized phases.
\end{theorem}
\begin{proof}
At a good prime $p>c$ its normalized local values are
$c^a\mathcal R(p^{-1},p^{-a},a)$; at $p\leq c$ the factor $c^a$
is omitted.  The integral realization remains a positive integer
and is multiplicative.  The previous bounds give
$\rho_{c,\mathcal R}(n)=O(n^{\theta_c+\epsilon})$ for every
$\epsilon>0$ (take $\theta_1=0$).  Choose
$0<\epsilon<1-\theta_c$ and use
Lemma~\ref{lem:v5-sublinear-growth}.

The logarithmic order is $\kappa-1$, where
$\kappa=c\kappa_{\mathcal R}$.  Indeed
$g=\rho_{c,\mathcal R}*\mu^{*\kappa}$ has $g(p)=O(p^{-1})$ at
good primes and $g(p^a)=O(c^a(1+a)^D)$ at higher powers.
Its Dirichlet series converges absolutely at every
$\sigma>\max(1/2,\theta_c)$; the finitely many bad local factors
have only polynomial coefficient growth.  The character-convolution
proof of Lemma~\ref{lem:v5-log-profile-mean} therefore applies.
In the nonunit reduction, a modification at a fixed good prime has
values bounded by a constant times $c^a(1+a)^D$, so its local series
still converges at the same $\sigma<1$; at a bad prime it has polynomial
growth.  Thus that part of the proof also applies.

At good primes the rational local functions are
\[
 L_{c,\mathcal R}(z)
  =1+\sum_{a\geq1}c^a\mathcal R(z,z^a,a)z^a,
 \qquad U_{c,\mathcal R}(z)=(1-z)L_{c,\mathcal R}(z).
\]
Rationality follows term by term from
$\sum_{a\geq1}a^m(cz^{j+1})^a$.  The leading Euler product is
strictly positive: its local factors are
$(1-p^{-1})^\kappa L_{c,\mathcal R}(1/p)$ at good primes,
and $(1-p^{-1})^\kappa\mathcal L_{\mathcal R}(1/p)$ at bad primes.
Every factor is positive and the good factors are $1+O(p^{-2})$.
The fresh-prime response components are
\[
 \delta_{c,\mathcal R,t}(z)=
 \frac{c^t\mathcal R(z,z^t,t)-1}{U_{c,\mathcal R}(z)}.
\]

These vectors are injective in the pair $(c,\mathcal R)$.
The identity
\[
 \sum_{t\geq1}\delta_{c,\mathcal R,t}(z)z^t
  =\frac{U_{c,\mathcal R}(z)-1}
         {(1-z)U_{c,\mathcal R}(z)}
\]
first recovers $U$.  Equality of two vectors consequently gives
\[
 c^t\mathcal R(z,z^t,t)
   ={c'}^t\mathcal R'(z,z^t,t)\qquad(t\geq1).
\]
Write $\mathcal R(z,w,u)=\sum_j r_j(z,u)w^j$ and likewise for
$\mathcal R'$.  Specialize $z$ to a sufficiently small positive real
number avoiding the finitely many coefficient poles and zeros needed
to retain a nonzero coefficient.  If $c\ne c'$, also avoid the finitely
many equations $cz^j=c'z^{j'}$.  The displayed equality is then an
identity of polynomial-exponential sequences with distinct positive
bases $cz^j$ and $c'z^{j'}$.  Such sequences are independent: divide
by the largest base having a nonzero polynomial coefficient; the
coefficient would tend to zero exponentially, impossible for a nonzero
polynomial.  Removing that term inductively proves the assertion.
Thus $c=c'$, and the same argument applied to $\mathcal R-\mathcal R'$
gives $\mathcal R=\mathcal R'$.  Comparison with $(1,1)$ also shows
that only the excluded pair has zero response vector.

At the largest retained logarithmic order the fresh-prime interpolation
of Lemma~\ref{lem:v5-log-separation} now applies to these distinct
rational response vectors.  All primes at most the largest retained
$c$ are avoided.  This proves the required isolated-shift separation
and completes the proof.
\end{proof}

In particular, on the unit classes modulo $M_c$ the numerator is the
unmodified product $c^{\Omega(n)}f_{\mathcal R,h}(n)$.
Thus one may simultaneously include such products with powers of
$\tau$, ordered-factorization functions, divisor sums, Jordan functions,
and the other proved polynomial valuation profiles.  Products that give
the same polynomial profile must first be grouped.  The construction
does not license arbitrary products of the $W_c$ without rechecking
growth: for example, the local value of $W_2W_3$ at $5^a$ is $6^a$,
which exceeds the strictly sublinear regime at that prime.

\begin{corollary}[Unmodified total-multiplicity sums on rough classes]
\label{cor:v5-ordinary-omega}
Define the entire function
\[
 B_c(z)=\sum_{\substack{n\geq1\\(n,M_c)=1}}
              \frac{c^{\Omega(n)}z^n}{n!},\qquad c\geq2.
\]
At every permitted positive $x$, all values $B_c^{[j]}(x)$,
$c\geq2$, $j\in\Z$, are jointly linearly independent with $1,e^x$
and all the arithmetic two-sided jets from the polynomial valuation
families.  The same holds for individual unit residue classes modulo
$M_c$, after retaining independent residue coordinates.  Every one
of these values is irrational.
\end{corollary}

\begin{proof}
On the indicated residue classes $W_c(n)=c^{\Omega(n)}$ exactly.
Thus $B_c$ is a nonempty sum of residue-restricted $A_{W_c}$ functions.
For any finite proposed relation take a common multiple of the finitely
many $M_c$.  Each coefficient of $B_c^{[j]}$ is then copied to all
its unit residue coordinates.  The supports for different base-and-order
pairs are disjoint, so a nonzero coefficient cannot vanish in this
conversion.  The preceding theorem gives the contradiction.
One must not list both a sum of residue coordinates and all of its
summands as an independent family; the assertion uses either independent
individual coordinates or the displayed sums.
\end{proof}

For example, each of
\[
 \sum_{\substack{n\geq1\\n\ \mathrm{odd}}}
       \frac{2^{\Omega(n)+n}}{n!},\qquad
 \sum_{\substack{n\geq1\\(n,6)=1}}
       \frac{3^{\Omega(n)}2^n}{n!},\qquad
 \sum_{\substack{n\geq1\\(n,6)=1}}
       \frac{4^{\Omega(n)}2^n}{n!}
\]
is irrational, and the three are jointly independent with $1,e^2$
and all the polynomial valuation profile jets at $2$.
Polynomial moments, arbitrary normalized integration orders, and
the permitted rational weights may be added in their canonical bases.
The removal of the small primes is essential to the proved growth
bound: this section does not assert the analogous unrestricted theorem
for $\sum c^{\Omega(n)}x^n/n!$.

\section{Every sublinear integer coefficient sequence}
\label{sec:v5-sublinear}

There is a complementary principle which requires no multiplicativity,
positivity, or progression means.  It supplies an exact rationality test
for every integer coefficient sequence of sublinear growth.  Its proof
uses the same dyadic denominator lemma, but only at the particularly
favourable indices $n=2^k$.  The principle and its consequences are
presented as elementary extensions of the current toolkit, without a
claim of historical priority.

\begin{theorem}[An exact sublinear-coefficient criterion]
\label{thm:v5-sublinear}
Let $a_n\in\Z$ for $n\geq0$, and suppose $a_n=o(n)$ as $n\to\infty$.
For every $x=\varepsilon2/b$, $b\geq1$ an integer and
$\varepsilon\in\{-1,1\}$,
\[
 S_a(x)=\sum_{n\geq0}\frac{a_nx^n}{n!}\in\Q
 \quad\Longleftrightarrow\quad
 a_n=0\text{ for all sufficiently large }n.
\]
The same assertion holds for rational coefficients with a common fixed
integer denominator.  In particular every non-eventually-zero integer
sequence satisfying $a_n=O(n^\beta)$ for some $\beta<1$ gives an
irrational value.
\end{theorem}
\begin{proof}
Finite support plainly gives a rational value.  Conversely suppose
$S_a(x)=u/v\in\Q$, $v\geq1$, and set
\[
 T(n)=n!x^{-n}\sum_{m>n}\frac{a_mx^m}{m!}.
\]
The growth assumption implies polynomial growth, so
Lemma~\ref{lem:dyadic-height} applies.  At $n=2^k$ it gives
$2vT(2^k)\in\Z$.

Write $y=|x|$.  For every $\eta>0$, the inequality
$|a_m|\leq\eta m$ holds beyond a fixed threshold.  For all sufficiently
large $n$,
\begin{align*}
 |T(n)|
 &\leq\eta\sum_{j\geq1}(n+j)y^j\frac{n!}{(n+j)!}\\
 &=\eta y\sum_{k\geq0}y^k\frac{n!}{(n+k)!}
 \leq\frac{\eta y}{1-y/(n+1)}.
\end{align*}
Since $\eta$ is arbitrary, $T(n)\to0$.  Thus $T(2^k)=0$ for all
sufficiently large $k$.

It remains to exclude cancellation of an infinite signed tail.
Choose $\eta>0$ so small that $2\eta y<1$, and take a large index
$n$ with $T(n)=0$.  If there were a first $m>n$ with $a_m\ne0$,
then, after factoring out $x^m/m!$, its tail would equal
\[
 a_m+\sum_{j\geq1}a_{m+j}x^j\frac{m!}{(m+j)!}.
\]
The second term has absolute value at most
$\eta y/(1-y/(m+1))<1$ for large $m$, whereas the nonzero integer
$a_m$ has absolute value at least one.  The tail could not vanish.
Therefore no such $m$ exists, proving eventual zero.  Clearing the
fixed denominator proves the rational-coefficient version.
\end{proof}

The growth threshold cannot simply be replaced by $O(n)$.  For the
same $x=\varepsilon2/b$, the non-eventually-zero integer sequence
$a_n=bn-2\varepsilon$ satisfies
\[
 \sum_{n\geq0}\frac{(bn-2\varepsilon)x^n}{n!}
 =(bx-2\varepsilon)e^x=0.
\]
This example identifies a genuine boundary of the universal statement;
the arithmetic-profile theorems overcome it only through additional
structure of their coefficients.

\begin{corollary}[Sequence relations classify all value relations]
\label{cor:v5-sequence-classification}
Let $a^{(i)}_n$ be any collection of integer sequences, each $o(n)$,
and put $S_i(x)=\sum_{n\geq0}a^{(i)}_nx^n/n!$.
For finitely supported rational coefficients $c_i$,
\[
 \sum_i c_iS_i(x)\in\Q
 \quad\Longleftrightarrow\quad
 \sum_i c_i a^{(i)}_n=0
 \text{ for every sufficiently large }n.
\]
Consequently $\{1\}\cup\{S_i(x)\}$ is linearly independent over
$\Q$ exactly when the coefficient sequences are linearly independent
over $\Q$ modulo eventual equality.  When the displayed sum is rational,
its value is the corresponding finite initial factorial sum.
\end{corollary}
\begin{proof}
Every finite rational combination has a common denominator and remains
$o(n)$.  Apply Theorem~\ref{thm:v5-sublinear}.
\end{proof}

The criterion applies to derivatives and normalized antiderivatives
without multiplying the coefficients by powers of their indices.
Extend a sequence $a_n$, $n\geq0$, by $a_n=0$ for $n<0$, and set
$A(z)=\sum_{n\geq0}a_nz^n/n!$.  For every integer $j$,
\[
 A^{[j]}(x)=\sum_{n\geq0}\frac{a_{n+j}x^n}{n!},
\]
where negative orders have the zero-origin normalization used earlier.
Each shifted sequence is still $o(n)$.  Thus a rational relation among
any finite number of these two-sided values is equivalent to an
eventual constant-coefficient relation among the corresponding shifts
of $a_n$.  This is an exact criterion; it does not assume that all shifts
of every sequence are independent.

\begin{theorem}[The full hierarchy of every bounded aperiodic sequence]
\label{thm:v5-aperiodic-hierarchy}
Let $(a_n)_{n\geq0}$ be a bounded integer sequence which is not eventually
periodic, and put $A(z)=\sum_{n\geq0}a_nz^n/n!$.
For every fixed $q\geq1$ and every permitted $x=\pm2/b$, the family
\[
 \{1\}\cup\{E_{q,r}(x):0\leq r<q\}
 \cup\{A^{[j]}(x):j\in\Z\}
\]
is linearly independent over $\Q$.
In particular all derivatives and all normalized antiderivatives of
$A$, together with $1,e^x$, are jointly independent.
\end{theorem}
\begin{proof}
By Corollary~\ref{cor:v5-sequence-classification}, a rational relation
would give, for every sufficiently large $n$,
\[
 w(n)+\sum_{j\in J}c_j a_{n+j}=0,
\]
where $w$ is rational-valued and $q$-periodic, and $J$ is finite.
If every $c_j$ is zero then $w$ is eventually zero, hence identically
zero, so all its residue coefficients vanish.

Otherwise translate the index and retain the largest shift to obtain
an affine recurrence
\[
 a_{n+d}=-\frac{w'(n)+\sum_{i=0}^{d-1}c'_i a_{n+i}}{c'_d},
 \qquad c'_d\ne0,
\]
with a periodic forcing $w'$.  If $d=0$, this says directly that $a_n$
is eventually periodic.  If $d\geq1$, the state
$(n\bmod q,a_n,\ldots,a_{n+d-1})$ belongs to a finite set, because
the coefficients $a_n$ are bounded integers.  The recurrence determines
its next state uniquely.  Two equal states therefore have identical
subsequent trajectories; an infinite trajectory in this finite set
is eventually periodic.  Again $a_n$ is eventually periodic, a
contradiction.  All coefficients in the proposed relation must vanish.
\end{proof}

This theorem supplies complete two-sided hierarchies for several
familiar coefficient sequences:
\begin{enumerate}
\item The prime indicator $a_n=\mathbf1_{n\text{ prime}}$.
There are infinitely many ones, while factorial blocks give arbitrarily
long intervals of zeros at arbitrarily large indices.  Such a sequence
cannot be eventually periodic.
\item The M\"obius function $a_n=\mu(n)$ and its square
$a_n=\mu(n)^2$, with value zero assigned at $n=0$.
For every desired block length choose distinct primes $p_i$ and impose
$N+i\equiv0\pmod{p_i^2}$ by CRT.  This gives arbitrarily late zero
blocks, while primes supply infinitely many nonzero terms.
\item The Liouville function $a_n=(-1)^{\Omega(n)}$ for $n\geq1$,
again assigning any integer at zero.  If it had eventual period $q$,
its values on all sufficiently large multiples $qn$ would be constant.
Complete multiplicativity would then make $a_n$ eventually constant.
The identity $a_{2n}=-a_n$ contradicts this.
\item For each fixed integer $k\geq2$, the indicator of perfect
$k$th powers; or the indicator of the set $\{2^m:m\geq0\}$.
Each has infinitely many ones and gaps of zeros tending to infinity.
\end{enumerate}
For example the prime factorial function
\[
 \mathcal P(z)=\sum_{p\text{ prime}}\frac{z^p}{p!}
\]
has $\{1,e^2,\mathcal P^{[j]}(2):j\in\Z\}$ linearly independent
over $\Q$.  The same assertion holds separately for each other
function just listed.  These separate families must not be united
without checking eventual relations between their coefficient
sequences.  Nor does the theorem automatically adjoin such arbitrary
sequences to the positive arithmetic-profile family: their combined
coefficients need not be sublinear.

\subsection{A simultaneous signed multiplicative family}

The sequence criterion also allows several of the preceding arithmetic
functions to be combined, provided their eventual relations are checked.
For $c\geq1$ define
\[
 B_c(n)=c^{\omega(n)},\quad
 L_c(n)=\lambda(n)c^{\omega(n)},\quad
 S_c(n)=\mu(n)^2c^{\omega(n)},\quad
 M_c(n)=\mu(n)c^{\omega(n)}.
\]
Write $W_c(n)=c^{\Omega(n_{>c})}$ for $c\geq2$, where $n_{>c}$ is
the product of all prime-power factors of $n$ with prime greater than
$c$.  The growth bound for this function is also proved in the
rough-prime extension above.  For any coefficient function $f$ here,
write $A_f(z)=\sum_{n\geq1}f(n)z^n/n!$.

\begin{theorem}[A joint signed arithmetic block]
\label{thm:v5-signed-block}
At every permitted $x=\pm2/b$, the family
\[
\begin{gathered}
 1,\ e^x,\ \mathcal P(x),\\
 A_{d_r}(x)\ (r\geq2),\qquad
 A_{B_c}(x)\ (c\geq2),\qquad A_{W_c}(x)\ (c\geq2),\\
 A_{L_c}(x),\ A_{S_c}(x),\ A_{M_c}(x)\qquad(c\geq1)
\end{gathered}
\]
is linearly independent over $\Q$.  Here $\lambda(n)=(-1)^{\Omega(n)}$
and $\mathcal P$ is the prime factorial function defined above.
\end{theorem}
\begin{proof}
All the coefficient sequences are $o(n)$.  For $d_r$ and $c^{\omega}$
this follows from the polynomial-valuation growth bound; multiplying by
$\mu$, $\mu^2$, or $\lambda$ does not increase absolute values.
If $p_c$ is the least prime greater than $c$, then
$W_c(n)\leq n^{\log c/\log p_c}=o(n)$.
By Corollary~\ref{cor:v5-sequence-classification} it suffices to exclude
a finite nontrivial eventual relation between these coefficient sequences,
including the constant sequence $B_1=1$ for $e^x$.

Choose a prime $p$ larger than every $c$ occurring in the proposed
relation, and evaluate at $n=p^a$ as $a\to\infty$.
The $W_c$ terms are $c^a$, the $d_r$ terms are the polynomials
$\binom{a+r-1}{r-1}$, the $B_c$ terms are constants $c$, and the
$L_c$ terms are $c(-1)^a$.  The prime, $S_c$, and $M_c$ terms vanish
for $a\geq2$.  Successively dividing by the largest exponential $c^a$
shows that all the $W_c$ coefficients vanish.  The remaining polynomial
plus an alternating constant can vanish for all large $a$ only if its
polynomial is constant and its alternating coefficient is zero.
The distinct positive degrees of the $d_r$ polynomials then force all
their coefficients to vanish.

Now choose $k$ distinct primes.  Give them all exponent two, and then
change one exponent to three.  The numbers can be chosen arbitrarily
large.  In both cases $\omega(n)=k$, while $\lambda(n)$ has opposite
signs and all $S_c,M_c$, and prime-indicator terms are zero.
Adding and subtracting the two eventual relations gives
\[
 \sum_c b_c c^k=0,\qquad \sum_c \ell_c c^k=0\qquad(k\geq1).
\]
Finite Vandermonde matrices on the distinct positive bases $c$ are
invertible, so all $B_c$ and $L_c$ coefficients vanish, including $b_1$.

For a product of $k\geq2$ distinct primes, the remaining relation is
\[
 \sum_c s_c c^k+(-1)^k\sum_c m_c c^k=0.
\]
Even $k$ and a Vandermonde matrix on the bases $c^2$ give $s_c+m_c=0$;
odd $k\geq3$ give $s_c-m_c=0$.  Thus all $S_c$ and $M_c$ coefficients
vanish.  Evaluation at arbitrarily large primes finally kills the prime
indicator coefficient.  This proves the asserted eventual independence
and hence the value independence.
\end{proof}

In particular the constants
\[
 \sum_{n\geq1}\frac{2^n\mu(n)}{n!},\quad
 \sum_{n\geq1}\frac{2^n\lambda(n)}{n!},\quad
 \sum_{n\geq1}\frac{2^{n+\omega(n)}\mu(n)^2}{n!},\quad
 \sum_{p\text{ prime}}\frac{2^p}{p!}
\]
belong to one explicit joint independent family together with all the
$d_r$, $B_c$, and $W_c$ values in the theorem.  This is a base-value
statement.  Combining the entire derivative hierarchies of different
bounded sequences would require additional eventual-shift independence,
which has not been assumed or asserted here.

The theorem proves arithmetic properties of the displayed factorial
values.  It does not prove digit aperiodicity of an unrelated named
constant, and it does not imply transcendence of these values.

\section{All cyclotomic orders and two quadratic coefficient fields}
\label{sec:cyclotomic-v5}

Fix a positive permitted radius $x=2/b$, with an integer $b\geq1$.  The residue-class
theorem has an infinite angular consequence: one can simultaneously use
every cyclotomic order, provided that the angular values are expressed in
their canonical rational coordinates.  At orders four and six this also
gives linear independence of the actual complex values over a quadratic
number field.  These statements do not require a new denominator estimate.

\subsection{Canonical coordinates at every root of unity}

Let $\Phi_d(T)$ be the $d$th cyclotomic polynomial, let
$\zeta_d=\exp(2\pi i/d)$, and put $\varphi_d=\deg\Phi_d$.
For an entire function $H$ with real Taylor coefficients define the
polynomial of degree less than $\varphi_d$
\begin{equation}\label{eq:cyclotomic-interpolant-v5}
 \mathcal C_d[H](T)=
 \sum_{\substack{1\leq k\leq d\\(k,d)=1}}
 H(\zeta_d^k x)
 \frac{\Phi_d(T)}{(T-\zeta_d^k)\Phi_d'(\zeta_d^k)}
 =\sum_{\ell=0}^{\varphi_d-1}C_{d,\ell}[H](x)T^\ell.
\end{equation}
For $d=1$, this means $C_{1,0}[H](x)=H(x)$.
Conjugating the interpolation formula pairs $k$ with $-k$, so every
$C_{d,\ell}[H](x)$ is real.  These are explicitly defined finite linear
combinations of actual angular values, with algebraic coefficients.

\begin{theorem}[All cyclotomic coordinates simultaneously]
\label{thm:all-cyclotomic-v5}
For the canonical base family $\mathcal B$ of
Theorem~\ref{thm:two-sided}, the real numbers
\[
 \{1\}\ \cup\
 \{C_{d,\ell}[e^z](x):d\geq1,\ 0\leq\ell<\varphi_d\}
 \ \cup\
 \{C_{d,\ell}[A_f^{[j]}](x):
       f\in\mathcal B,\ j\in\Z,\ d\geq1,\
       0\leq\ell<\varphi_d\}
\]
are linearly independent over $\Q$.
In particular, every displayed coordinate is irrational, and every
nonzero finite rational combination of the coordinates is irrational.
\end{theorem}

\begin{proof}
Consider a finite proposed relation and take a positive integer $q$
divisible by every cyclotomic order that occurs.  For each $f,j$ put
$U_a=A_{f,q,a}^{[j]}(x)$, $0\leq a<q$, and define
\[
 P_{f,j}(T)=\sum_{a=0}^{q-1}U_a T^{\langle a-j\rangle_q},
\]
where $\langle\cdot\rangle_q$ denotes the representative in
$\{0,\ldots,q-1\}$.  The derivative and antiderivative series give
\[
 P_{f,j}(\zeta_d^k)=A_f^{[j]}(\zeta_d^k x)
 \qquad(d\mid q).
\]
Indeed a derivative changes the exponent of $z$ from $n$ to $n-j$;
the same exponent formula applies to a normalized antiderivative.
Thus the remainder of $P_{f,j}$ modulo $\Phi_d$ is exactly the
interpolant in~\eqref{eq:cyclotomic-interpolant-v5}.

The map
\[
 \Q[T]/(T^q-1)\longrightarrow
 \bigoplus_{d\mid q}\Q[T]/(\Phi_d),
 \qquad P\longmapsto(P\bmod\Phi_d)_{d\mid q},
\]
is an isomorphism over $\Q$.  To see this without a dimension
assumption, the cyclotomic factors are pairwise coprime, so a polynomial
in the kernel is divisible by their product $T^q-1$; the dimensions
on both sides are $q$, since $\sum_{d\mid q}\varphi_d=q$.
Taking remainders therefore gives an invertible rational matrix on
the $q$ coefficient coordinates.  The initial permutation
$a\mapsto\langle a-j\rangle_q$ is invertible as well.

For the exponential use $P_0(T)=\sum_a E_{q,a}(x)T^a$ and the same
remainder map.  Hence, block by block, all coordinates of orders
dividing $q$ are an invertible rational transform of the family in
Corollary~\ref{cor:two-sided-jet}.  That family, together with $1$,
is linearly independent.  Its transformed family is independent,
and so is the subfamily in the proposed relation.  The order $q$
can be chosen for every finite relation; this proves the assertion
for all $d$ simultaneously.
\end{proof}

The theorem is not a claim that arbitrary algebraic-coefficient
combinations of the original real constants are independent.  The
interpolants are special: reduction modulo $\Phi_d$ identifies their
coefficients as a rational, invertible transform of residue-class
coordinates.  Nor does a decomposition
$H(\zeta_d x)=\sum_\ell C_{d,\ell}[H](x)\zeta_d^\ell$, by itself,
prove transcendence of $H(\zeta_d x)$.

\subsection{Six directions and twelve explicit real transforms}

Write $\omega=\exp(i\pi/3)$.  For any $H$ as above, the six real
transforms
\begin{equation}\label{eq:six-transforms-v5}
 H(x),\quad H(-x),\quad
 \Re H(\omega x),\quad\frac{\Im H(\omega x)}{\sqrt3},\quad
 \Re H(\omega^2x),\quad\frac{\Im H(\omega^2x)}{\sqrt3}
\end{equation}
are an invertible rational transform of the cyclotomic coordinates
of orders $1,2,3,6$.  For example, at order six the two coordinates
satisfy
\[
 H(\omega x)=C_{6,0}+\omega C_{6,1},\qquad
 \Re H(\omega x)=C_{6,0}+\tfrac12 C_{6,1},\qquad
 \frac{\Im H(\omega x)}{\sqrt3}=\tfrac12 C_{6,1}.
\]
At order three the sign of $C_{3,1}/2$ in the real part is negative.
Consequently all six transforms of $e^z$, and all six transforms
of every $A_f^{[j]}$, are jointly independent together with $1$.
For the exponential these are
\[
 e^x,\ e^{-x},\quad
 e^{x/2}\cos\frac{\sqrt3x}{2},\quad
 \frac{e^{x/2}}{\sqrt3}\sin\frac{\sqrt3x}{2},\quad
 e^{-x/2}\cos\frac{\sqrt3x}{2},\quad
 \frac{e^{-x/2}}{\sqrt3}\sin\frac{\sqrt3x}{2}.
\]

One can retain these six transforms, add the order-four pair
$\Re H(ix),\Im H(ix)$, and add four further transforms.  Put
$\eta=\exp(i\pi/6)$ and define
\begin{align*}
 D_0[H]&=\frac{\Re H(\eta x)+\Re H(\eta^5x)}2,&
 D_1[H]&=\frac{\Re H(\eta x)-\Re H(\eta^5x)}{\sqrt3},\\
 D_2[H]&=\frac{\Im H(\eta x)-\Im H(\eta^5x)}{\sqrt3},&
 D_3[H]&=\frac{\Im H(\eta x)+\Im H(\eta^5x)}2.
\end{align*}
Since $\Phi_{12}(T)=T^4-T^2+1$, writing its interpolant as
$C_0+C_1T+C_2T^2+C_3T^3$ gives
\[
 D_0=C_0+\tfrac12 C_2,\qquad D_1=C_1,\qquad
 D_2=C_2,\qquad D_3=\tfrac12 C_1+C_3.
\]
This rational transformation is invertible.  Therefore the twelve
transforms just listed, for every arithmetic function and every
integer derivative order, are jointly independent with the twelve
exponential transforms and $1$.
The last four exponential transforms have the familiar forms
\begin{gather*}
 \cosh\frac{\sqrt3x}{2}\cos\frac x2,\qquad
 \frac2{\sqrt3}\sinh\frac{\sqrt3x}{2}\cos\frac x2,\\
 \frac2{\sqrt3}\sinh\frac{\sqrt3x}{2}\sin\frac x2,\qquad
 \cosh\frac{\sqrt3x}{2}\sin\frac x2.
\end{gather*}
The factors $1/\sqrt3$ are part of the stated coordinates and must
not be silently removed when asserting joint independence.
The general theorem supplies higher cyclotomic orders by the same
explicit interpolation procedure, without a fixed upper bound on $d$.

\subsection{Gaussian and Eisenstein linear independence}

\begin{lemma}[Imaginary quadratic scalar extension]
\label{lem:imaginary-quadratic-v5}
If a family of real numbers is linearly independent over $\Q$, it
is linearly independent over every imaginary quadratic field
$K=\Q(\sqrt{-D})$, $D\in\N$.
\end{lemma}
\begin{proof}
Write each coefficient as $a+i\sqrt D\,b$ with $a,b\in\Q$.
The real and imaginary parts of a finite relation are two rational
relations among the same real numbers.  Both are trivial.
\end{proof}

\begin{corollary}[Actual angular values over quadratic fields]
\label{cor:quadratic-angular-v5}
For $q=4$, put $K=\Q(i)$ and $\zeta=i$; for $q=6$, put
$K=\Q(\sqrt{-3})$ and $\zeta=\omega$.
In either case the complex numbers
\[
 \{1\}\ \cup\ \{e^{\zeta^k x}:0\leq k<q\}
 \ \cup\
 \{A_f^{[j]}(\zeta^k x):f\in\mathcal B,\ j\in\Z,
                              0\leq k<q\}
\]
are linearly independent over $K$.
In particular, each displayed arithmetic value lies outside $K$,
and every nonzero finite $K$-linear combination of the exponential
and arithmetic values lies outside $K$.
\end{corollary}

\begin{proof}
The residue-class family in Corollary~\ref{cor:two-sided-jet} is real.
Lemma~\ref{lem:imaginary-quadratic-v5} extends its independence to $K$.
For each fixed $f,j$, the evaluation map is the discrete Fourier map
\[
 A_f^{[j]}(\zeta^k x)=
 \zeta^{-kj}\sum_{a=0}^{q-1}\zeta^{ka}A_{f,q,a}^{[j]}(x).
\]
The Fourier matrix is invertible over $K$, with inverse obtained
by conjugating its entries and dividing by $q$.  Its diagonal
phase factors $\zeta^{-kj}$ are nonzero elements of $K$.
The exponential block has the same Fourier matrix without those
phase factors.  Apply these invertible transformations to each
finite collection of blocks.  Independence together with $1$
also proves the final assertion.
\end{proof}

This quadratic-field statement is stronger than individual
irrationality, but it is still a linear-independence result, not an
algebraic-independence or transcendence theorem.  For a general
cyclotomic field its real subfield is nontrivial, so the scalar-extension
lemma does not apply.  The canonical real-coordinate theorem remains
valid at every order.

Finally, all arguments in this section concern a common absolute
radius $x$.  For a relation combining unrelated rational radii,
rationality of the sum does not imply rationality of each summand;
one cannot apply Lemma~\ref{lem:dyadic-height} separately at each
radius.  No independence across unrelated radii is claimed here.

\section{Logarithmic kernels and positive integrals in every dimension}
\label{sec:v5-kernel-calculus}

The rational-weight theorem has an exact integral calculus.  It includes
logarithmic kernels after the indicated differential factors are supplied,
as well as positive integrals over cubes of arbitrarily large dimension.
The results in this section are consequences of the earlier theorem; different
integral representations are not counted as independent constants unless
their rational multipliers are independent.

Throughout, $\mathcal B$ is the canonical arithmetic family in
Theorem~\ref{thm:two-sided}, $x=\pm2/b$ with $b\geq1$, and
$A_f(z)=\sum_{n\geq1}f(n)z^n/n!$.  Write $\theta=z\,d/dz$.
For integers $a,r\geq1$ and $P\in\Q[X]$, define
\[
 \mathscr L_{a,r,P}[A_f](x)=
 \frac1{(r-1)!}\int_0^1
 t^{a-1}(-\log t)^{r-1}(P(\theta)A_f)(xt)\,dt.
\]

\begin{proposition}[An exact test for logarithmic integral expressions]
\label{prop:v5-log-kernels}
Consider finitely many polynomials $P_{f,a,r}\in\Q[X]$ and put
\[
 R_f(X)=\sum_{a,r}\frac{P_{f,a,r}(X)}{(X+a)^r}.
\]
Suppose each reduced rational function $R_f$ has at most simple poles.
If some $R_f$ is nonzero, then for all $u,v\in\Q$ the number
\[
 u+v e^x+\sum_{f,a,r}\mathscr L_{a,r,P_{f,a,r}}[A_f](x)
\]
is irrational.  The same conclusion holds with separate rational
multipliers on each residue class, adjoining the corresponding
periodic exponential functions.

The simple-pole condition can be checked by finite rational arithmetic:
for every $f,a$ and every $j\geq2$, it is precisely
\begin{equation}
 \sum_{r\geq j}
 \frac{P_{f,a,r}^{(r-j)}(-a)}{(r-j)!}=0.
 \label{eq:v5-principal-parts}
\end{equation}
Only the finitely many indices present in the expression need be tested.
\end{proposition}
\begin{proof}
For $n\geq1$, substitution $t=e^{-u}$ and repeated integration by parts
give
\[
 \frac1{(r-1)!}\int_0^1t^{n+a-1}(-\log t)^{r-1}\,dt
 =\frac1{(n+a)^r}.
\]
Moreover $P(\theta)A_f(z)=\sum_{n\geq1}P(n)f(n)z^n/n!$.
Termwise integration is absolutely justified: the arithmetic coefficients
have polynomial growth, the factorial series converges uniformly on the
compact set of arguments, and its terms vanish at least linearly at zero.
Consequently
\[
 \mathscr L_{a,r,P}[A_f](x)
 =\sum_{n\geq1}\frac{P(n)f(n)x^n}{n!(n+a)^r}.
\]
The entire expression therefore has multiplier $R_f(n)$ in channel $f$.
Apply Corollary~\ref{cor:simple-poles}.  Taylor expansion of each
$P_{f,a,r}$ at $-a$ gives exactly~\eqref{eq:v5-principal-parts} as the
coefficient of $(X+a)^{-j}$, proving the finite test.  Residue restrictions
commute with all these operations and use the residue version of the same
corollary.
\end{proof}

In particular, for every $a,r\geq1$,
\begin{equation}
 \frac1{(r-1)!}\int_0^1t^{a-1}(-\log t)^{r-1}
       ((\theta+a)^{r-1}A_f)(xt)\,dt
 =M_{f,a}(x)\notin\Q.
 \label{eq:v5-log-example}
\end{equation}
This includes arbitrary logarithmic powers, with an explicit differential
operator.  For fixed $f,a$, changing $r$ in~\eqref{eq:v5-log-example}
gives the same number.  Without that operator, the repeated-pole
multiplier $(n+a)^{-r}$ is outside the asserted theorem when $r>1$.
Failure of~\eqref{eq:v5-principal-parts} therefore means that this particular
criterion does not apply; it does not imply rationality.

\begin{theorem}[Product kernels and a jointly independent tower]
\label{thm:v5-product-kernels}
Let $a_1,a_2,\ldots$ be any sequence of pairwise distinct positive integers.
For $m\geq1$ define
\[
 I_{f,m}(x)=\int_{[0,1]^m}
       \left(\prod_{j=1}^m t_j^{a_j-1}\right)
       A_f\!\left(x\prod_{j=1}^m t_j\right)
       \,dt_1\cdots dt_m.
\]
All numbers $I_{f,m}(x)$, for $f\in\mathcal B$ and $m\geq1$, are
jointly linearly independent over $\Q$ together with $1,e^x$ and
all polynomial moments $\sum_{n\geq1}n^d f(n)x^n/n!$.
They may also be restricted jointly to every residue class of a fixed
modulus.  Each is irrational.
\end{theorem}
\begin{proof}
Fubini's theorem and absolute convergence give
\[
 I_{f,m}(x)=\sum_{n\geq1}
   \frac{f(n)x^n}{n!\prod_{j=1}^m(n+a_j)}.
\]
The multiplier $R_m(X)=\prod_{j=1}^m(X+a_j)^{-1}$ has only simple
negative-integer poles.  The family consisting of all $R_m$ and all
polynomials is linearly independent over $\Q$: in a finite relation,
the pole at $-a_M$ occurs only in the term with largest index $M$;
its residue is nonzero, forcing that coefficient to vanish.  Continue
downward, and then use polynomial independence.
Corollary~\ref{cor:simple-poles} now rules out every proposed relation
of the displayed numbers, separately in each canonical arithmetic
channel and residue class.
\end{proof}

For example, taking $a_j=j$ gives the particularly simple identity
\[
 I_{f,m}(x)=\sum_{n\geq1}\frac{f(n)x^n}{(n+m)!}.
\]
More generally,
\[
 I_{f,m}(x)=\sum_{j=1}^m
 \frac{M_{f,a_j}(x)}{\prod_{\ell\ne j}(a_\ell-a_j)}.
\]
These are exact alternative representations, not additional independent
coordinates on top of all the $M_{f,a}$.  Multiplication by
$\prod_j a_j$ gives the expectation of $A_f(x\prod_j U_j)$ when
the $U_j$ are independent and have densities $a_j t^{a_j-1}$ on
$(0,1)$.  Thus the corresponding expectation constants inherit the
same independence statement.

Arbitrary logarithmic powers can also be inserted simultaneously in
this cube integral.  For positive integers $r_1,\ldots,r_m$, replace
the weight by
\[
 \prod_{j=1}^m
 \frac{t_j^{a_j-1}(-\log t_j)^{r_j-1}}{(r_j-1)!}
\]
and replace $A_f$ by
$\prod_{j=1}^m(\theta+a_j)^{r_j-1}A_f$.
The resulting integral is exactly $I_{f,m}(x)$, as is seen term by
term.  This supplies a positive-dimensional integral identity even
when the logarithmic kernel is of arbitrarily high order.

\begin{theorem}[Positive bi-Vandermonde cube integrals]
\label{thm:v5-vandermonde-integrals}
Fix a positive integer $a$.  For $m\geq1$, set
\begin{align*}
 V_m(t)&=\prod_{1\leq i<j\leq m}(t_j-t_i),\\
 W_m(t)&=\prod_{1\leq i<j\leq m}(t_j^m-t_i^m),\\
 K_{f,a,m}(x)&=\int_{[0,1]^m}
  \left(\prod_{j=1}^m t_j^{a-1}\right)V_m(t)W_m(t)
  A_f\!\left(x\prod_{j=1}^m t_j\right)\,dt.
\end{align*}
Write $(y)_d=y(y+1)\cdots(y+d-1)$ and
\[
 C_m=m!\,m^{m(m-1)/2}\prod_{j=1}^{m-1}(j!)^2.
\]
Then
\begin{equation}
 K_{f,a,m}(x)=C_m\sum_{n\geq1}
          \frac{f(n)x^n}{n!(n+a)_{m^2}}.
 \label{eq:v5-vandermonde-value}
\end{equation}
For this fixed $a$, all $K_{f,a,m}(x)$, as $m\geq1$ and
$f\in\mathcal B$ vary, are jointly linearly independent over $\Q$
together with $1,e^x$ and all polynomial moments.  In particular they
are irrational.  The kernel multiplying $A_f$ is nonnegative on the
cube and positive away from equal coordinates and the boundary.
\end{theorem}
\begin{proof}
Expansion of the two Vandermonde determinants and Fubini's theorem
give, for each integer $n\geq0$,
\[
 \int_{[0,1]^m}\left(\prod_jt_j^{n+a-1}\right)V_m(t)W_m(t)\,dt
 =m!\det\left[\frac1{n+a+mi+j}\right]_{0\leq i,j<m}.
\]
Indeed, expand both determinants over permutations; integrating the
product of monomials gives this determinant identity after fixing one
permutation and reindexing the other.  The Cauchy determinant formula
then yields
\[
 m!\det\left[\frac1{n+a+mi+j}\right]
 =\frac{m!\prod_{0\leq i<j<m}m(j-i)^2}
        {\prod_{i,j=0}^{m-1}(n+a+mi+j)}
 =\frac{C_m}{(n+a)_{m^2}}.
\]
Here the integers $mi+j$ run through $0,\ldots,m^2-1$ exactly once.
The Cauchy identity itself follows by multiplying the determinant by
its common denominator: the numerator is alternating separately in
the row and column parameters, its degree gives the two Vandermonde
factors, and a residue at one pole gives constant one by induction.
Termwise integration now proves~\eqref{eq:v5-vandermonde-value}.

For a finite linear relation, the multiplier with largest $m$ has a
simple pole at $-a-m^2+1$ that none of the smaller multipliers has.
Its coefficient therefore vanishes, and descending induction proves
independence of these rational multipliers together with polynomials.
Apply Corollary~\ref{cor:simple-poles}.  Finally,
$(t_j-t_i)(t_j^m-t_i^m)\geq0$ on $[0,1]^2$, with strict inequality
for distinct positive coordinates, proving the kernel assertion.
\end{proof}

At $a=1$, formula~\eqref{eq:v5-vandermonde-value} simplifies to
\[
 K_{f,1,m}(x)=C_m\sum_{n\geq1}\frac{f(n)x^n}{(n+m^2)!}.
\]
For example, with $m=2$ one obtains the explicit double integral
\[
 \int_0^1\!\int_0^1
  (v-u)^2(u+v)A_f(xuv)\,du\,dv
 =4\sum_{n\geq1}\frac{f(n)x^n}{(n+4)!}\notin\Q.
\]
The weight has total mass $C_m/(a)_{m^2}$, a positive rational number.
Dividing by this mass therefore gives another probabilistic
interpretation without changing irrationality.  As throughout this
section, the integral representations inherit the proved hierarchy;
they are not claims to new historical resolutions of named problems.

\section{Polynomial integral operators and exact projective classification}
\label{sec:operator-projective}

This section extracts further consequences of the established independent
family.  The claims are closure and classification results, rather than new
irrationality criteria for unrelated constants.  Throughout, $q\geq1$ is
fixed, $\mathcal B$ is a canonical base family satisfying
Theorem~\ref{thm:two-sided}, and $x=\pm2/b\ne0$ is a permitted rational
point.  Write $A_{f,q,a}^{[j]}$ for the two-sided jets in
Corollary~\ref{cor:two-sided-jet}.

\subsection{A normal form for arbitrarily nested polynomial kernels}

Let $D$ denote differentiation, let
\[
 (JG)(z)=\int_0^zG(t)\,dt,
\]
and allow multiplication by any polynomial in $\Q[z]$.  Start with the
functions $1$, $E_{q,a}$, and $A_{f,q,a}$, and apply finitely many of these
operations, rational linear combinations, and compositions.  In particular,
this permits repeated polynomial-kernel integrals
\[
 G(z)\longmapsto\int_0^z K(z,t)G(t)\,dt,
 \qquad K\in\Q[z,t],
\]
interspersed with differentiation.

\begin{theorem}[Polynomial-operator normal form]
\label{thm:operator-normal-form}
Every resulting function has a unique finite representation
\begin{equation}\label{eq:operator-normal-form}
 P_{\mathrm c}(z)+\sum_{a=0}^{q-1}P_a(z)E_{q,a}(z)
   +\sum_{f,a,j}P_{f,a,j}(z)A_{f,q,a}^{[j]}(z),
 \qquad P_{\mathrm c},P_a,P_{f,a,j}\in\Q[z].
\end{equation}
At a permitted point $x$, its value is rational if and only if
\[
 P_a(x)=0\quad\hbox{for every }a,
 \qquad P_{f,a,j}(x)=0\quad\hbox{for every }f,a,j.
\]
In that case the value is exactly $P_{\mathrm c}(x)$.
\end{theorem}
\begin{proof}
The displayed module is closed under polynomial multiplication and $D$.
Indeed, $DE_{q,a}=E_{q,a-1}$, with residues read modulo $q$, and
$DA_{f,q,a}^{[j]}=A_{f,q,a}^{[j+1]}$ for every integer $j$.
For integration, repeated integration by parts gives, for a polynomial $P$,
\begin{equation}\label{eq:polynomial-ibp}
 \int_0^z P(t)A_{f,q,a}^{[j]}(t)\,dt
 =\sum_{\ell=0}^{\deg P}(-1)^\ell
 \left[P^{(\ell)}(t)A_{f,q,a}^{[j-\ell-1]}(t)\right]_{t=0}^{t=z}.
\end{equation}
All lower-endpoint terms are rational: the value of a nonpositive jet
at zero is zero, whereas for $k\geq1$,
\[
 A_{f,q,a}^{[k]}(0)=f(k)\mathbf1_{k\equiv a\ (q)}\in\Z.
\]
Also $JE_{q,a}=E_{q,a+1}-\mathbf1_{a+1\equiv0\ (q)}$;
the same integration-by-parts argument treats polynomial multiples of
these functions.  Thus the module is closed under $J$ as well.
Expanding a polynomial kernel in its two variables reduces it to these
operations.  This proves existence and gives a finite exact algorithm.

For uniqueness, any nonzero rational-function relation among the displayed
generators could be evaluated at one of the infinitely many points $2/b$,
avoiding every pole and the zeros of a nonzero coefficient.  It would
contradict Corollary~\ref{cor:two-sided-jet}.  In fact the generators are
linearly independent over $\Q(z)$, which is stronger than the required
uniqueness over $\Q[z]$.  At a fixed permitted $x$, their rational linear
independence gives exactly the stated rationality test.
\end{proof}

For example, for $q=1$ and $A=A_f$, an actual nested expression reduces to
\begin{align*}
 &\int_0^x t^2\left(\int_0^t(s+1)A^{(3)}(s)\,ds\right)dt\\
 &\quad=(x^3+x^2)A'(x)-(4x^2+2x)A(x)
       +(8x+2)A^{[-1]}(x)-8A^{[-2]}(x)
       -\frac{x^3}{3}\bigl(f(2)-f(1)\bigr).
\end{align*}
The coefficient of $A^{[-2]}(x)$ is $-8$, so this number is irrational
at every permitted nonzero $x$.  At $x=2$ and $f=\sigma_5$, the formula is
\[
 12A_5'(2)-20A_5(2)+18A_5^{[-1]}(2)-8A_5^{[-2]}(2)-\frac{256}{3}.
\]
More generally, normal-form coefficient vectors decide all rational
linear relations among any finite collection of nested expressions.
Identities such as $DJG=G$ are automatically respected; the theorem does
not incorrectly regard different operator descriptions as independent.

\subsection{All polynomial modes, beta averages, and selected multiple integrals}

For $P\in\Q[t]$ define
\[
 \mathcal L_{f,a,P}(x)=\int_0^1P(t)A_{f,q,a}(xt)\,dt.
\]
If $P(t)=\sum_{k=0}^d p_kt^k$, then
\begin{equation}\label{eq:kernel-rational-weight}
 \mathcal L_{f,a,P}(x)=
 \sum_{\substack{n\geq1\\n\equiv a\ (q)}}
 \frac{f(n)x^n}{n!}\sum_{k=0}^d\frac{p_k}{n+k+1}.
\end{equation}
The map $P\mapsto\sum_kp_k/(X+k+1)$ is injective, as is seen from
the residues at its distinct poles.  Corollary~\ref{cor:simple-poles}
therefore gives the following exact criterion: within each fixed $(f,a)$
block, polynomial kernels produce independent values precisely when the
kernels themselves are linearly independent over $\Q$.  Independent
kernel families in all blocks may be adjoined jointly to $1$, the
exponential residue values, and every polynomial arithmetic moment.

One familiar choice is the complete family of shifted Legendre polynomials
\[
 L_m(t)=\frac1{m!}\frac{d^m}{dt^m}\bigl[t^m(t-1)^m\bigr],
 \qquad m\geq0.
\]
Their degrees are $m$ with nonzero rational leading coefficient, so they
form a basis of $\Q[t]$.  Thus all the Legendre coefficients
\[
 \Lambda_{f,a,m}(x)=\int_0^1L_m(t)A_{f,q,a}(xt)\,dt
 \qquad(f\in\mathcal B,\ 0\leq a<q,\ m\geq0)
\]
are jointly independent together with those moments and exponential
values.  At a positive permitted point they are also strictly positive:
integrating by parts $m$ times gives
\[
 \Lambda_{f,a,m}(x)=\frac{x^m}{m!}
   \int_0^1t^m(1-t)^m A_{f,q,a}^{(m)}(xt)\,dt>0.
\]
There are infinitely many positive coefficients in every residue class,
so the last derivative is positive for $0<t<1$.  The same kernel-rank
criterion applies to any other rational polynomial basis, including
orthogonal bases for polynomial weights, without a separate arithmetic
argument.

For integers $u,v\geq1$, the normalized beta average is
\begin{align*}
 \mathcal B_{u,v}A_{f,q,a}(x)
 &=\frac{(u+v-1)!}{(u-1)!(v-1)!}
     \int_0^1t^{u-1}(1-t)^{v-1}A_{f,q,a}(xt)\,dt\\
 &=\frac{(u+v-1)!}{(u-1)!}
   \sum_{\substack{n\geq1\\n\equiv a\ (q)}}
   \frac{f(n)x^n}{n!(n+u)(n+u+1)\cdots(n+u+v-1)}.
\end{align*}
For each fixed $u$, all $v\geq1$ give an independent family, jointly
with the polynomial moments and exponential coordinates: the kernels
$t^{u-1}(1-t)^{v-1}$ have distinct degrees.  These are also the expectations
of $A_{f,q,a}(xT_0)$ under integer-parameter Dirichlet distributions on
any simplex, with $u$ the parameter of $T_0$ and $v$ the sum of the other
parameters.  This interpretation does not create additional independent
values when the two parameters $u,v$ agree.

There is a useful multiple-integral version with a different geometry.
For distinct positive integers $u_1,\ldots,u_s$, termwise integration gives
\[
 \int_{[0,1]^s}\left(\prod_{j=1}^s t_j^{u_j-1}\right)
 A_{f,q,a}(x t_1\cdots t_s)\,dt_1\cdots dt_s
 =\sum_{\substack{n\geq1\\n\equiv a\ (q)}}
       \frac{f(n)x^n}{n!\prod_{j=1}^s(n+u_j)}.
\]
Its rational weight is nonzero with simple negative-integer poles,
so the value is irrational.  Partial fractions decide independence of
several such integrals; different subsets of poles need not give
independent weights.  Distinctness of the $u_j$ is essential for this
particular deduction from Corollary~\ref{cor:simple-poles}.

\subsection{Every quotient of rational affine forms}

Let $(\beta_i)_{i\in I}$ denote any one of the independent coordinate
families proved above, with $\beta_0=1$.  Coefficient vectors below always
have finite support in this fixed family, and
$L_v=\sum_i v_i\beta_i$ for $v\in\Q^{(I)}$.

\begin{theorem}[Exact projective and rank tests]
\label{thm:projective-rank}
Every nonzero vector $v$ has $L_v\ne0$.  For $v\ne0$,
\[
 \frac{L_u}{L_v}\in\Q
 \quad\Longleftrightarrow\quad u\in\Q v.
\]
More generally, for $v_0\ne0$, the numbers
\[
 1,\quad\frac{L_{v_1}}{L_{v_0}},\ldots,
          \frac{L_{v_m}}{L_{v_0}}
\]
are linearly independent over $\Q$ if and only if
$v_0,v_1,\ldots,v_m$ are linearly independent over $\Q$.
\end{theorem}
\begin{proof}
The map $v\mapsto L_v$ is injective by the underlying independence.
The first equivalence follows by applying this injection to $u-rv$,
where $r\in\Q$.  For the last assertion, multiply a proposed rational
relation by the nonzero common denominator $L_{v_0}$ and use the same
injection.  Both implications follow.
\end{proof}

Consequently, whenever $L_u,L_v$ have independent coefficient vectors,
every rational M\"obius transform
\[
 \frac{aL_u+bL_v}{cL_u+dL_v},
 \qquad a,b,c,d\in\Q,\quad ad-bc\ne0,
\]
is well defined and irrational.  This includes arbitrary rational affine
shifts in the original coordinates, not merely quotients of positive sums.
It does not assert a result for general nonlinear rational functions.

For $q=1$, division of the two-sided hierarchy by $e^x$ proves joint
independence of
\[
 \{1,e^{-x}\}\ \cup\
 \{e^{-x}A_f^{[j]}(x):f\in\mathcal B,\ j\in\Z\}.
\]
At a positive permitted $x$, fix one $f_*$ and divide instead by the
positive number $A_{f_*}(x)$.  Then
\[
 1,\quad\frac1{A_{f_*}(x)},\quad
 \frac{e^x}{A_{f_*}(x)},\quad
 \frac{A_f^{[j]}(x)}{A_{f_*}(x)}
 \quad\bigl((f,j)\ne(f_*,0)\bigr)
\]
are jointly independent.  In particular $A_f'(x)/A_f(x)$, every ratio of
two distinct jets of the same base, and every ratio of two distinct
Legendre modes are irrational.  Higher derivatives of $\log A_f$ involve
products of coordinates and are not included by this linear argument.

\subsection{Truncated and arithmetic-biased Poisson laws}

Let $N$ be Poisson with positive permitted mean $x$, fix $M\geq1$, and
let $S$ be a nonempty set of residues modulo $q$.  Condition on
\[
 C=\{N\geq M,\ N\bmod q\in S\}.
\]
Its unnormalized mass is
\[
 D_C=\sum_{a\in S}E_{q,a}(x)
   -\sum_{\substack{0\leq n<M\\n\bmod q\in S}}\frac{x^n}{n!}>0.
\]
The second sum is rational.  Removing finitely many terms from any
arithmetic coordinate likewise subtracts a rational number.
Theorem~\ref{thm:projective-rank}, applied to this common denominator,
therefore proves joint independence of $1$ and all the numbers
\begin{gather*}
 \mathbb E\!\left[N^d f(N)\mathbf1_{N\equiv a\ (q)}\mid C\right],
       \qquad f\in\mathcal B,\ a\in S,\ d\geq0,\\
 \mathbb E\!\left[\frac{f(N)}{N+u}
          \mathbf1_{N\equiv a\ (q)}\mid C\right],
       \qquad f\in\mathcal B,\ a\in S,\ u\geq1.
\end{gather*}
The moment and simple-pole coordinates are independent by
Corollary~\ref{cor:simple-poles}; each retains its own coefficient after
the finite truncation, whereas $D_C$ has a nonzero exponential coordinate.
This verifies the required rank directly.  Hence arbitrary finite cutoffs
and residue conditioning preserve these joint irrationality statements.

There is a further family using an arithmetic weight in the probability
law itself.  Let $f$ have a nontrivial admissible profile and let $m\geq1$.
On the same infinite support $C$, define
\[
 \Pr_m(N=n)=\frac{f(n)^m x^n/n!}
 {\displaystyle\sum_{\substack{k\geq M\\k\bmod q\in S}}
                       f(k)^m x^k/k!}.
\]
All its moments exist.  Then
\begin{equation}\label{eq:biased-integer-moments}
 \{1\}\ \cup\
 \{\mathbb E_m[f(N)^j]:j\in\Z,\ j\geq-m,\ j\ne0\}
 \quad\hbox{is linearly independent over }\Q.
\end{equation}
Indeed, its numerator at exponent $j$ is the truncated factorial sum
for $f^{m+j}$.  The case $m+j=0$ is an exponential residue sum minus
a rational finite sum.  All positive powers have admissible profiles
$R,R^2,R^3,\ldots$, which are pairwise distinct: a nontrivial positive
rational profile cannot satisfy $R^a=R^b$ for $a\ne b$.
The numerator vectors, together with the denominator vector for $f^m$,
are consequently independent.  The projective rank test proves
\eqref{eq:biased-integer-moments}.

For example, for the law on positive integers proportional to
$\sigma_5(n)^m2^n/n!$, all positive integer moments of $\sigma_5(N)$
and the reciprocal moments of orders $1,\ldots,m$ are irrational and
jointly independent together with $1$.  The same conclusion applies
to the totient, the radical, and every other nontrivial admissible base.
This is a statement about linear relations among expectations.  Variances,
cumulants, entropy, and products of the expectations require additional
arguments and are not inferred here.


\section{A classical companion: all inverse exponential moments}
\label{sec:v5-classical-inverse}

There is a separate, classical transcendence theorem for the pure
exponential channel.  It reaches every nonzero algebraic argument and
arbitrary repeated poles, including a pole at zero.  Its essential input
is Beukers's refinement of the Siegel--Shidlovskii theorem
\cite[Corollary~1.4]{Beukers2006V5}; it is not a new result of the
fiber-cancellation method.  No joint independence between the constants
of this section and the nontrivial arithmetic channels is asserted.

For $s\geq1$ define the entire inverse-moment functions
\[
 \mathcal E_s(z)=\sum_{n\geq1}\frac{z^n}{n!n^s},
 \qquad \mathcal E_0(z)=e^z-1.
\]
They satisfy $\theta\mathcal E_s=\mathcal E_{s-1}$.

\begin{theorem}[Classical inverse-moment independence]
\label{thm:v5-classical-inverse}
At every nonzero algebraic $\alpha$, the infinite family
\[
 1,\quad e^\alpha,\quad
 \mathcal E_1(\alpha),\mathcal E_2(\alpha),\ldots
\]
is linearly independent over $\overline{\Q}$.  In particular every
$\mathcal E_s(\alpha)$ is transcendental.  Every nonconstant
$\overline{\Q}$-linear combination of these entries is transcendental;
here nonconstant means that some coefficient other than that of $1$
is nonzero.
\end{theorem}
\begin{proof}
We check explicitly the hypotheses of the cited theorem for any finite
initial segment.  The functions $1,e^z,\mathcal E_1,\ldots,\mathcal E_S$
are $E$-functions.  Indeed their factorial-series coefficients are rational,
bounded in absolute value, and the denominators through index $N$ divide
$\operatorname{lcm}(1,\ldots,N)^S=\exp(O_S(N))$.  For completeness, the
elementary bound used here follows by letting
$\vartheta(y)=\sum_{p\leq y}\log p$: primes in $(m,2m]$ divide
$\binom{2m}{m}$, so dyadic summation gives $\vartheta(y)=O(y)$.
Consequently
\[
 \log\operatorname{lcm}(1,\ldots,N)
 =\sum_{j\geq1}\vartheta(N^{1/j})
 =O(N)+O(\sqrt N\log N)=O(N).
\]
Finally they satisfy the homogeneous first-order system
\[
 1'=0,\qquad (e^z)'=e^z,\qquad
 \mathcal E_1'=(e^z-1)/z,\qquad
 \mathcal E_s'=\mathcal E_{s-1}/z\quad(s\geq2).
\]
Its matrix has coefficients in $\Q[z,1/z]$, so every nonzero
algebraic argument is an ordinary point of the system.

It remains to prove functional linear independence.  As $t\to+\infty$,
\begin{equation}
 \mathcal E_s(-t)
 =-\frac{(\log t)^s}{s!}+O_s((\log t)^{s-1})
 \qquad(s\geq1).
 \label{eq:v5-inverse-log-asymptotic}
\end{equation}
For $s=1$ this follows from
$\mathcal E_1(-t)=\int_0^t(e^{-u}-1)\,du/u$ by splitting at $u=1$.
The identity $d\mathcal E_s(-t)/dt=\mathcal E_{s-1}(-t)/t$
then proves the formula inductively.
Suppose a linear relation over $\mathbb C(z)$ exists and clear its
denominators, obtaining polynomial coefficients.  Among the coefficients
of $1,\mathcal E_1,\ldots,\mathcal E_S$, take their largest degree
$D$, and among those of degree $D$ take the largest index $s$, with
$1$ assigned index zero.  Along $z=-t$, division by
$t^D(\log t)^s$ leaves a nonzero limit by
\eqref{eq:v5-inverse-log-asymptotic}.  Every smaller polynomial degree,
every smaller logarithmic degree at the same polynomial degree, and
the polynomial multiple of $e^{-t}$ vanish in this normalization.
This is impossible.  Thus all these polynomial coefficients vanish,
and the coefficient of $e^z$ vanishes as well.

Beukers's linear-independence theorem for $E$-function systems now
transfers this functional independence to
$\overline{\Q}$-linear independence at $\alpha\ne0$.
If a nonconstant algebraic linear combination had an algebraic value,
moving that value to the coefficient of $1$ would contradict this
independence.
\end{proof}

The theorem also makes quotients explicit.  If $V,W$ are two distinct
entries other than $1$ in the displayed independent family, then
$V/W$ is transcendental; otherwise $V-(V/W)W=0$ would be an algebraic
linear relation.  The same argument gives transcendence of
$\mathcal E_s(\alpha)/e^\alpha$.  These are statements over the
algebraic numbers, stronger than irrationality, but they are classical
consequences of the cited theorem rather than new priority claims.

\begin{corollary}[All integer-pole rational weights in the pure channel]
\label{cor:v5-pure-rational-classification}
Let $R\in\overline{\Q}(X)$ have poles only at $0,-1,-2,\ldots$, of
arbitrary finite orders, and let $\alpha\ne0$ be algebraic.  The value
\[
 V_R(\alpha)=\sum_{n\geq1}\frac{R(n)\alpha^n}{n!}
\]
has an effective finite expression
\[
 V_R(\alpha)=c_0+c_e e^\alpha+
                \sum_{s=1}^S c_s\mathcal E_s(\alpha),
 \qquad c_0,c_e,c_s\in\overline{\Q}.
\]
It is algebraic if and only if $c_e=c_1=\cdots=c_S=0$; otherwise
it is transcendental.  When algebraic, its value is the explicitly
computed number $c_0$.
\end{corollary}
\begin{proof}
Partial fractions reduce the calculation to polynomial weights and
$(n+a)^{-r}$, where $a\geq0$ and $r\geq1$.  A pole at zero gives
$\mathcal E_r(\alpha)$ directly.  For $a\geq1$, write the falling
factorial polynomial as
\[
 X(X-1)\cdots(X-a+1)=\sum_{j=1}^a c_{a,j}X^j.
\]
For every integer $\ell$, put
$\mathcal F_\ell(z)=\sum_{m\geq1}z^m/(m!m^\ell)$.
Thus $\mathcal F_\ell=\mathcal E_\ell$ for $\ell\geq1$,
$\mathcal F_0=e^z-1$, and
$\mathcal F_{-d}=\theta^d(e^z-1)$ for $d\geq1$.
Reindexing by $m=n+a$ yields the exact identity
\begin{equation}
 \sum_{n\geq1}\frac{\alpha^n}{n!(n+a)^r}
 =\alpha^{-a}\sum_{j=1}^a c_{a,j}
                      \mathcal F_{r-j}(\alpha)-\frac1{a^r}.
 \label{eq:v5-all-poles-reduction}
\end{equation}
In extending the reindexed sum to $m\geq1$, the falling factorial
vanishes at $1,\ldots,a-1$; the remaining boundary term $m=a$
is exactly $1/a^r$.  Every $\mathcal F_{-d}(\alpha)$ is an algebraic
multiple of $e^\alpha$: the relevant polynomials are generated by
$Q_0(z)=1$ and $Q_{d+1}(z)=z(Q_d(z)+Q_d'(z))$.
Polynomial weights are treated by the same Euler operator, with the
constant weight contributing $e^\alpha-1$.
These finite operations give the claimed coefficients.  The final
classification is exactly Theorem~\ref{thm:v5-classical-inverse}.
\end{proof}

In particular, for every $a\geq1$, $r\geq2$ and nonzero algebraic
$\alpha$,
\begin{equation}
 \sum_{n\geq1}\frac{\alpha^n}{n!(n+a)^r}
 \quad\hbox{is transcendental}.
 \label{eq:v5-repeated-poles-transcendental}
\end{equation}
Indeed its coefficient of $\mathcal E_{r-1}(\alpha)$ in
\eqref{eq:v5-all-poles-reduction} is
$\alpha^{-a}c_{a,1}=\alpha^{-a}(-1)^{a-1}(a-1)!\ne0$, and no
larger inverse moment occurs.  Simple-pole values have genuine
exceptions, which the criterion detects; for example
\[
 \sum_{n\geq1}\frac1{n!(n+2)}=\frac12.
\]
This example also explains why the nontrivial arithmetic-profile
hypothesis cannot simply be discarded from the earlier
rational-weight theorem.

\section{Exact finite verification and reproducibility}\label{sec:verification}
The proof for every $K$ is the argument above, not a numerical
experiment. The checks below test indexing, arithmetic normalizations,
nonzero output, and the progression calculation.

\subsection{An entirely explicit one-exponent example}
For $K=1$, the general construction gives
\[
 L=2,\quad b=2,\quad F=D=2,\quad B=5,
 \qquad(p_0,t_0)=(5,0),\quad(p_1,t_1)=(7,2).
\]
The integer weights $(c_0,c_1)=(21,-20)$ satisfy
\[
 21\frac5{\sigma_1(5)}-20\frac7{\sigma_1(7)}
       =21\frac56-20\frac78=0.
\]
The lower shifts both equal $5$, and the terminal shifts are $10$
and $12$. Their surviving coefficients are $175/2$ and $-245/2$.
Thus the zero-limit contradiction is already visible in a two-vertex
example. This example recovers an individual result; it is not the
new joint conclusion.

\subsection{A smaller two-exponent certificate}
For a compact actual computation, we discard zero-coefficient
constraints and also retain the lower-degree residual, using a smaller
grid with the required shift and coprimality properties. At $K=2$, the own-degree expansions have the terms
\[
 T_1(n)=\frac{\sigma_1(n+1)}{n+1}
        +\frac{\sigma_1(n+2)}{(n+2)^2}+O(n^{-3/2}),
\]
and
\[
 T_2(n)=\frac{\sigma_2(n+1)}{n+1}
       +\frac{\sigma_2(n+2)}{(n+2)^2}
       +\frac{\sigma_2(n+2)}{(n+2)^3}
       +\frac{\sigma_2(n+3)}{(n+3)^3}+O(n^{-3/2}).
\]
For the check, impose the three array types $(k,i,j)=(1,0,0)$,
$(2,0,0)$, and $(2,1,1)$ on a $4\times4$ grid. Unlike the
stronger general construction, this smaller certificate leaves the
lower-degree residual as well. It is checked as a mixed survivor;
the isolated terminal shift belongs only to $k=2$. The progression
comparison still works term by term using~\eqref{eq:refined}.
This is an audit of the mechanism, not a claim that the smaller
matrix satisfies every equation of Lemma~\ref{lem:weights}.

Take radix $7$, $I_{\max}=24$,
\[
 D=\prod_{\substack{\ell\leq24\\\ell\ \mathrm{prime}}}\ell=223092870,
 \qquad B=1+2D\cdot24=10708457761.
\]
The pairwise-coprimality argument only requires all primes up to
$\max(I_{\max},K+1)$ to divide $D$, so a primorial is sufficient.
Using exact rational elimination gives the following observed outputs.

\begin{center}
\begin{tabular}{lr}
\toprule
Quantity & Verified value\\
\midrule
Constraint matrix & $12\times16$\\
Exact rank / nullity & $12\,/\,4$\\
Nonzero entries of the selected integer vector & $12$\\
Maximum decimal digits in an integer weight & $281$\\
Isolated terminal vertex & $(0,0)$\\
Isolated terminal shift & $32125373283$\\
Separating prime & $38148880801$\\
Exact arithmetic checks & Passed\\
\bottomrule
\end{tabular}
\end{center}

The program also reconstructs each multiplier from its prime
factorization, computes $\sigma_1$ and $\sigma_2$, checks every
matrix equation, constructs the simultaneous CRT residue, verifies
all multiplicativity coprimalities, and computes a nonzero difference
of refined progression means divided by $\zeta(3)$ using rational
Euler factors. The infinite mean identity itself is proved in
Lemma~\ref{lem:average}.

The complete verifier is embedded below using \texttt{filecontents*}.
Compiling this single source writes \texttt{check\_joint\_grid.py};
the PDF does not require that script to be executed. Reproduce with
\begin{lstlisting}
pdflatex -interaction=nonstopmode Simultaneous_Fiber_Cancellation.tex
pdflatex -interaction=nonstopmode Simultaneous_Fiber_Cancellation.tex
python3 check_joint_grid.py
\end{lstlisting}
The verifier uses only the Python standard library and writes the
full weights, factorizations, CRT data, and rational mean certificate
to \texttt{joint\_grid\_K2.json}. It does not approximate the constants
$\alpha_k$ or infer linear independence from floating-point data.

\subsection{Exact checks for mixed-prime interpolation}
The expanded source also embeds \texttt{check\_prime\_responses.py}.
It uses only integers and Python's standard-library \texttt{Fraction}
class. The following results were obtained from executing it:
\begin{center}
\begin{tabular}{lr}
\toprule
Finite check & Result\\
\midrule
Divisor/Jordan response ranks, $H=1,\ldots,6$ & $2,4,6,8,10,12$\\
Falling-factorial inverse transforms & Degrees $0$--$12$\\
Exact polynomial evaluation checks & $546$\\
Single-prime progression-mean checks & $588$\\
Six-channel single-prime rank & $5$\\
Six-channel mixed-prime isolation rank & $6$\\
Nonempty-subset on/off contrast checks & $1152$\\
\bottomrule
\end{tabular}
\end{center}
The six channels are $\sigma_1,\varphi,\psi,J_3,\varphi^3,\psi^2$,
whose normalized response functions are independent of any polynomial
weights needed to align their growth degrees. The single-prime rank
defect is genuine: symbolic polynomial arithmetic verifies
$\delta_\psi=-\delta_\varphi+2\delta_{J_3}$.
For each target, the checker expands the interpolation expression
into $32$ nonempty-subset terms and verifies that every other channel
is zero while the target is nonzero. It also checks the canonical
identity $(1-X)(1+X)=1-X^2$, so $\varphi\psi$ and $J_2$ are not
mistakenly counted as distinct functions.

These are finite auxiliary checks, not a computational proof of
Theorem~\ref{thm:allmoments}. The general proof is
Lemma~\ref{lem:multimean} and Theorem~\ref{thm:transfer}.
After compiling the standalone source, reproduce the additional
checks with
\begin{lstlisting}
python3 check_prime_responses.py
\end{lstlisting}
The program writes \texttt{prime\_response\_checks.json}, containing
the exact matrices and interpolation certificates. These verification
scripts are embedded in this one source; no external data or package
is required to execute them.

\subsection{Exact denominator checks and the concrete value}
The additional embedded program \texttt{check\_dyadic\_height.py}
checks $129$ factorial valuations, $8385$ factorial-quotient
valuations, $100001$ binary-height bounds, and $60$ extremal binary
cases. It verifies $6888$ scaled partial-sum terms and $2624$ scaled
rational-total terms for both signs and several denominators,
including even denominators. All checks passed.
It also computes the exact rational enclosure of $C_*$ stated in
Section~\ref{sec:dyadic}; the lower and upper bounds have the same
first $50$ decimal places. Execute it after compiling with
\begin{lstlisting}
python3 check_dyadic_height.py
\end{lstlisting}
The output \texttt{dyadic\_height\_checks.json} stores the exact
fractions and check counts. These finite computations supplement,
and do not replace, Lemma~\ref{lem:dyadic-height} and
Theorem~\ref{thm:dyadic}.


\subsection{Exact checks for the integrated extension}
The fourth embedded standard-library program,
\texttt{check\_extensions\_v4.py}, checks the new structures using
rational arithmetic. Its independent formulas compare divisor,
radical, lattice, and lcm coefficients with their prime-power profiles;
test the local-response formulas and mixed-prime isolation; and verify
the residue and offset-grid transformations.
\begin{center}
\begin{tabular}{lr}
\toprule
Finite check & Passed instances\\
\midrule
Arithmetic coefficients versus prime-power profiles & $960$\\
Least-common-multiple identities & $80$\\
Local means and valuation responses & $720$\\
Radical valuation-depth checks & $10$\\
Mixed-prime profile isolation & $144$\\
Modulus-four inverse transformations & $625$\\
Periodic-grid pairs & $268$\\
Negative-offset-grid pairs & $386$\\
Shifted-factorial coefficient expansions & $54$\\
Periodic polynomial rearrangements & $155$\\
\bottomrule
\end{tabular}
\end{center}
In particular, the radical profile has a zero depth-one response and
a nonzero depth-two response, demonstrating why the valuation-vector
extension is needed. The oscillatory transformation has determinant
$-8$. The offset checks include the repaired congruence
$n_e\equiv R+1\pmod{p_e}$, including negative shifts.
Run the additional verifier after compiling the standalone source:
\begin{lstlisting}
python3 check_extensions_v4.py
\end{lstlisting}
It writes \texttt{extensions\_v4\_checks.json}. These finite checks
supplement the uniform proofs; they are not proof-assistant certification.


\subsection{Exact checks for the expanded edition}
Four further standard-library Python verifiers are embedded as
\texttt{filecontents*} blocks. Compiling this single source writes all
eight verifier scripts beside the build. Run the four original commands
above and the following commands:
\begin{lstlisting}[language=bash]
python3 check_profile_universe_v5.py
python3 check_projective_functionals.py
python3 check_rational_kernels_v5.py
python3 check_scope_v5.py
\end{lstlisting}
The new checks compare weighted factorizations, gcd and unit-gcd sums,
and subgroup profiles with exact formulas and direct finite counts;
verify nested-integral and Legendre identities coefficient by coefficient;
check product-kernel, bi-Vandermonde, and inverse-moment reductions; and
check thirty cyclotomic CRT matrices and thirty finite-ring counts.
A $42\times19$ exact sample matrix for the signed arithmetic family has
rank $19$. The general signed-family proof uses eventual relations and
Vandermonde arguments, not this finite rank test. All numerical quantities
in these scripts are integers or exact rational numbers.

The new checks validate identities and finite constructions. The universal
growth, progression-mean, and independence claims are proved in the text;
they are not inferred from finite samples. None of the scripts is a Lean
formalization or a substitute for external mathematical review.

\section{Attribution, novelty, and remaining verification}
\begin{center}
\small
\begin{tabular}{p{0.29\textwidth}p{0.63\textwidth}}
\toprule
Component & Provenance and role\\
\midrule
Original joint target & Deajim--Siksek (2011); corroborated by Pratt's introduction. The constants and all-rank quantifier are retained.\\[4pt]
Individual all-exponent architecture & Tokengrinder (2026): factorial expansion, coprime grids, integer limits, and isolated-shift means. These are attributed and reproved here.\\[4pt]
Added joint step & Rational fiber constraints include every arithmetic normalization simultaneously; a wider radix makes all terminal shifts injective.\\[4pt]
Added separation step & Mixed differences and interpolation separate local responses; exact-valuation refinements extend this to full prime-power profile vectors.\\[4pt]
Denominator extension & Binary digit sums bound the common denominator by $O(N)$; one extra cancellation order covers reduced numerators $\pm2$.\\[4pt]
New application range & Polynomial valuation profiles of every logarithmic order; arithmetic convolutions and rough multiplicity; signed sublinear sequences; all cyclotomic coordinates; integral operators and quotient classification.\\[4pt]
Logical status & Theorems~\ref{thm:main}, \ref{thm:profile-transfer}, \ref{thm:periodic-transfer}, \ref{thm:taumoments}, \ref{thm:two-sided}, \ref{thm:v5-log-profile-main}, \ref{thm:v5-rough-product}, and \ref{thm:v5-sublinear} have complete unconditional arguments in this manuscript. No unproved estimate or conjecture is an input.\\[4pt]
Review status & Separate AI proof reviews of the original and integrated arguments, and eight exact finite arithmetic verifiers; no external human referee or Lean formalization of this new joint proof.\\[4pt]
Historical status & No identical unconditional joint resolution was located in the searches performed. Exhaustive absence of prior work and a historically unprecedented foundational principle are not certified.\\
\bottomrule
\end{tabular}
\end{center}

The bounded literature check on 28 October 2026 used the exact title
and author names of~\cite{DS}, searches for unconditional joint
independence of divisor-sum factorial series, Pratt's primary
introduction, the original and extension directories of~\cite{TG},
and a scholarly citation search. The full 2011 article was not
accessible in the retrieved sources; its mathematical target is
attested by the primary abstract and by~\cite{Pratt}. The extension
directory was read at commit
\nolinkurl{46b9fd26361faf669bcd7e4495e63678bfa6a735}.
The source and limits of the priority check are stated so that a
mathematical proof is not confused with a historical certification.
The additional search examined Jordan and Dedekind factorial series,
their powers, and weighted divisor-sum moments. Pratt's introduction
provides prior motivation for the moment hierarchy; it does not state
a named conjecture identical to our full polynomial-Euler theorem.
The degree-one overlap with Erd\H{o}s--Straus~\cite{ES74} is credited
explicitly. No assertion is made that every specialization was
previously open.

Accordingly, this manuscript does not claim that every requested
success condition has already been independently established. It
provides a complete proof for an unchanged, pre-existing,
author-identified target and a concrete additional method. Confirmation
of worldwide priority and external acceptance remain separate tasks.
No solution of the other proposed Erd\H{o}s problems, and no
transcendence conclusion for an individual $\alpha_k$, is asserted.
Neither Robin's inequality nor irrationality of Euler's constant follows
from the results proved here.

The additional bounded literature check on 29 October 2026 examined
factorial series with products, periodic coefficients, and shifted
denominators, and the primary sources for the classical counting
functions used here. Han\v{c}l--Tijdeman~\cite{HT10} already study
higher differences and affine-product denominators; these general
ideas are not claimed as new. The specific offset-grid argument and
the joint profile-and-residue theorem are the statements supplied here.
The search neither located an identical full theorem nor established
historical open status for every individual application. Subpolynomial
numerators at $x=1$ have elementary classical factorial-tail treatments
and are not presented as novelty flagships.


The expanded edition separates three kinds of results: new arguments
developed in this manuscript, explicit applications of those arguments,
and classical results reproduced with attribution. The sublinear criterion
and finite-state recurrence argument are elementary and are not presented
as historically new discoveries. Classical subgroup counts and Menon-type
identities are credited to their literature; their factorial-value
applications are verified by the stated profile conditions. The entire
inverse-exponential-moment section is explicitly based on Beukers's
classical theorem, including the transcendence statements there.
No claim is made that adding these families resolves the unrelated
irrationality questions for Euler's constant, odd zeta values, or rational
Lambert values. The enlarged arithmetic statements have received separate
internal proof reviews, but remain unrefereed mathematical claims.

\begin{thebibliography}{99}
\bibitem{DS}
A.~Deajim and S.~Siksek,
\emph{On the $\Q$-linear independence of the sums
$\sum_{n=1}^{\infty}\sigma_k(n)/n!$},
Journal of Number Theory \textbf{131} (2011), no.~4, 745--749.
\href{https://doi.org/10.1016/j.jnt.2010.11.009}{doi:10.1016/j.jnt.2010.11.009}.
Primary abstract: \url{https://wrap.warwick.ac.uk/id/eprint/41557/}.

\bibitem{Pratt}
K.~Pratt, \emph{The irrationality of a divisor function series of
Erd\H{o}s and Kac}, arXiv:2209.11124v1, 22 September 2022,
Introduction (the Deajim--Siksek criterion, weighted moment functions,
and Theorem~1). Published in Acta Arithmetica \textbf{211} (2023), 193--228,
\href{https://doi.org/10.4064/aa220927-1-9}{doi:10.4064/aa220927-1-9}.
\url{https://arxiv.org/html/2209.11124v1}.

\bibitem{ES74}
P.~Erd\H{o}s and E.~G.~Straus,
\emph{On the irrationality of certain series},
Pacific Journal of Mathematics \textbf{55} (1974), 85--92,
Theorem~3.7.
\href{https://doi.org/10.2140/pjm.1974.55.85}{doi:10.2140/pjm.1974.55.85}.
\url{https://msp.org/pjm/1974/55-1/pjm-v55-n1-p08-p.pdf}.

\bibitem{TG}
Tokengrinder (GitHub account \texttt{tokengr1nder}),
\emph{Erd\H{o}s 252: irrationality of the factorial divisor-sum series},
public Lean development, September 2026.
\url{https://github.com/tokengr1nder/Erdos252}.
Original reviewed source commit:
\nolinkurl{dc071aafce41bbae41caf4c015499db6dafafd11}.
The present proof reconstructs the mathematics and does not claim
to inherit formal verification from that source.

\bibitem{TGgen}
Tokengrinder, \emph{Genuine extensions of the E252 argument},
\texttt{Generalizations/README.md} in~\cite{TG}, inspected at commit
\nolinkurl{46b9fd26361faf669bcd7e4495e63678bfa6a735}.
The section on polynomial divisor sums explicitly distinguishes
joint independence from the proved polynomial-numerator extensions.

\bibitem{Corpus}
KEIO, \emph{Verifiable Mathematical Reasoning Corpus},
snapshot of 28 October 2026, commit
\nolinkurl{510a32c9a3fd2aa0b5c78d3f6812ded723b9ac8f}.
\url{https://github.com/QEEDDD/Verifiable-Mathematical-Reasoning-Corpus}.
Selected relevant drafts: [013], \emph{Profinite Recognition of Mahler
Equations and D-Finite Witt Escape}; [078], \emph{Mixed Mahler
Congruences}. These are methodological context, not assumed theorems.

\bibitem{HT10}
J.~Han\v{c}l and R.~Tijdeman,
\emph{On the irrationality of factorial series II},
Journal of Number Theory \textbf{130} (2010), 595--607.
\href{https://doi.org/10.1016/j.jnt.2009.10.005}{doi:10.1016/j.jnt.2009.10.005}.
\url{https://pub.math.leidenuniv.nl/~tijdeman/hantijd4.pdf}.

\bibitem{Zou06}
Y.~M.~Zou, \emph{Gaussian binomials and the number of sublattices},
Acta Crystallographica A \textbf{62} (2006), 409--410.
\href{https://doi.org/10.1107/S010876730602455X}{doi:10.1107/S010876730602455X}.
\url{https://arxiv.org/abs/math/0610684}.

\bibitem{Toth12}
L.~T\'oth, \emph{On the number of cyclic subgroups of a finite Abelian group},
Bulletin of the Mathematical Society of the Sciences Math\'ematiques
de Roumanie \textbf{55(103)} (2012), 423--428.
\url{https://arxiv.org/abs/1203.6201}.

\bibitem{IM14}
S.~Ikeda and K.~Matsuoka, \emph{On the Lcm-Sum Function},
Journal of Integer Sequences \textbf{17} (2014), Article 14.1.7.
\url{https://cs.uwaterloo.ca/journals/JIS/VOL17/Ikeda/ikeda4.pdf}.


\bibitem{TothMenon23}
L.~T\'oth, \emph{Proofs, generalizations and analogs of Menon's identity:
a survey}, Acta Universitatis Sapientiae, Mathematica
\textbf{15} (2023), 142--197.
\url{https://doi.org/10.2478/ausm-2023-0009}.
\url{https://arxiv.org/abs/2110.07271}.

\bibitem{Hampejs14}
M.~Hampejs, N.~Holighaus, L.~T\'oth, and C.~Wiesmeyr,
\emph{Representing and counting the subgroups of the group
$\Z_m\times\Z_n$}, Journal of Numbers (2014), Article 491428.
\url{https://doi.org/10.1155/2014/491428}.
\url{https://arxiv.org/abs/1211.1797}.

\bibitem{Beukers2006V5}
F.~Beukers, \emph{A refined version of the Siegel--Shidlovskii theorem},
Annals of Mathematics \textbf{163} (2006), 369--379.
\url{https://doi.org/10.4007/annals.2006.163.369}.
\url{https://annals.math.princeton.edu/wp-content/uploads/annals-v163-n1-p08.pdf}.

\end{thebibliography}

\clearpage
\appendix
\section{Complete exact verifier}
The following code is included in the single source. Its role is the
finite audit described in Section~\ref{sec:verification}, not a substitute for the general
proof. Every primality check uses exact trial division.
\lstinputlisting[language=Python]{check_joint_grid.py}
\end{document}
